\chapter{The base theories and their bridge}
\label{prov:ch:base}
\label{base:chap}

Three formal systems recur in every chapter of the article: the native
scalar core $\Nat$, a calculus with eight rules and no connectives; the
ordinary theory $\Teq$, intuitionistic logic with equality over the same
terms; and $\iPRA$, intuitionistic first-order primitive recursive
arithmetic with quantifier-free induction.  This chapter states the three
systems, develops the derived rules of $\Nat$ and the arithmetic of $\Teq$,
and proves that the three systems are bi-interpretable with whole domains,
actual equality and identity composites
(\cref{base:thm:bridge}).  It also fixes the notion of a clause context and
the protocol by which proof checkers are measured.  The later chapters use
these statements without repeating them.

Notation that is fixed here and used throughout: $1\defeq\Suc 0$,
$2\defeq\Suc\Suc 0$, $3\defeq\Suc\Suc\Suc 0$; a term or formula written
$h(v)$ has a distinguished variable $v$, and $h(s)$ abbreviates
$h\subst{s}{v}$; $A\dashv\vdash B$ means that each of $A$, $B$ is derivable
from the other in the system under discussion.

% ---------------------------------------------------------------------
\section{Terms and substitution}
\label{base:sec:terms}

This section draws on the reference statement of the inherited basis filed
with \pkg{certified\_zero} and on the syntax section of \pkg{raw\_clausal}.

\begin{definition}[Terms]
\label{base:def:terms}
There is one numerical sort.  The terms are
\[
 t ::= x \mid 0 \mid \Suc t \mid \It{n}{a}{u}{h},
\]
where $x$ ranges over a countable set of variables.  In the \emph{iterator}
$\It{n}{a}{u}{h}$ the term $n$ is the \emph{counter}, $a$ is the \emph{seed}
and $h$ is the \emph{body}.  The iterator binds the variable $u$ in $h$ and
nowhere else; $u$ is not bound in $n$ or in $a$.  The body may contain other
free variables, which act as fixed numerical parameters.  Terms that differ
only in the names of bound variables are identified.  $\FV(t)$ is the set of
free variables of $t$.
\end{definition}

A \emph{substitution} $\theta$ is a map from variables to terms that is the
identity outside a finite set.  The term $t[\theta]$ is defined by
\[
 x[\theta]=\theta(x),\qquad 0[\theta]=0,\qquad
 (\Suc t)[\theta]=\Suc(t[\theta]),\qquad
 \It{n}{a}{u}{h}[\theta]=\It{n[\theta]}{a[\theta]}{u}{h[\theta]},
\]
where in the last clause the bound variable $u$ is first renamed so that
$\theta(u)=u$ and $u\notin\FV(\theta(x))$ for every $x\neq u$.  The
substitution that sends $x$ to $s$ and fixes all other variables is written
$\subst{s}{x}$.  Substitution is capture-free by construction, and the
following standard facts hold by induction on $t$.

\begin{lemma}[Substitution]
\label{base:lem:subst}
\begin{enumerate}[label=\textup{(\roman*)}]
\item If $x\notin\FV(t)$ then $t\subst{s}{x}=t$.
\item $t[\theta][\theta']=t[\theta'']$, where $\theta''(x)=\theta(x)[\theta']$.
\item If $x\neq z$ and $x\notin\FV(r)$, then
 $t\subst{s}{x}\subst{r}{z}=t\subst{r}{z}\subst{s\subst{r}{z}}{x}$.
\end{enumerate}
\end{lemma}

\begin{definition}[Selector, zero test, parity]
\label{base:def:sel}
For terms $c,a,b$ put
\[
 \sel(c,a,b)\defeq\It{c}{a}{u}{b}\quad(u\notin\FV(b)),\qquad
 Z(c)\defeq\sel(c,1,0),\qquad
 N(n)\defeq\It{n}{0}{u}{Z(u)} .
\]
\end{definition}

The italic letter $N$ names the parity term; the calculus is written with
the sans-serif letter $\Nat$.

The two \emph{computation equations} of the iterator are
\begin{align}
 \It{0}{a}{u}{h}&=a, \label{base:eq:I0}\\
 \It{\Suc n}{a}{u}{h}&=h\subst{\It{n}{a}{u}{h}}{u}. \label{base:eq:IS}
\end{align}
They are equations between distinct terms, to be derived or assumed in a
calculus; no term is identified with its value.  For the selector they read
$\sel(0,a,b)=a$ and $\sel(\Suc c,a,b)=b$, because the bound variable does not
occur in the body.

\begin{remark}[What the language does not contain]
\label{base:rem:nolang}
Every term is finite syntax.  No numerical variable ranges over bodies or
over descriptions of functions, and there is no term that evaluates
descriptions uniformly.  The body of an iterator receives the previous value
only; it does not receive the counter.  The reference statements present the
same language with function symbols instead of a binder: at stage zero the
symbols are $0$ and $\Suc$, and for every earlier-stage term $h(u,\vec p)$ a
symbol $i_h(n,a,\vec p)$ is added with the equations
$i_h(0,a,\vec p)=a$ and $i_h(\Suc n,a,\vec p)=h(i_h(n,a,\vec p),\vec p)$.
The formal parameters of a symbol are distinct and its actual arguments are
arbitrary.  The two presentations have the same terms up to this renaming,
and the article uses the binder.
\end{remark}

% ---------------------------------------------------------------------
\section{Native formulas and their ordinary reading}
\label{base:sec:reading}

This section draws on the reference statements filed with
\pkg{certified\_zero} (the inherited basis and the inherited ordinary
bridge).

\begin{definition}[Native formulas]
\label{base:def:formulas}
The native formulas are
\[
 A ::= (s=t)\mid \Q{y}A \mid \D{x}{y}A .
\]
$\Q{y}$ binds $y$; $\D{x}{y}$ binds the two distinct variables $x$ and $y$.
Formulas that differ only in the names of bound variables are identified,
and $A[\theta]$ is capture-free substitution.  A \emph{literal equation} is
a formula of the form $s=t$; an equation below a binder is not a literal
equation.  The designated contradiction is the equation $0=\Suc 0$.
\end{definition}

There are no connectives, no falsity constant and no ordinary quantifier in
the native language.

\begin{definition}[Ordinary formulas and the ordinary reading]
\label{base:def:reading}
The \emph{ordinary formulas} are the first-order formulas over the terms of
\cref{base:def:terms}, built from equations by $\bot$, $\wedge$, $\vee$,
$\to$, $\forall$ and $\exists$.  The \emph{ordinary reading} $A^\circ$ of a
native formula $A$ is the ordinary formula defined by
\begin{align*}
 (s=t)^\circ &\defeq (s=t),\\
 (\Q{y}B)^\circ &\defeq B^\circ(0)\to\forall y\,B^\circ(y),\\
 (\D{x}{y}B)^\circ &\defeq
   \exists x\,\bigl(B^\circ(x,0)\to\forall y\,B^\circ(x,y)\bigr).
\end{align*}
For a list $\Gamma$ of native formulas, $\Gamma^\circ$ is the list of the
readings of its members.
\end{definition}

Thus $\Q{y}B$ says that $B$ holds everywhere provided it holds at zero, and
$\D{x}{y}B$ says that some $x$ has this property.  The reading uses the whole
numerical domain and reads the equality sign as equality.  By induction on
$A$,
\begin{equation}
 (A[\theta])^\circ = A^\circ[\theta]. \label{base:eq:reading-subst}
\end{equation}
In the Lean formalization the native formulas are \Lean{NFm} and the reading
is \Lean{NFm.toFm}.

% ---------------------------------------------------------------------
\section{The native scalar core}
\label{base:sec:kernel}

This section draws on the reference statement of the inherited basis filed
with \pkg{certified\_zero}; the two rules \rulename{Tr} and \rulename{QI}
are taken in the forms of \pkg{one\_context} and \pkg{direct\_closure}, as in
the Lean kernel.

A \emph{sequent} is $\Gamma\proves A$, where the \emph{context} $\Gamma$ is a
finite list of native formulas and $A$ is a native formula.

\begin{definition}[The native scalar core $\Nat$]
\label{base:def:kernel}
The derivable sequents of $\Nat$ are generated by the following eight rules.
In \rulename{Tr}, $h(v)$ is a term with a distinguished variable $v$.
\begin{gather*}
 \infrule{A\in\Gamma}{\Gamma\proves A}\ \rulename{Hyp}
 \qquad
 \infrule{}{\Gamma\proves\It{0}{a}{u}{h}=a}\ \rulename{I0}
 \\[1.2ex]
 \infrule{}{\Gamma\proves\It{\Suc n}{a}{u}{h}=h\subst{\It{n}{a}{u}{h}}{u}}\
 \rulename{IS}
 \\[1.2ex]
 \infrule{\Gamma\proves s=t\qquad\Gamma\proves h(s)=r}
         {\Gamma\proves r=h(t)}\ \rulename{Tr}
 \qquad
 \infrule{\Gamma\proves\Q{y}B(y)\qquad\Gamma\proves B(0)}
         {\Gamma\proves B(t)}\ \rulename{QE}
 \\[1.2ex]
 \infrule{\Gamma\proves\D{x}{y}B(x,y)\qquad
          \Gamma,\ \Q{y}B(b,y)\proves A}
         {\Gamma\proves A}\ \rulename{DE}
 \\[1.2ex]
 \infrule{\Gamma,\ B(t,0)\proves B(t,\Suc z)}
         {\Gamma\proves\D{x}{y}B(x,y)}\ \rulename{DI}
 \qquad
 \infrule{\Gamma,\ E(z)\proves E(\Suc z)}
         {\Gamma\proves\Q{y}E(y)}\ \rulename{QI}
\end{gather*}
Side conditions.  In \rulename{DE} the variable $b$ is not free in $\Gamma$,
in $A$ or in $\D{x}{y}B$.  In \rulename{DI} the variable $z$ is not free in
$\Gamma$, in $t$ or in $\D{x}{y}B$; the term $t$ is the \emph{witness}.  In
\rulename{QI} the variable $z$ is not free in $\Gamma$ or in $\Q{y}E$, and
$E$ is a literal equation.  The displayed assumptions $\Q{y}B(b,y)$,
$B(t,0)$ and $E(z)$ are discharged.
\end{definition}

The rules are named as follows: hypothesis; the two computation equations;
transport in one term; elimination of $Q$; elimination of $D$;
introduction of $D$ \emph{by cases}; introduction of $Q$ \emph{by
progress}.  The matrix of \rulename{DI} is an arbitrary native formula, and
its premise has one assumption, the instance at zero.  The matrix of
\rulename{QI} is a literal equation, and its premise has one assumption, the
instance at the eigenvariable.

\begin{remark}[The complete inventory]
\label{base:rem:inventory}
The eight rules are all there is.  There is no reflexivity axiom and no
other equality rule; no axiom states that a number is zero or a successor,
that the successor is injective or that $\Suc x\neq 0$; there is no
induction rule other than \rulename{QI}, whose matrix is one equation; there
is no logical rule for a connective, since there is no connective; and there
are no structural rules.  Weakening, substitution and cut are admissible
(\cref{base:lem:struct}).  Every one of the omitted principles is derived in
\cref{base:sec:structural} or, through the bridge, in \cref{base:sec:arith}.
In Lean the calculus is the inductive predicate \Lean{NPrf}, with de Bruijn
indices: a fresh variable is the new index $0$ over a context whose formulas
are shifted.
\end{remark}

\begin{remark}[The calculus as the reports state it]
\label{base:rem:reports}
The research packages state the calculus with two rules in a stronger form:
transport in an arbitrary native formula template,
\[
 \infrule{\Gamma\proves s=t\qquad\Gamma\proves F(s)}{\Gamma\proves F(t)},
\]
and a witnessed introduction by progress, which infers $\D{x}{y}E(x,y)$ from
$\Gamma,E(t,z)\proves E(t,\Suc z)$ for a literal equation $E$.  Both are
derived rules of $\Nat$ (\cref{base:cor:tr,base:cor:combined}).  Conversely,
in the calculus of the reports reflexivity follows by transporting
$\It{0}{t}{u}{0}=t$ along itself in the template $v=t$, symmetry and
transitivity follow by transport in the templates $v=s$ and $s=v$, the rule
\rulename{Tr} follows from the transport of $r=h(s)$ in the template
$r=h(v)$, and \rulename{QI} follows by introducing $\D{x}{y}E(y)$ with $x$
not free in $E$ and opening it by \rulename{DE}, which yields the hypothesis
$\Q{y}E(y)$ itself.  The two calculi therefore derive the same sequents
(Lean: \Lean{iprf\_iff\_nprf}).  The reports also write the calculus with
clause contexts instead of the binder $Q$; that form is described in
\cref{base:sec:clauses}.  The comparison of rule bases is the subject of
later chapters and is not pursued here.
\end{remark}

% ---------------------------------------------------------------------
\section{Structural properties and derived rules}
\label{base:sec:structural}

This section draws on the reference statement of the inherited basis filed
with \pkg{certified\_zero} for the logical interface, and follows the Lean
modules \Lean{NativeStruct}, \Lean{NativeEq}, \Lean{NativeGate} and
\Lean{NativeLogic} in the organization of the proofs.  Throughout the
section $\proves$ is derivability in $\Nat$.

\subsection{Weakening, substitution and cut}

\begin{lemma}[Weakening and substitution]
\label{base:lem:struct}
If $\Gamma\proves A$, then $\Delta\proves A[\theta]$ for every substitution
$\theta$ and every context $\Delta$ that contains $B[\theta]$ for each
$B\in\Gamma$.  In particular $\Gamma\proves A$ implies
$\Gamma[\theta]\proves A[\theta]$, and $\Gamma\proves A$ with
$\Gamma\subseteq\Delta$ implies $\Delta\proves A$.
\end{lemma}

\begin{proof}
By induction on the derivation of $\Gamma\proves A$, for all $\theta$ and
$\Delta$ at once.

\rulename{Hyp}: $A\in\Gamma$ gives $A[\theta]\in\Delta$.

\rulename{I0}, \rulename{IS}: an instance of a computation equation under
$\theta$ is again an instance, for \eqref{base:eq:IS} by
\cref{base:lem:subst}:
$\bigl(h\subst{\It{n}{a}{u}{h}}{u}\bigr)[\theta]
 = h[\theta]\subst{\It{n[\theta]}{a[\theta]}{u}{h[\theta]}}{u}$.

\rulename{Tr}: choose the template variable $v$ fixed by $\theta$ and not
free in its values.  Then $h(s)[\theta]=h[\theta](s[\theta])$ and
$h(t)[\theta]=h[\theta](t[\theta])$, so the two premises given by the
induction hypothesis are premises of \rulename{Tr} with the template
$h[\theta]$.

\rulename{QE}: $B(t)[\theta]=B[\theta](t[\theta])$ and $B(0)[\theta]=
B[\theta](0)$, and the induction hypothesis gives
$\Delta\proves\Q{y}B[\theta]$ and $\Delta\proves B[\theta](0)$.

\rulename{DE}: choose $b'$ not free in $\Delta$, $A[\theta]$,
$\D{x}{y}B[\theta]$ and the values of $\theta$, and let $\theta'$ be $\theta$
changed to send $b$ to $b'$.  Because $b$ is not free in $\Gamma$, $A$ and
$\D{x}{y}B$, the substitutions $\theta$ and $\theta'$ agree on them, and
$\bigl(\Q{y}B(b,y)\bigr)[\theta']=\Q{y}B[\theta](b',y)$.  The induction
hypothesis for the second premise, with $\theta'$ and the context
$\Delta,\Q{y}B[\theta](b',y)$, gives
$\Delta,\Q{y}B[\theta](b',y)\proves A[\theta]$; for the first premise it gives
$\Delta\proves\D{x}{y}B[\theta]$.  \rulename{DE} applies.

\rulename{DI}: choose $z'$ fresh and let $\theta'$ send $z$ to $z'$ and
agree with $\theta$ elsewhere.  The induction hypothesis with $\theta'$ and
the context $\Delta,B[\theta](t[\theta],0)$ gives
$\Delta,B[\theta](t[\theta],0)\proves B[\theta](t[\theta],\Suc z')$, and
\rulename{DI} with the witness $t[\theta]$ applies.

\rulename{QI}: as for \rulename{DI}, with the premise
$\Delta,E[\theta](z')\proves E[\theta](\Suc z')$; the formula $E[\theta]$ is
a literal equation because substitution preserves the root of a formula.
\end{proof}

\begin{lemma}[Cut]
\label{base:lem:cut}
If $\Gamma\proves A$ and $\Delta\proves B$ for every $B\in\Gamma$, then
$\Delta\proves A$.
\end{lemma}

\begin{proof}
The following more general statement is proved by induction on the
derivation of $\Gamma\proves A$: if $\Delta\proves B[\theta]$ for every
$B\in\Gamma$, then $\Delta\proves A[\theta]$.  \rulename{Hyp} is the
assumption.  \rulename{I0} and \rulename{IS} are instances.  \rulename{Tr}
and \rulename{QE} follow from the induction hypotheses for their premises as
in \cref{base:lem:struct}.  For \rulename{DE} choose $b'$ and $\theta'$ as in
that proof.  Every $C\in\Gamma$ satisfies $C[\theta']=C[\theta]$, hence
$\Delta,\Q{y}B[\theta](b',y)\proves C[\theta']$ by weakening, and the new
hypothesis is derivable from itself.  The induction hypothesis for the second
premise gives $\Delta,\Q{y}B[\theta](b',y)\proves A[\theta]$, and
\rulename{DE} applies.  \rulename{DI} and \rulename{QI} are treated in the
same way.
\end{proof}

In Lean the three properties are \Lean{NPrf.weaken}, \Lean{NPrf.subst} and
\Lean{NPrf.cut}.  They are used below without mention.  For a context
$\Gamma$, formulas $A$ and $B$ are \emph{interderivable over $\Gamma$} if
$\Gamma,A\proves B$ and $\Gamma,B\proves A$; for the empty $\Gamma$ this is
$A\dashv\vdash B$.  By cut, interderivable formulas may replace each other
as hypotheses and as conclusions.

\subsection{Closure by cases and numerical cases}

\begin{lemma}[A witness that is not used]
\label{base:lem:dummy}
If $x\notin\FV(B)$, then $\D{x}{y}B\dashv\vdash\Q{y}B$.
\end{lemma}

\begin{proof}
From $\D{x}{y}B$, \rulename{DE} with a fresh $b$ adds the hypothesis
$\Q{y}B\subst{b}{x}$, which is $\Q{y}B$, and \rulename{Hyp} derives the
conclusion $\Q{y}B$.  From $\Q{y}B$, \rulename{DI} with the witness $0$ asks
for $\Q{y}B,\ B(0)\proves B(\Suc z)$, which is \rulename{QE} at $\Suc z$.
\end{proof}

\begin{lemma}[Introduction of $Q$ by cases]
\label{base:lem:qI}
If $\Gamma,B(0)\proves B(\Suc z)$ with $z$ not free in $\Gamma$ and
$\Q{y}B$, then $\Gamma\proves\Q{y}B(y)$.  Here $B$ is an arbitrary native
formula.
\end{lemma}

\begin{proof}
Choose $x$ not free in $B$.  \rulename{DI} with the witness $0$ and the
matrix $B$ derives $\Gamma\proves\D{x}{y}B$ from the given premise, and
\cref{base:lem:dummy} gives $\Gamma\proves\Q{y}B$.
\end{proof}

\begin{corollary}[Numerical cases]
\label{base:cor:cases}
If $\Gamma\proves B(0)$ and $\Gamma,B(0)\proves B(\Suc z)$ with $z$ fresh,
then $\Gamma\proves B(t)$ for every term $t$.  In particular, if
$\Gamma\proves B(0)$ and $\Gamma\proves B(\Suc z)$ with $z$ fresh, then
$\Gamma\proves B(t)$.
\end{corollary}

\begin{proof}
\Cref{base:lem:qI} and \rulename{QE}.
\end{proof}

\begin{lemma}[Introduction of $D$ from a closure]
\label{base:lem:dI}
If $\Gamma\proves\Q{y}B(t,y)$, then $\Gamma\proves\D{x}{y}B(x,y)$.
\end{lemma}

\begin{proof}
\rulename{DI} with the witness $t$ asks for
$\Gamma,B(t,0)\proves B(t,\Suc z)$, which is \rulename{QE} at $\Suc z$
applied to the given closure and the assumption.
\end{proof}

\begin{corollary}[The witnessed and the combined introduction]
\label{base:cor:combined}
Let $E(x,y)$ be a literal equation and $z$ a fresh variable.
\begin{enumerate}[label=\textup{(\roman*)}]
\item If $\Gamma,E(t,z)\proves E(t,\Suc z)$, then
 $\Gamma\proves\D{x}{y}E(x,y)$.
\item If $\Gamma,E(t,0),E(t,z)\proves E(t,\Suc z)$, then
 $\Gamma\proves\D{x}{y}E(x,y)$.
\end{enumerate}
\end{corollary}

\begin{proof}
(i) \rulename{QI} gives $\Gamma\proves\Q{y}E(t,y)$, and \cref{base:lem:dI}
applies.  (ii) Over the context $\Gamma,E(t,0)$, \rulename{QI} gives
$\Q{y}E(t,y)$, and \rulename{QE} at $\Suc z$ with the hypothesis $E(t,0)$
gives $E(t,\Suc z)$.  This is the premise of \rulename{DI} with the witness
$t$.
\end{proof}

In Lean the rules of this subsection are \Lean{NPrf.qI}, \Lean{NPrf.cases}
and \Lean{NPrf.dIProg}.

\begin{remark}[The older paired introduction]
\label{base:rem:olddi}
Part (ii) is the introduction rule of the oldest reports, which retains the
instance at zero and, for an equation, the instance at the eigenvariable.
The inherited basis replaces it by a choice between the two assumptions, one
per inference, and proves that the consequence relation is unchanged.  In one
direction each choice is the old rule after a weakening.  In the other
direction only an old introduction at an equation needs a construction, and
it is the one in the proof of (ii).  The witness $t$ is kept as a parameter
of the matrix $E(t,y)$, and the closure is introduced without a witness; an
introduction and elimination of $\D{x}{y}E(x,y)$ itself would return a fresh
witness, not $t$.  The reference counts inference nodes for its own
presentation, in which the inner step is a witnessed introduction with an
unused coordinate followed by \rulename{DE}: an old derivation with $n$
nodes, $r$ of them introductions at an equation, becomes a derivation with
exactly $n+5r\le 6n$ nodes, the six nodes being the outer introduction by
cases, the inner introduction by progress, \rulename{DE}, \rulename{QE} and
two uses of \rulename{Hyp}.  In $\Nat$ the inner introduction and its
elimination are the single inference \rulename{QI}, so four nodes replace
one.  The bound counts nodes of the inference tree; it is not a bound on the
size of formulas or of serialized certificates.
\end{remark}

\subsection{Equality}

\begin{lemma}[Equality rules]
\label{base:lem:eq}
The following rules are derivable, each by the stated number of applications
of \rulename{Tr}.
\begin{enumerate}[label=\textup{(\roman*)}]
\item Reflexivity: $\Gamma\proves t=t$ (one application, on two instances of
 \rulename{I0}).
\item Symmetry: $\Gamma\proves s=t$ implies $\Gamma\proves t=s$ (one more).
\item Transitivity: $\Gamma\proves s=t$ and $\Gamma\proves t=r$ imply
 $\Gamma\proves s=r$ (one more).
\item Congruence: $\Gamma\proves s=t$ implies $\Gamma\proves h(s)=h(t)$ for
 every term template $h(v)$ (one more).
\item Transport in an equation template: $\Gamma\proves s=t$ and
 $\Gamma\proves a(s)=b(s)$ imply $\Gamma\proves a(t)=b(t)$ (two).
\end{enumerate}
\end{lemma}

\begin{proof}
Write $\rulename{Tr}_h$ for the rule with the template $h(v)$.
(i) Apply $\rulename{Tr}_v$ to the premises $\It{0}{t}{u}{u}=t$ and
$\It{0}{t}{u}{u}=t$, both instances of \rulename{I0}; the conclusion is
$t=t$.
(ii) Apply $\rulename{Tr}_v$ to $s=s$ and $s=t$; the conclusion is $t=s$.
(iii) Apply $\rulename{Tr}_v$ to $t=r$ and $t=s$, the latter by (ii); the
conclusion is $s=r$.
(iv) Apply $\rulename{Tr}_h$ to $s=t$ and $h(s)=h(s)$; the conclusion is
$h(s)=h(t)$.
(v) $\rulename{Tr}_a$ applied to $s=t$ and $a(s)=b(s)$ gives $b(s)=a(t)$,
and $\rulename{Tr}_b$ applied to $s=t$ and $b(s)=a(t)$ gives $a(t)=b(t)$.
\end{proof}

By (ii)--(iv), provable equality over a fixed context is a congruence for
every term template, and chains of equations are used below in the usual
way.  The selector equations
\begin{equation}
 \Gamma\proves\sel(0,a,b)=a,\qquad\Gamma\proves\sel(\Suc c,a,b)=b
 \label{base:eq:sel}
\end{equation}
are instances of \rulename{I0} and \rulename{IS}.

\begin{lemma}[Simultaneous congruence]
\label{base:lem:cong}
If $\Gamma\proves s_i=t_i$ for $i=1,\dots,k$, then
$\Gamma\proves a\subst{\vec s}{\vec x}=a\subst{\vec t}{\vec x}$ for every
term $a$.
\end{lemma}

\begin{proof}
Let $a_j$ be $a$ with $t_1,\dots,t_j,s_{j+1},\dots,s_k$ substituted for
$\vec x$.  With a fresh variable $v$ and the template
$h_j(v)$ obtained by substituting $t_1,\dots,t_{j-1},v,s_{j+1},\dots,s_k$,
one has $a_{j-1}=h_j(s_j)$ and $a_j=h_j(t_j)$ as terms, so
\cref{base:lem:eq}(iv) gives $\Gamma\proves a_{j-1}=a_j$.  Transitivity
joins $a_0=a\subst{\vec s}{\vec x}$ to $a_k=a\subst{\vec t}{\vec x}$.
\end{proof}

\begin{lemma}[Equation induction]
\label{base:lem:eqind}
Let $E(y)$ be a literal equation and $z$ a variable not free in $\Gamma$ and
$\Q{y}E$.  If $\Gamma\proves E(0)$ and $\Gamma,E(z)\proves E(\Suc z)$, then
$\Gamma\proves E(t)$ for every term $t$.
\end{lemma}

\begin{proof}
\rulename{QI} and \rulename{QE}.
\end{proof}

In Lean this rule is \Lean{NPrf.eqInd}.

\subsection{Congruence of the binders and the Leibniz rule}

\begin{lemma}[Congruence of $Q$ and $D$]
\label{base:lem:bindercong}
Let $B$ and $B'$ be interderivable over $\Gamma$.
\begin{enumerate}[label=\textup{(\roman*)}]
\item If $y$ is not free in $\Gamma$, then $\Q{y}B$ and $\Q{y}B'$ are
 interderivable over $\Gamma$.
\item If $x$ and $y$ are not free in $\Gamma$, then $\D{x}{y}B$ and
 $\D{x}{y}B'$ are interderivable over $\Gamma$.
\end{enumerate}
\end{lemma}

\begin{proof}
Both parts rest on one derivation.  Let $C(y)$ and $C'(y)$ be
interderivable over $\Gamma$, with $y$ not free in $\Gamma$, and let $z$ be
fresh.  Then
\begin{equation}
 \Gamma,\ \Q{y}C(y),\ C'(0)\proves C'(\Suc z). \label{base:eq:qstep}
\end{equation}
Indeed, substituting $0$ for $y$ in $\Gamma,C'\proves C$ gives
$\Gamma,C'(0)\proves C(0)$, because $y$ is not free in $\Gamma$;
\rulename{QE} at $\Suc z$ gives $C(\Suc z)$; and substituting $\Suc z$ for
$y$ in $\Gamma,C\proves C'$ gives $\Gamma,C(\Suc z)\proves C'(\Suc z)$.

(i) \Cref{base:lem:qI} applied to \eqref{base:eq:qstep} for $C=B$ and
$C'=B'$ gives $\Gamma,\Q{y}B\proves\Q{y}B'$.  The converse is symmetric.

(ii) Over $\Gamma,\D{x}{y}B$ apply \rulename{DE} with a fresh $b$; the new
hypothesis is $\Q{y}B(b,y)$.  Substituting $b$ for $x$ shows that $B(b,y)$
and $B'(b,y)$ are interderivable over $\Gamma$.  \rulename{DI} with the
witness $b$ and the matrix $B'$ asks for
$\Gamma,\Q{y}B(b,y),B'(b,0)\proves B'(b,\Suc z)$, which is
\eqref{base:eq:qstep} for $C=B(b,y)$ and $C'=B'(b,y)$.  Hence
$\Gamma,\D{x}{y}B\proves\D{x}{y}B'$, and the converse is symmetric.
\end{proof}

The binder $Q$ is contravariant in its instance at zero, which is why
both directions of the hypothesis are used in each direction of the
conclusion.

\begin{theorem}[Leibniz rule]
\label{base:thm:leibniz}
Let $A$ be a native formula and let $\Gamma\proves s_i=t_i$ for
$i=1,\dots,k$.  Then $A\subst{\vec s}{\vec x}$ and $A\subst{\vec t}{\vec x}$
are interderivable over $\Gamma$.  In particular
$\Gamma\proves A\subst{\vec s}{\vec x}$ implies
$\Gamma\proves A\subst{\vec t}{\vec x}$.
\end{theorem}

\begin{proof}
By induction on $A$, for all $\Gamma$, $\vec s$ and $\vec t$.  If $A$ is
$a=b$, then \cref{base:lem:cong} gives
$\Gamma\proves a\subst{\vec s}{\vec x}=a\subst{\vec t}{\vec x}$ and the same
for $b$, and from the hypothesis
$a\subst{\vec s}{\vec x}=b\subst{\vec s}{\vec x}$ symmetry and transitivity
give $a\subst{\vec t}{\vec x}=b\subst{\vec t}{\vec x}$; the converse is the
same argument with $\vec s$ and $\vec t$ exchanged.  If $A$ is $\Q{y}B$,
rename $y$ so that it is not free in $\Gamma$, $\vec s$, $\vec t$ and
differs from $\vec x$.  The induction hypothesis makes
$B\subst{\vec s}{\vec x}$ and $B\subst{\vec t}{\vec x}$ interderivable over
$\Gamma$, and \cref{base:lem:bindercong}(i) applies.  The case $\D{x}{y}B$
is the same with \cref{base:lem:bindercong}(ii).
\end{proof}

\begin{corollary}[Transport in a formula template]
\label{base:cor:tr}
If $\Gamma\proves s=t$ and $\Gamma\proves F(s)$, then $\Gamma\proves F(t)$,
for every native formula template $F(v)$.
\end{corollary}

In Lean these are \Lean{NPrf.leibniz} and \Lean{NPrf.tr}; the two-step rule
of \cref{base:lem:eq}(v) is \Lean{NPrf.trEq}.  Transport never has to be
primitive below a binder: the template is taken apart by \rulename{QE},
\rulename{DE} and \rulename{DI}.

\subsection{Gates and selectors of formulas}

A case distinction on a term cannot be expressed by a connective.  It is
pushed into a formula, down to its equations.

\begin{definition}[Selector and gate of formulas]
\label{base:def:gate}
For a term $v$ and native formulas $A$, $B$ the native formula
$\sel(v,A,B)$ is defined by recursion, first on $A$ and, for an equation
$A$, on $B$.  Bound variables are chosen fresh for $v$ and for the other
formula.
\begin{align*}
 \sel(v,\ s=t,\ r=w) &\defeq \bigl(\sel(v,s,r)=\sel(v,t,w)\bigr),\\
 \sel(v,\ s=t,\ \Q{y}B) &\defeq
   \Q{y}\,\sel\bigl(v,\ \sel(y,0,s)=\sel(y,0,t),\ B\bigr),\\
 \sel(v,\ s=t,\ \D{x}{y}B) &\defeq
   \D{x}{y}\,\sel\bigl(v,\ \sel(y,0,s)=\sel(y,0,t),\ B\bigr),\\
 \operatorname{gate}(v,B) &\defeq \sel(v,\ 0=0,\ B),\\
 \sel(v,\ \Q{y}A,\ B) &\defeq \Q{y}\,\sel\bigl(v,\ A,\ \operatorname{gate}(y,B)\bigr),\\
 \sel(v,\ \D{x}{y}A,\ B) &\defeq
   \D{x}{y}\,\sel\bigl(v,\ A,\ \operatorname{gate}(y,B)\bigr).
\end{align*}
\end{definition}

The intention is that $\sel(v,A,B)$ is $A$ when $v$ is zero and $B$ when $v$
is a successor, and that $\operatorname{gate}(v,B)$ is trivially derivable
when $v$ is zero and is $B$ when $v$ is a successor.  When a binder of one
formula has no counterpart in the other, the other formula is guarded by the
bound variable, so that the instance of the binder at zero is trivial.  By
induction on the definition, substitution commutes with both operations:
\begin{equation}
 \sel(v,A,B)[\theta]=\sel(v[\theta],A[\theta],B[\theta]),\qquad
 \operatorname{gate}(v,B)[\theta]=\operatorname{gate}(v[\theta],B[\theta]).
 \label{base:eq:sel-subst}
\end{equation}

\begin{lemma}[A guarded equation under a binder]
\label{base:lem:qsel}
If $x$ and $y$ are not free in $s$ and $t$, then
\[
 \Q{y}\bigl(\sel(y,0,s)=\sel(y,0,t)\bigr)\dashv\vdash (s=t)
 \dashv\vdash\D{x}{y}\bigl(\sel(y,0,s)=\sel(y,0,t)\bigr).
\]
\end{lemma}

\begin{proof}
From the closure: its instance at $0$ is $\sel(0,0,s)=\sel(0,0,t)$, which is
derivable because both sides equal $0$ by \eqref{base:eq:sel}.
\rulename{QE} at $1$ gives $\sel(1,0,s)=\sel(1,0,t)$, and
\eqref{base:eq:sel} turns it into $s=t$.  From $s=t$:
\eqref{base:eq:sel} gives
$\sel(\Suc z,0,s)=s=t=\sel(\Suc z,0,t)$, which is the premise of
\cref{base:lem:qI}.  The second equivalence follows from the first by
\cref{base:lem:dummy}.
\end{proof}

\begin{lemma}[Laws of the gate and the selector]
\label{base:lem:gate}
For all terms $v$ and native formulas $A$, $B$:
\begin{enumerate}[label=\textup{(\roman*)}]
\item $\proves\operatorname{gate}(0,B)$;
\item $\operatorname{gate}(\Suc v,B)\dashv\vdash B$;
\item $\Q{y}\operatorname{gate}(y,B)\dashv\vdash B\dashv\vdash
 \D{x}{y}\operatorname{gate}(y,B)$ for $x,y\notin\FV(B)$;
\item $\sel(0,A,B)\dashv\vdash A$;
\item $\sel(\Suc v,A,B)\dashv\vdash B$.
\end{enumerate}
\end{lemma}

\begin{proof}
First let $A$ be an equation $s=t$.  The two statements
\[
 \sel(0,\ s=t,\ B)\dashv\vdash(s=t),\qquad
 \sel(\Suc v,\ s=t,\ B)\dashv\vdash B
\]
are proved by induction on $B$, for all $s$, $t$ and $v$.  If $B$ is $r=w$,
they hold by \eqref{base:eq:sel} and the equality rules.  If $B$ is
$\Q{y}B'$, the induction hypothesis for $B'$ and the equation
$\sel(y,0,s)=\sel(y,0,t)$ gives
\begin{gather*}
 \sel\bigl(0,\ \sel(y,0,s)=\sel(y,0,t),\ B'\bigr)\dashv\vdash
 \bigl(\sel(y,0,s)=\sel(y,0,t)\bigr),\\
 \sel\bigl(\Suc v,\ \sel(y,0,s)=\sel(y,0,t),\ B'\bigr)\dashv\vdash B' .
\end{gather*}
By \cref{base:lem:bindercong}(i) the first formula of the statement is
interderivable with $\Q{y}(\sel(y,0,s)=\sel(y,0,t))$, hence with $s=t$ by
\cref{base:lem:qsel}, and the second with $\Q{y}B'$.  The case
$\D{x}{y}B'$ is the same with \cref{base:lem:bindercong}(ii) and the second
equivalence of \cref{base:lem:qsel}.

(i) and (ii) are the two statements for the equation $0=0$, which is
derivable by reflexivity.

(iii) From $\Q{y}\operatorname{gate}(y,B)$: the instance at $0$ is
derivable by (i), \rulename{QE} at $1$ gives $\operatorname{gate}(1,B)$, and
(ii) gives $B$.  From $B$: (ii) gives $\operatorname{gate}(\Suc z,B)$, the
premise of \cref{base:lem:qI}.  The second equivalence follows by
\cref{base:lem:dummy}.

(iv) and (v) are proved by induction on $A$, for all $v$ and $B$.  For an
equation they are the two statements above.  If $A$ is $\Q{y}A'$, then
$\sel(0,A,B)=\Q{y}\sel(0,A',\operatorname{gate}(y,B))$ is interderivable
with $\Q{y}A'$ by the induction hypothesis and
\cref{base:lem:bindercong}(i); and
$\sel(\Suc v,A,B)=\Q{y}\sel(\Suc v,A',\operatorname{gate}(y,B))$ is
interderivable with $\Q{y}\operatorname{gate}(y,B)$ by the induction
hypothesis and \cref{base:lem:bindercong}(i), hence with $B$ by (iii).  The
case $\D{x}{y}A'$ is the same with \cref{base:lem:bindercong}(ii).
\end{proof}

\begin{lemma}[Selector at an arbitrary index]
\label{base:lem:selIE}
\begin{enumerate}[label=\textup{(\roman*)}]
\item If $\Gamma\proves A$ and $\Gamma\proves B$, then
 $\Gamma\proves\sel(w,A,B)$ for every term $w$.
\item If $\Gamma\proves v=w$ and $\Gamma\proves\sel(v,A,B)$, then
 $\Gamma\proves\sel(w,A,B)$.
\end{enumerate}
\end{lemma}

\begin{proof}
(i) Apply \cref{base:cor:cases} to the formula $\sel(y,A,B)$ with $y$ fresh:
its instance at $0$ follows from $A$ by \cref{base:lem:gate}(iv), and its
instance at $\Suc z$ follows from $B$ by \cref{base:lem:gate}(v).  The
instances are computed by \eqref{base:eq:sel-subst}.
(ii) \Cref{base:cor:tr} with the template $\sel(y,A,B)$.
\end{proof}

\subsection{The connectives as native formulas}

\begin{definition}[Native connectives]
\label{base:def:macros}
With fresh bound variables and the parity term $N$ of \cref{base:def:sel},
\begin{align*}
 \bot_* &\defeq (0=\Suc 0),\\
 A\Rightarrow_* B &\defeq \Q{y}\,\sel(y,A,B),\\
 A\wedge_* B &\defeq
   \Q{y}\,\operatorname{gate}\bigl(y,\ \sel(N(y),A,B)\bigr),\\
 \exists_* x\,A &\defeq \D{x}{y}\,\operatorname{gate}(y,A),\\
 A\vee_* B &\defeq \exists_* x\,\sel(x,A,B),\\
 \forall_* x\,A &\defeq A(0)\wedge_*\Q{x}A .
\end{align*}
\end{definition}

By \eqref{base:eq:sel-subst} each of these operations commutes with
substitution.  The parity term satisfies
\begin{equation}
 \proves N(0)=0,\qquad \proves N(1)=1,\qquad \proves N(2)=0,
 \label{base:eq:par012}
\end{equation}
by \rulename{I0}, by $N(\Suc 0)=Z(N(0))=Z(0)=\sel(0,1,0)=1$, and by
$N(\Suc 1)=Z(N(1))=Z(1)=\sel(\Suc 0,1,0)=0$.  The conjunction reads its two
operands off the instances at $1$ and at $2$.

\begin{theorem}[Natural deduction for the native connectives]
\label{base:thm:logic}
The following rules are derivable in $\Nat$, for all native formulas.
\begin{enumerate}[label=\textup{(\roman*)}]
\item $\bot_*$E: $\Gamma\proves\bot_*$ implies $\Gamma\proves A$.
\item $\Rightarrow_*$I: $\Gamma,A\proves B$ implies
 $\Gamma\proves A\Rightarrow_* B$.\quad
 $\Rightarrow_*$E: $\Gamma\proves A\Rightarrow_* B$ and $\Gamma\proves A$
 imply $\Gamma\proves B$.
\item $\wedge_*$I: $\Gamma\proves A$ and $\Gamma\proves B$ imply
 $\Gamma\proves A\wedge_* B$.\quad
 $\wedge_*$E: $\Gamma\proves A\wedge_* B$ implies $\Gamma\proves A$ and
 $\Gamma\proves B$.
\item $\exists_*$I: $\Gamma\proves A(t)$ implies
 $\Gamma\proves\exists_* x\,A(x)$.\quad
 $\exists_*$E: $\Gamma\proves\exists_* x\,A(x)$ and $\Gamma,A(b)\proves C$,
 with $b$ not free in $\Gamma$, $C$ and $\exists_* x\,A$, imply
 $\Gamma\proves C$.
\item $\vee_*$I: each of $\Gamma\proves A$, $\Gamma\proves B$ implies
 $\Gamma\proves A\vee_* B$.\quad
 $\vee_*$E: $\Gamma\proves A\vee_* B$, $\Gamma,A\proves C$ and
 $\Gamma,B\proves C$ imply $\Gamma\proves C$.
\item $\forall_*$I: $\Gamma\proves A(x)$ with $x$ not free in $\Gamma$
 implies $\Gamma\proves\forall_* x\,A(x)$.\quad
 $\forall_*$E: $\Gamma\proves\forall_* x\,A(x)$ implies
 $\Gamma\proves A(t)$.
\end{enumerate}
\end{theorem}

\begin{proof}
(i) From $0=\Suc 0$ every equation $s=t$ follows:
$s=\sel(0,s,t)=\sel(\Suc 0,s,t)=t$, by \eqref{base:eq:sel} and congruence
in the first argument of the selector.  Every formula follows by induction
on the number of its binders, a measure that substitution does not change.
For $\Q{y}B$ use \cref{base:lem:qI}, whose premise $B(\Suc z)$ holds by the
induction hypothesis; for $\D{x}{y}B$ use \rulename{DI} with the witness
$0$, whose premise $B(0,\Suc z)$ holds by the induction hypothesis.

(ii) Introduction.  By \cref{base:lem:qI} it suffices to derive
$\Gamma,\sel(0,A,B)\proves\sel(\Suc z,A,B)$.  \Cref{base:lem:gate}(iv) gives
$A$, the premise gives $B$, and \cref{base:lem:gate}(v) gives
$\sel(\Suc z,A,B)$.  Elimination.  \Cref{base:lem:gate}(iv) turns $A$ into
$\sel(0,A,B)$, \rulename{QE} at $1$ gives $\sel(1,A,B)$, and
\cref{base:lem:gate}(v) gives $B$.

(iii) Introduction.  By \cref{base:lem:qI} it suffices to derive
$\operatorname{gate}(\Suc z,\sel(N(\Suc z),A,B))$ from $\Gamma$, and by
\cref{base:lem:gate}(ii) it suffices to derive $\sel(N(\Suc z),A,B)$, which
is \cref{base:lem:selIE}(i).  Elimination.  The instance of the matrix at
$0$ is derivable by \cref{base:lem:gate}(i).  \rulename{QE} at $\Suc c$ and
\cref{base:lem:gate}(ii) give $\Gamma\proves\sel(N(\Suc c),A,B)$ for every
term $c$.  For $c=1$, \eqref{base:eq:par012} and \cref{base:lem:selIE}(ii)
give $\sel(0,A,B)$, hence $A$.  For $c=0$ they give $\sel(1,A,B)$, hence
$B$.

(iv) Introduction.  \rulename{DI} with the witness $t$ asks for
$\operatorname{gate}(\Suc z,A(t))$, which follows from $A(t)$ by
\cref{base:lem:gate}(ii).  Elimination.  \rulename{DE} with the fresh
variable $b$ adds the hypothesis $\Q{y}\operatorname{gate}(y,A(b))$, which
gives $A(b)$ by \cref{base:lem:gate}(iii); cut with the second premise gives
$C$.

(v) Introduction is (iv) with the witnesses $0$ and $1$ and
\cref{base:lem:gate}(iv), (v).  Elimination.  By (iv) it suffices to derive
$C$ from $\Gamma,\sel(b,A,B)$ with $b$ fresh.  Apply \cref{base:cor:cases}
to the formula $\sel(y,A,B)\Rightarrow_* C$ with $y$ fresh.  Its instance at
$0$ follows by (ii) from $\Gamma,\sel(0,A,B)\proves C$, which holds by
\cref{base:lem:gate}(iv) and the premise $\Gamma,A\proves C$; its instance at
$\Suc z$ follows in the same way from \cref{base:lem:gate}(v) and
$\Gamma,B\proves C$.  The instance at $b$ and (ii) give $C$.

(vi) Introduction.  Substituting $0$ for $x$ in the premise gives
$\Gamma\proves A(0)$, and substituting $\Suc z$ gives
$\Gamma\proves A(\Suc z)$, hence $\Gamma\proves\Q{x}A$ by
\cref{base:lem:qI}; (iii) combines the two.  Elimination.  By (iii) the
premise gives $A(0)$ and $\Q{x}A$, and \rulename{QE} at $t$ gives $A(t)$.
\end{proof}

\begin{remark}[Which rules the logical interface uses]
\label{base:rem:noprogress}
No derivation in this section applies \rulename{QI}, except the derivations
of \cref{base:lem:eqind} and \cref{base:cor:combined} and, through
\cref{base:thm:pred}, the elimination rule of the single-binder universal
(\cref{base:prop:single-binder}(ii)).  The equality rules,
the Leibniz rule, numerical cases and the whole of \cref{base:thm:logic} are
derived from the seven other rules.  Cases supply the logic and progress
supplies induction for equations; the two do not mix.  If \rulename{QI}
were allowed at a compound matrix, the proof of \cref{base:lem:eqind} would
give induction for every native formula.  The universal quantifier needs no
numerical case distinction: its elimination is \rulename{QE}.  The
reference statement of the inherited basis, which has no binder $Q$,
defines implication and the universal quantifier through $D$ with an
unused coordinate and eliminates the universal quantifier through the
derived coverage $x=0\vee_*\exists_* v\,(x=\Suc v)$.  Its abbreviations
differ in detail from the ones above (the selector aligns the binders of its
two formulas, and the matrix of the universal quantifier is gated), but they
satisfy the same introduction and elimination rules and are therefore
interderivable with them.
In Lean the connectives and their rules are the module \Lean{NativeLogic}.
\end{remark}

\subsection{A universal quantifier with one binder}
\label{base:sec:single-binder}

This subsection draws on \pkg{single\_binder}, which writes the universal
quantifier with one paired binder over one gated copy of its matrix, at the
price of the predecessor term $P$ of \cref{base:thm:pred}.

\begin{definition}[The single-binder universal]
\label{base:def:single-binder}
Let $P$ be the predecessor term of \cref{base:thm:pred}; its only free
variable is its argument.  For a native formula $A(x)$ and distinct
variables $u$, $z$ not free in $A$, put
\[
 \mathcal U_x A(x)\defeq
 \D{u}{z}\operatorname{gate}\bigl(z,\ A(P(z))\bigr),
\]
where $A(P(z))$ is $A\subst{P(z)}{x}$.  The variable $u$ does not occur in
the matrix, and $\FV(\mathcal U_xA)=\FV(A)\setminus\{x\}$.
\end{definition}

The predecessor is derived in the equational fragment $\mathcal E$ of
\cref{base:sec:arith}, whose rules are rules or derived rules of $\Nat$
obtained without the connectives (\cref{base:lem:E}(iv),
\cref{base:lem:eqind}); no result of this subsection is used there, so the
forward reference is not circular.  The consequence used below is
\begin{equation}
 \Gamma\proves P(\Suc t)=t\qquad\text{for every context $\Gamma$ and term $t$,}
 \label{base:eq:pred-any}
\end{equation}
by substitution of $t$ for the variable of \cref{base:thm:pred} and
weakening.  Because $P(z)$ has no free variable other than $z$,
\eqref{base:eq:sel-subst} gives
$(\mathcal U_xA)[\theta]=\mathcal U_x(A[\theta])$ for every substitution
$\theta$ that fixes $x$, $u$ and $z$ and whose values do not contain them.
Since $u$ does not occur in the matrix, \cref{base:lem:dummy} gives
\begin{equation}
 \mathcal U_xA(x)\dashv\vdash\Q{z}\operatorname{gate}\bigl(z,\ A(P(z))\bigr),
 \label{base:eq:single-binder-q}
\end{equation}
a formula with the same single binder.

\begin{proposition}[Rules of the single-binder universal]
\label{base:prop:single-binder}
For every native formula $A(x)$, context $\Gamma$ and term $t$:
\begin{enumerate}[label=\textup{(\roman*)}]
\item $\mathcal U$I: if $\Gamma\proves A(x)$ and $x$ is not free in
 $\Gamma$, then $\Gamma\proves\mathcal U_xA(x)$;
\item $\mathcal U$E: if $\Gamma\proves\mathcal U_xA(x)$, then
 $\Gamma\proves A(t)$;
\item $\mathcal U_xA\dashv\vdash\forall_* x\,A$.  The derivation of
 $\forall_* x\,A\proves\mathcal U_xA$ does not apply \rulename{QI}; the
 derivation of $\mathcal U_xA\proves\forall_* x\,A$ uses
 \eqref{base:eq:pred-any} once.
\end{enumerate}
\end{proposition}

\begin{proof}
Choose $u$ and $z$ as in \cref{base:def:single-binder} and write
$B(z)\defeq\operatorname{gate}(z,A(P(z)))$.  For every term $r$,
\[
 B(r)=\operatorname{gate}\bigl(r,\ A(P(r))\bigr),
\]
by \eqref{base:eq:sel-subst} and \cref{base:lem:subst}, because $z$ is not
free in $A$ and $P(z)\subst{r}{z}=P(r)$.

(i) Let $z'$ be a variable not free in $\Gamma$ and $A$.  Since $x$ is not
free in $\Gamma$, substituting $P(\Suc z')$ for $x$ in the premise
(\cref{base:lem:struct}) gives $\Gamma\proves A(P(\Suc z'))$, and
\cref{base:lem:gate}(ii) at the term $z'$ gives
$\Gamma\proves\operatorname{gate}(\Suc z',A(P(\Suc z')))$, that is,
$\Gamma\proves B(\Suc z')$.  Weakening by $B(0)$ gives
$\Gamma,B(0)\proves B(\Suc z')$.  This is the premise of \rulename{DI}
with the matrix $B$, in which $u$ does not occur, the witness $0$ and the
eigenvariable $z'$; the side conditions hold because $z'$ is not free in
$\Gamma$, in $0$ or in $\mathcal U_xA$.  \rulename{DI} gives
$\Gamma\proves\D{u}{z}B=\mathcal U_xA$.  The assumption $B(0)$ is not used.

(ii) Let $b$ be a variable not free in $\Gamma$, $A$ and $t$.  By
\rulename{DE} it suffices to derive $A(t)$ over the context $\Gamma'$ that
extends $\Gamma$ by the hypothesis $\Q{z}B(z)$, which is the opened
formula $\Q{z}B\subst{b}{u}$ because $u$ does not occur in $B$; the side
conditions of \rulename{DE} hold because $b$ is not free in $\Gamma$, in
$A(t)$ or in $\mathcal U_xA$.  Over $\Gamma'$:
$B(0)=\operatorname{gate}(0,A(P(0)))$ is derivable by
\cref{base:lem:gate}(i); \rulename{QE} at $\Suc t$ gives
$B(\Suc t)=\operatorname{gate}(\Suc t,A(P(\Suc t)))$;
\cref{base:lem:gate}(ii) at the term $t$ gives $A(P(\Suc t))$; and
\eqref{base:eq:pred-any} with \cref{base:cor:tr} for the template $A(x)$
gives $A(t)$.

(iii) By $\forall_*$E of \cref{base:thm:logic}(vi) at the variable $x$,
$\forall_* x\,A\proves A(x)$, and $x$ is not free in $\forall_* x\,A$, so (i)
gives $\forall_* x\,A\proves\mathcal U_xA$.  The rules used are
$\forall_*$E, substitution, \cref{base:lem:gate}(ii) and \rulename{DI};
none of them applies \rulename{QI} (\cref{base:rem:noprogress}).
Conversely, (ii) at the term $x$ gives $\mathcal U_xA\proves A(x)$, and
$x$ is not free in $\mathcal U_xA$, so $\forall_*$I gives
$\mathcal U_xA\proves\forall_* x\,A$.  This direction uses
\eqref{base:eq:pred-any}, and the derivation of \cref{base:thm:pred}
applies equation induction, that is, \rulename{QI}.
\end{proof}

Thus $\mathcal U$ satisfies the same introduction and elimination rules as
$\forall_*$ and may replace it wherever only these rules are used; the
replacement in the bridge is \cref{base:cor:single-binder-bridge}.  The
interderivability holds in $\Nat$ with all eight rules.  Whether it holds
without \rulename{QI} is open (\cref{base:q:single-binder-noprogress}).

\begin{remark}[What the single binder saves]
\label{base:rem:single-binder-size}
Let the \emph{binder rank} $r(A)$ be the number of occurrences of $Q$ and
$D$ in $A$: $r(s=t)=0$ and $r(\Q{y}B)=r(\D{x}{y}B)=1+r(B)$.  Substitution
preserves rank.  By induction on \cref{base:def:gate}, first for an
equation as first argument by induction on $B$ and then by induction on
$A$,
\[
 r\bigl(\operatorname{gate}(v,B)\bigr)=r(B),\qquad
 r\bigl(\sel(v,A,B)\bigr)=r(A)+r(B):
\]
the selector of \cref{base:def:gate} stacks the binders of its two formulas
and does not align them.  Hence, by \cref{base:def:macros},
$r(A\wedge_* B)=1+r(A)+r(B)$ and
\[
 r(\forall_* x\,A)=1+r(A(0))+r(\Q{x}A)=2\,r(A)+2,\qquad
 r(\mathcal U_xA)=r(A)+1 .
\]
The formula $\forall_* x\,A$ contains two copies of $A$, the instance $A(0)$
and the matrix of $\Q{x}A$; the formula $\mathcal U_xA$ contains one,
$A(P(z))$.  For $n$ universal quantifiers over one equation $E$ the
recursion $r_n=2r_{n-1}+2$, $r_0=0$, gives the rank $2^{n+1}-2$ and $2^n$
copies of $E$ for $\forall_*$, against the rank $n$ and one copy of $E$
for $\mathcal U$.  The cost is moved into the terms: every free occurrence
of $x$ in $A$ is replaced by the fixed term $P(z)$, and the gate replaces
each side $s$ of the terminal equation of $A$ by $\sel(z,a,s)$, where $a$
is a nest of $r(A)$ selectors over $0$, one for each binder of $A$.  No
formula is duplicated.  In the measure of \cref{base:def:termsize},
$|P(z)|=75$, so each occurrence of $x$ adds $74$ nodes.  Any term with the
two properties named in \cref{base:q:single-binder-inverse} may take the
place of $P$; the predecessor $P_{45}$ of \cref{base:prop:pred45} has them and
$|P_{45}(z)|=45$, so with it each occurrence adds $44$ nodes
(\cref{base:rem:pred57}).
\end{remark}

\begin{remark}[The rules as \pkg{single\_binder} states them]
\label{base:rem:single-binder-sources}
\pkg{single\_binder} states the definition for the formulas of the
inherited basis, built from equations and $D$ alone, with the gate
$\operatorname{gate}(v,s=t)\defeq(\sel(v,0,s)=\sel(v,0,t))$ and
$\operatorname{gate}(v,\D{a}{b}B)\defeq\D{a}{b}\operatorname{gate}(v,B)$,
and with a selector that aligns the binders of its two formulas; there the
old universal adds two to the rank instead of $r(A)+2$, and the package
proves that its translation of an ordinary formula has binder rank exactly
its logical depth.  Its introduction transports $A(v)$ along
$P(\Suc v)=v$; the substitution of $P(\Suc z')$ for the eigenvariable in
the proof of \cref{base:prop:single-binder}(i) makes that transport
unnecessary, so that only the elimination uses a law of the predecessor,
and only $P(\Suc t)=t$; the equation $P(0)=0$ is never used.  The package
adds that an introduction by progress serves as well, because the
assumption of the premise is not used.  In $\Nat$ this applies only when
the matrix $\operatorname{gate}(z,A(P(z)))$ is a literal equation, that
is, when $A$ is an equation, since \rulename{QI} is restricted to literal
equations; for a compound $A$ the introduction is by cases.  The package
obtains the preservation of the formula round trips by an induction on
ordinary formulas whose congruence steps it does not write out; they are
given in the proof of \cref{base:cor:single-binder-bridge}.
\end{remark}

% ---------------------------------------------------------------------
\section{The ordinary theory \texorpdfstring{$\Teq$}{Teq}}
\label{base:sec:teq}

This section draws on the reference statement of the inherited ordinary
bridge filed with \pkg{certified\_zero} and on the section on the ordinary
realization of \pkg{raw\_clausal}.

\begin{definition}[The logic]
\label{base:def:logic}
\emph{Intuitionistic logic with equality} over a language of terms is the
natural deduction calculus for sequents $\Gamma\proves\varphi$ of ordinary
formulas with the rules: hypothesis; $\bot$E; introduction and elimination
for $\wedge$, $\vee$, $\to$, $\forall$ and $\exists$, with the usual
eigenvariable conditions; reflexivity $\Gamma\proves t=t$; and the Leibniz
rule, from $\Gamma\proves s_i=t_i$ for all $i$ and
$\Gamma\proves\varphi\subst{\vec s}{\vec x}$ infer
$\Gamma\proves\varphi\subst{\vec t}{\vec x}$.  A \emph{theory} is a set of
formulas, its axioms; every substitution instance of an axiom is derivable
in every context.  Derivability is closed under weakening, substitution and
cut.
\end{definition}

\begin{definition}[$\Teq$]
\label{base:def:teq}
$\Teq$ is the theory in intuitionistic logic with equality over the terms of
\cref{base:def:terms} with the following axioms.
\begin{enumerate}[label=\textup{(T\arabic*)}]
\item The computation equations \eqref{base:eq:I0} and \eqref{base:eq:IS},
 for all terms.
\item Coverage: $x=0\ \vee\ \exists y\,(x=\Suc y)$.
\item Separation: $0=\Suc 0\to\bot$.
\item Induction for literal equations: for every equation $E(z)$,
 \[
  E(0)\to\bigl(\forall z\,(E(z)\to E(\Suc z))\to\forall z\,E(z)\bigr).
 \]
\end{enumerate}
\end{definition}

Injectivity of the successor, $\Suc x\neq 0$ and induction for compound
formulas are not axioms.  In Lean the axioms are \Lean{TeqAx}.

\begin{lemma}[Cases and contextual induction in $\Teq$]
\label{base:lem:teqcases}
\begin{enumerate}[label=\textup{(\roman*)}]
\item For every ordinary formula $\varphi(y)$: if
 $\Gamma\proves_{\Teq}\varphi(0)$ and
 $\Gamma\proves_{\Teq}\varphi(\Suc z)$ with $z$ not free in $\Gamma$ and
 $\varphi$, then $\Gamma\proves_{\Teq}\forall y\,\varphi(y)$.
\item For every equation $E(z)$: if $\Gamma\proves_{\Teq}E(0)$ and
 $\Gamma,E(z)\proves_{\Teq}E(\Suc z)$ with $z$ not free in $\Gamma$, then
 $\Gamma\proves_{\Teq}E(t)$ for every term $t$.
\end{enumerate}
\end{lemma}

\begin{proof}
(i) Let $y$ be fresh and eliminate the disjunction (T2) at $x=y$.  In the
branch $y=0$ the Leibniz rule turns $\varphi(0)$ into $\varphi(y)$.  In the
other branch eliminate the existential quantifier at the fresh variable $z$;
from $y=\Suc z$ the Leibniz rule turns $\varphi(\Suc z)$ into $\varphi(y)$.
Generalize over $y$.
(ii) Discharge $E(z)$, generalize over $z$, apply (T4) and instantiate at
$t$.  The context $\Gamma$ is arbitrary, so the rule holds with quantified
side hypotheses.
\end{proof}

\begin{remark}[Separation is an axiom]
\label{base:rem:sep}
The computation equations, coverage and equation induction hold in the
structure with one element, in which every equation is true.  Separation
fails there.  Hence separation does not follow from the other axioms, and
the native contradiction $0=\Suc 0$ has to be linked to $\bot$ explicitly:
the reading of the native calculus proves every formula from $0=\Suc 0$,
and the translation of $\bot$ is that equation, but neither fact removes
(T3) from $\Teq$.
\end{remark}

% ---------------------------------------------------------------------
\section{Intuitionistic primitive recursive arithmetic}
\label{base:sec:ipra}

This section fixes the usual presentation, as in the reference statement on
equality reflection filed with \pkg{certified\_zero} and in the Lean module
\Lean{PR}.

\begin{definition}[$\iPRA$]
\label{base:def:ipra}
The \emph{function symbols} are generated as follows, each with an arity:
$\mathrm{zero}$ of arity $0$; $\mathrm{succ}$ of arity $1$; the projections
$\pi^k_i$ of arity $k$, for $1\le i\le k$; the composition
$\mathrm{comp}(f;g_1,\dots,g_m)$ of arity $k$, for $f$ of arity $m$ and
$g_1,\dots,g_m$ of arity $k$; and the recursion $\mathrm{prec}(f,g)$ of
arity $k+1$, for $f$ of arity $k$ and $g$ of arity $k+2$.  The terms are the
variables and the applications $f(t_1,\dots,t_k)$; $0$ is
$\mathrm{zero}()$ and $\Suc t$ is $\mathrm{succ}(t)$.  $\iPRA$ is the
theory in intuitionistic logic with equality over these terms with the
axioms
\begin{align*}
 \pi^k_i(\vec x)&=x_i, &
 \mathrm{comp}(f;\vec g)(\vec x)&=f(g_1(\vec x),\dots,g_m(\vec x)),\\
 \mathrm{prec}(f,g)(0,\vec x)&=f(\vec x), &
 \mathrm{prec}(f,g)(\Suc y,\vec x)&=g\bigl(y,\mathrm{prec}(f,g)(y,\vec x),\vec x\bigr),\\
 \Suc x=0&\to\bot, & \Suc x=\Suc y&\to x=y,
\end{align*}
and induction for quantifier-free formulas: for every quantifier-free
$\varphi(z)$,
\[
 \varphi(0)\to\bigl(\forall z\,(\varphi(z)\to\varphi(\Suc z))\to
 \forall z\,\varphi(z)\bigr).
\]
\end{definition}

The recursion argument of $\mathrm{prec}(f,g)$ comes first.  Distinct
descriptions are distinct symbols, even when they compute the same function
on the natural numbers.  The quantifiers of the logic are unrestricted; only
the induction formula is quantifier free.  In Lean the axioms are
\Lean{PRAAx}.

% ---------------------------------------------------------------------
\section{Arithmetic inside \texorpdfstring{$\Teq$}{Teq}}
\label{base:sec:arith}

This section draws on the reference statements filed with
\pkg{certified\_zero} (scalar predecessor, equality reflection, primitive
recursion) and, for the predecessor, on the equational derivations of
\pkg{uniform\_raw}, \pkg{newest\_parameter}, \pkg{parameter\_witness} and
\pkg{no\_witness}, with a smaller predecessor from \pkg{479\_predecessor57},
\pkg{absorbing\_orbit}, \pkg{absorbing\_predecessor}, \pkg{pred50} and
\pkg{fused\_iteration}, the replacement of coverage from \pkg{zero\_test\_split},
and the diagonal counter and the orbit equality test from
\pkg{orbit\_equality}.  The language
has no predecessor and no pairing, and the
body of an iterator does not see the counter.  The arithmetic is developed
in this order: a
predecessor term; truncated subtraction and a term that reflects equality;
induction for quantifier-free formulas; pairing; primitive recursion.  No
later item is used for an earlier one.

\subsection{The equational fragment}

\begin{definition}[The fragment $\mathcal E$]
\label{base:def:E}
A sequent of $\mathcal E$ is $\Gamma\proves s=t$ with $\Gamma$ a finite list
of equations.  The rules of $\mathcal E$ are \rulename{Hyp}, \rulename{I0},
\rulename{IS}, \rulename{Tr} and \emph{equation induction}:
\[
 \infrule{\Gamma\proves E(0)\qquad\Gamma,E(z)\proves E(\Suc z)}
         {\Gamma\proves E(t)}\ \rulename{EI}
 \qquad(\text{$E$ an equation, $z$ not free in $\Gamma$ and in $E(0)$}).
\]
\end{definition}

\begin{lemma}[Properties of $\mathcal E$]
\label{base:lem:E}
\begin{enumerate}[label=\textup{(\roman*)}]
\item $\mathcal E$ is closed under weakening, substitution and cut.
\item Reflexivity, symmetry, transitivity and congruence
 (\cref{base:lem:eq,base:lem:cong}) are derived rules of $\mathcal E$.
\item \emph{Equational cases}: if $\Gamma\proves E(0)$ and
 $\Gamma\proves E(\Suc z)$ with $z$ fresh, then $\Gamma\proves E(t)$.
\item If $\Gamma\proves_{\mathcal E}s=t$, then $\Gamma\proves_{\Nat}s=t$ and
 $\Gamma\proves_{\Teq}s=t$.
\end{enumerate}
\end{lemma}

\begin{proof}
(i) The proofs of \cref{base:lem:struct,base:lem:cut} apply; \rulename{EI}
is treated like \rulename{QI}, with the eigenvariable renamed.  (ii) The
derivations use \rulename{I0} and \rulename{Tr} only.  (iii) Weaken the
second premise by $E(z)$ and apply \rulename{EI}.  (iv) \rulename{EI} is the
derived rule of \cref{base:lem:eqind} in $\Nat$ and of
\cref{base:lem:teqcases}(ii) in $\Teq$; \rulename{Tr} follows in $\Teq$ from
the Leibniz rule; the other rules are shared.
\end{proof}

By (i), an equation derived in $\mathcal E$ with variables that are not free
in the context may be instantiated at arbitrary terms.  Every case
distinction below proves an equation with an independent case variable and
is an instance of (iii); no hypothesis is ever specialized to a case.  The
fragment has no coverage axiom, no disjunction, no separation and no rule
for $D$.  Until \cref{base:thm:pred} all derivations are in
$\mathcal E$, in the empty context unless a context is shown.

\begin{lemma}[Uniqueness of iteration]
\label{base:lem:uniq}
Let $f(z)$ be a term template and $a$, $h(u)$ terms in which $z$ is not
free, and let $z$ be not free in $\Gamma$.  If
$\Gamma\proves f(0)=a$ and $\Gamma\proves f(\Suc z)=h(f(z))$, then
$\Gamma\proves f(c)=\It{c}{a}{u}{h}$ for every term $c$.
\end{lemma}

\begin{proof}
\rulename{EI} on the equation $f(z)=\It{z}{a}{u}{h}$.  The base is
$f(0)=a=\It{0}{a}{u}{h}$.  In the step,
$f(\Suc z)=h(f(z))=h(\It{z}{a}{u}{h})=\It{\Suc z}{a}{u}{h}$, by the
hypothesis, congruence on the induction equation, and \rulename{IS}.
\end{proof}

In Lean the rule is \Lean{NPrf.uniq}.

\begin{corollary}[Rotation and congruence of the body]
\label{base:cor:rot}
\begin{enumerate}[label=\textup{(\roman*)}]
\item $\proves\It{\Suc y}{a}{u}{h}=\It{y}{h(a)}{u}{h}$.
\item If $\Gamma\proves h(u)=h'(u)$ with $u$ not free in $\Gamma$, then
 $\Gamma\proves\It{c}{a}{u}{h}=\It{c}{a}{u}{h'}$.
\end{enumerate}
\end{corollary}

\begin{proof}
(i) Take $f(z)\defeq\It{\Suc z}{a}{u}{h}$.  Then
$f(0)=h(\It{0}{a}{u}{h})=h(a)$ and $f(\Suc z)=h(f(z))$ by \rulename{IS},
\rulename{I0} and congruence, and \cref{base:lem:uniq} gives the claim.
(ii) Take $f(z)\defeq\It{z}{a}{u}{h}$.  Then $f(0)=a$ and
$f(\Suc z)=h(f(z))=h'(f(z))$, the last step by instantiating the hypothesis
at $u=f(z)$.  \Cref{base:lem:uniq} with the body $h'$ gives the claim.
\end{proof}

Thus no extensionality rule for bodies is assumed; it is an equation
induction.

\subsection{Addition}

Put $a+b\defeq\It{b}{a}{u}{\Suc u}$, $2a\defeq a+a$ and
$3a\defeq(a+a)+a$; the same terms are also written $2\cdot a$ and
$3\cdot a$.

\begin{lemma}[Addition]
\label{base:lem:add}
$\mathcal E$ derives
\begin{gather*}
 a+0=a,\qquad a+\Suc b=\Suc(a+b),\qquad 0+b=b,\qquad \Suc a+b=\Suc(a+b),\\
 a+b=b+a,\qquad (a+b)+c=a+(b+c).
\end{gather*}
Consequently $\Suc a=a+1$, $2(\Suc a)=\Suc\Suc(2a)$,
$3(\Suc a)=\Suc\Suc\Suc(3a)$, $3(a+b)=3a+3b$, $2\cdot 0=0$ and $3\cdot 0=0$.
\end{lemma}

\begin{proof}
The first two equations are \rulename{I0} and \rulename{IS}.  The others
are equation inductions on the last variable shown.
$0+b=b$: the base is $0+0=0$, and $0+\Suc b=\Suc(0+b)=\Suc b$.
$\Suc a+b=\Suc(a+b)$: the base is $\Suc a+0=\Suc a=\Suc(a+0)$, and
$\Suc a+\Suc b=\Suc(\Suc a+b)=\Suc\Suc(a+b)=\Suc(a+\Suc b)$.
Commutativity: the base is $a+0=a=0+a$, and
$a+\Suc b=\Suc(a+b)=\Suc(b+a)=\Suc b+a$.
Associativity: the base is $(a+b)+0=a+b=a+(b+0)$, and
$(a+b)+\Suc c=\Suc((a+b)+c)=\Suc(a+(b+c))=a+\Suc(b+c)=a+(b+\Suc c)$.
The consequences follow by these equations.
\end{proof}

Sums are rearranged below by \cref{base:lem:add} without comment.  No
cancellation law is used, and there is no multiplication.

\subsection{Zero test, interpolation and parity}

\begin{lemma}[Zero test and interpolation]
\label{base:lem:interp}
$\mathcal E$ derives the following, for all terms and every term template
$F(v)$.
\begin{enumerate}[label=\textup{(\roman*)}]
\item $Z(0)=1$ and $Z(\Suc c)=0$.
\item $\sel(c,a,a)=a$ and $\sel(c,\Suc a,\Suc b)=\Suc\sel(c,a,b)$.
\item $Z(Z(Z(c)))=Z(c)$.
\item \emph{Equational interpolation}:
 $F(Z(c))=\sel\bigl(c,F(1),F(0)\bigr)$.
\item \emph{Two-point rule}: if $\Gamma\proves F(0)=F'(0)$,
 $\Gamma\proves F(1)=F'(1)$ and $\Gamma\proves Z(Z(b))=b$, then
 $\Gamma\proves F(b)=F'(b)$.
\end{enumerate}
\end{lemma}

\begin{proof}
(i) is \eqref{base:eq:sel}.  (ii)--(iv) are equational cases on a fresh
variable in the place of $c$.  For (ii): at $0$ both sides of each equation
are $a$, respectively $\Suc a$; at $\Suc z$ they are $a$, respectively
$\Suc b$.  For (iii): $Z(Z(Z(0)))=Z(Z(1))=Z(0)$, and
$Z(Z(Z(\Suc z)))=Z(Z(0))=Z(1)=0=Z(\Suc z)$.  For (iv): at $0$,
$F(Z(0))=F(1)=\sel(0,F(1),F(0))$; at $\Suc z$,
$F(Z(\Suc z))=F(0)=\sel(\Suc z,F(1),F(0))$.
(v) Instantiate (iv) for $F$ and for $F'$ at $c=Z(b)$:
\[
 F(Z(Z(b)))=\sel\bigl(Z(b),F(1),F(0)\bigr)
           =\sel\bigl(Z(b),F'(1),F'(0)\bigr)=F'(Z(Z(b))),
\]
the middle step by congruence from the two hypotheses.  Congruence along
$Z(Z(b))=b$ on both ends gives $F(b)=F'(b)$.
\end{proof}

The two-point rule is the replacement for a case distinction on a Boolean
value.  It is an equality rule, not an elimination of a disjunction: the
term $b$ and the context $\Gamma$ are arbitrary, and $\Gamma$ may contain
equations in which the variables of $b$ occur.

\begin{lemma}[Parity]
\label{base:lem:parity}
$\mathcal E$ derives
\begin{gather*}
 N(0)=0,\qquad N(\Suc n)=Z(N(n)),\qquad N(1)=1,\qquad Z(Z(N(n)))=N(n),\\
 N(n)+Z(N(n))=1,\qquad N(\Suc\Suc n)=N(n),\\
 N(2n)=0,\qquad N(\Suc(2n))=1,\qquad N(3n)=N(n).
\end{gather*}
\end{lemma}

\begin{proof}
The first two equations are computations and give $N(1)=Z(0)=1$.
$Z(Z(N(n)))=N(n)$ is an equation induction on $n$: the base is
$Z(Z(0))=Z(1)=0=N(0)$, and
$Z(Z(N(\Suc n)))=Z(Z(Z(N(n))))=Z(N(n))=N(\Suc n)$ by
\cref{base:lem:interp}(iii); the induction equation is not used.  The
equation $N(n)+Z(N(n))=1$ is the two-point rule for $F(v)\defeq v+Z(v)$,
$F'(v)\defeq 1$ and $b\defeq N(n)$: $F(0)=0+1=1$ and $F(1)=1+0=1$.  Two
computations give $N(\Suc\Suc n)=Z(Z(N(n)))=N(n)$.  $N(2n)=0$ is an
induction: $N(2\cdot 0)=N(0)=0$, and
$N(2(\Suc n))=N(\Suc\Suc(2n))=N(2n)=0$.  Then
$N(\Suc(2n))=Z(N(2n))=Z(0)=1$.  $N(3n)=N(n)$ is an induction:
$N(3\cdot 0)=N(0)$, and
$N(3(\Suc n))=N(\Suc\Suc\Suc(3n))=Z(N(\Suc\Suc(3n)))=Z(N(3n))=Z(N(n))
=N(\Suc n)$.
\end{proof}

No statement that a number is even or odd is involved.

\subsection{A three-cycle and a ramp}

Put
\[
 j(u)\defeq Z(u)+2N(u),\qquad M(n)\defeq\It{n}{0}{u}{j(u)},\qquad
 G(n)\defeq\It{n}{0}{u}{\Suc(u+M(u))} .
\]

\begin{lemma}[Three-cycle]
\label{base:lem:cycle}
$\mathcal E$ derives $j(0)=1$, $j(1)=2$, $j(2)=0$, $M(0)=0$ and
$M(\Suc n)=j(M(n))$.
\end{lemma}

\begin{proof}
$j(0)=Z(0)+2N(0)=1+(0+0)=1$; $j(1)=Z(1)+2N(1)=0+(1+1)=2$;
$j(2)=Z(\Suc 1)+2N(\Suc\Suc 0)=0+2N(0)=0$.  The last two equations are
\rulename{I0} and \rulename{IS}.
\end{proof}

Only these three values of $j$ are used; no theorem about remainders modulo
three is needed.

\begin{lemma}[Ramp]
\label{base:lem:ramp}
$\mathcal E$ derives
\begin{enumerate}[label=\textup{(\roman*)}]
\item $M(G(n))=N(n)$;
\item $G(0)=0$, $G(\Suc n)=\Suc(G(n)+N(n))$ and
 $G(\Suc\Suc n)=\Suc\Suc\Suc G(n)$;
\item $G(2n)=3n$;
\item $G(3n)=3G(n)+N(n)$.
\end{enumerate}
\end{lemma}

\begin{proof}
(i) Two identities are derived first, for variables $g$ and $b$, in the
context that consists of the one equation $Z(Z(b))=b$:
\begin{align}
 M(\Suc(g+b))&=\sel\bigl(b,\ j(M(g)),\ j(j(M(g)))\bigr),
   \label{base:eq:ramp-a}\\
 \sel\bigl(b,\ j(b),\ j(j(b))\bigr)&=Z(b). \label{base:eq:ramp-b}
\end{align}
Both are instances of the two-point rule.  For \eqref{base:eq:ramp-a} the
templates are $F(v)\defeq M(\Suc(g+v))$ and
$F'(v)\defeq\sel(v,j(M(g)),j(j(M(g))))$:
$F(0)=M(\Suc g)=j(M(g))=F'(0)$ and
$F(1)=M(\Suc\Suc g)=j(j(M(g)))=F'(1)$.  For \eqref{base:eq:ramp-b} the
templates are $\sel(v,j(v),j(j(v)))$ and $Z(v)$: at $0$ both sides are
$j(0)=1=Z(0)$, and at $1$ they are $j(j(1))=j(2)=0=Z(1)$, by
\cref{base:lem:cycle}.  By \cref{base:lem:parity} the hypothesis is
derivable for $b=N(n)$, so by substitution and cut both identities hold at
$g=G(n)$ and $b=N(n)$ in every context.

Now \rulename{EI} on $n$.  The base is $M(G(0))=M(0)=0=N(0)$.  In the step
the induction equation is $M(G(n))=N(n)$, and
\begin{align*}
 M(G(\Suc n))
 &=M\bigl(\Suc(G(n)+M(G(n)))\bigr)
  =M\bigl(\Suc(G(n)+N(n))\bigr)\\
 &=\sel\bigl(N(n),\ j(M(G(n))),\ j(j(M(G(n))))\bigr)\\
 &=\sel\bigl(N(n),\ j(N(n)),\ j(j(N(n)))\bigr)
  =Z(N(n))=N(\Suc n).
\end{align*}
The second and fourth steps use the induction equation by congruence, the
third is \eqref{base:eq:ramp-a} and the fifth is \eqref{base:eq:ramp-b}.
The value $N(n)$ is not split under the induction hypothesis.

(ii) The first equation is \rulename{I0}; the second is \rulename{IS} and
(i).  Applying it twice,
$G(\Suc\Suc n)=\Suc(\Suc(G(n)+N(n))+Z(N(n)))
 =\Suc\Suc\bigl(G(n)+(N(n)+Z(N(n)))\bigr)=\Suc\Suc(G(n)+1)
 =\Suc\Suc\Suc G(n)$, by \cref{base:lem:parity,base:lem:add}.

(iii) \rulename{EI}: $G(2\cdot 0)=G(0)=0=3\cdot 0$, and
$G(2(\Suc n))=G(\Suc\Suc(2n))=\Suc\Suc\Suc G(2n)=\Suc\Suc\Suc(3n)
 =3(\Suc n)$.

(iv) \rulename{EI}.  The base is $G(3\cdot 0)=0=3G(0)+N(0)$.  In the step,
by (ii), the induction equation and $N(3n)=N(n)$,
\[
 G(3(\Suc n))=G(\Suc\Suc\Suc(3n))=\Suc\Suc\Suc\,\Suc\bigl(G(3n)+N(3n)\bigr)
 =\Suc\Suc\Suc\Suc\bigl(3G(n)+N(n)+N(n)\bigr),
\]
and, by (ii) and $N(n)+Z(N(n))=1$,
\begin{align*}
 3G(\Suc n)+N(\Suc n)
 &=3\,\Suc\bigl(G(n)+N(n)\bigr)+Z(N(n))
  =\Suc\Suc\Suc\bigl(3G(n)+3N(n)+Z(N(n))\bigr)\\
 &=\Suc\Suc\Suc\bigl(3G(n)+N(n)+N(n)+1\bigr).
\end{align*}
The two right-hand sides are equal by \cref{base:lem:add}.
\end{proof}

\subsection{The factor map and the comparator}

Put
\[
 H(u)\defeq\sel\bigl(N(u),G(u),u\bigr),\qquad
 2^{x}\defeq\It{x}{1}{u}{u+u},\qquad
 L(x,y)\defeq N\bigl(\It{y}{2^{x}}{u}{H(u)}\bigr).
\]
The term $H$ replaces a factor two by a factor three and fixes the numbers
of parity one; $2^x$ is a term, not an operation of the language.

\begin{lemma}[Factor map]
\label{base:lem:factor}
$\mathcal E$ derives $H(2u)=3u$, $H(1)=1$ and $H(3u)=3H(u)$.
\end{lemma}

\begin{proof}
$H(2u)=\sel(N(2u),G(2u),2u)=\sel(0,3u,2u)=3u$ by
\cref{base:lem:parity,base:lem:ramp}.  $H(1)=\sel(N(1),G(1),1)=\sel(1,G(1),1)
=1$.  For the third equation the two-point rule gives, for variables $a$,
$c$ and $b$ in the context $Z(Z(b))=b$,
\[
 \sel(b,\ 3a+b,\ 3c)=3\,\sel(b,a,c):
\]
at $b=0$ both sides are $3a$, and at $b=1$ both sides are $3c$.
Instantiating at $b=N(u)$, $a=G(u)$ and $c=u$,
\[
 H(3u)=\sel\bigl(N(3u),G(3u),3u\bigr)
      =\sel\bigl(N(u),\ 3G(u)+N(u),\ 3u\bigr)=3\,\sel\bigl(N(u),G(u),u\bigr)
      =3H(u).
\]
The identity is proved with an independent variable $b$ and then
instantiated; the dependence of $N(u)$ on $u$ causes no difficulty.
\end{proof}

\begin{lemma}[Orbits of the factor map]
\label{base:lem:orbit}
$\mathcal E$ derives $\It{y}{3a}{u}{H(u)}=3\,\It{y}{a}{u}{H(u)}$ and
$\It{y}{1}{u}{H(u)}=1$.
\end{lemma}

\begin{proof}
Both are equation inductions on $y$ with $a$ fixed.  The bases are
\rulename{I0}.  The steps are
$\It{\Suc y}{3a}{u}{H(u)}=H(\It{y}{3a}{u}{H(u)})=H(3\,\It{y}{a}{u}{H(u)})
 =3H(\It{y}{a}{u}{H(u)})=3\,\It{\Suc y}{a}{u}{H(u)}$ and
$\It{\Suc y}{1}{u}{H(u)}=H(\It{y}{1}{u}{H(u)})=H(1)=1$, by
\cref{base:lem:factor}.
\end{proof}

\begin{lemma}[Comparator]
\label{base:lem:cmp}
$\mathcal E$ derives
\[
 L(0,y)=1,\qquad L(\Suc x,0)=0,\qquad L(\Suc x,\Suc y)=L(x,y),\qquad
 L(\Suc n,n)=0 .
\]
\end{lemma}

\begin{proof}
$2^{0}=1$ and $2^{\Suc x}=2\cdot 2^{x}$ are computations.  Hence
$L(0,y)=N(\It{y}{1}{u}{H(u)})=N(1)=1$ by \cref{base:lem:orbit}, and
$L(\Suc x,0)=N(2\cdot 2^{x})=0$ by \rulename{I0} and
\cref{base:lem:parity}.  For the third equation,
\begin{align*}
 L(\Suc x,\Suc y)
 &=N\bigl(\It{\Suc y}{2\cdot 2^{x}}{u}{H(u)}\bigr)
  =N\bigl(\It{y}{H(2\cdot 2^{x})}{u}{H(u)}\bigr)\\
 &=N\bigl(\It{y}{3\cdot 2^{x}}{u}{H(u)}\bigr)
  =N\bigl(3\,\It{y}{2^{x}}{u}{H(u)}\bigr)
  =N\bigl(\It{y}{2^{x}}{u}{H(u)}\bigr)=L(x,y),
\end{align*}
by rotation (\cref{base:cor:rot}), $H(2u)=3u$, \cref{base:lem:orbit} and
$N(3n)=N(n)$.  The fourth is an equation induction: $L(\Suc 0,0)=0$ by the
second equation, and $L(\Suc\Suc n,\Suc n)=L(\Suc n,n)$ by the third.
\end{proof}

On the natural numbers $L(x,y)$ is the indicator of $x\le y$.  That reading
is not used: only the four equations are.

\subsection{The saturating counter and the predecessor}

Put
\[
 \beta(k,u)\defeq\sel\bigl(L(\Suc\Suc u,k),\,u,\,\Suc u\bigr),\qquad
 W(t,k)\defeq\It{t}{0}{u}{\beta(k,u)},\qquad
 P(m)\defeq W(m,m).
\]
The counter $W(t,k)$ counts up to $t$ and stops below $k$.  Written out,
\[
 P(m)=\It{m}{0}{u}{\sel\bigl(L(\Suc\Suc u,m),\,u,\,\Suc u\bigr)} ,
\]
a term in the one variable $m$, which occurs both as the counter and in the
body.

\begin{theorem}[Predecessor]
\label{base:thm:pred}
The equational fragment derives, in the empty context,
\[
 \proves_{\mathcal E}P(0)=0,\qquad \proves_{\mathcal E}P(\Suc x)=x .
\]
Consequently $\Suc x=\Suc y\proves_{\mathcal E}x=y$, and the same sequents
are derivable in $\Nat$ and in $\Teq$.
\end{theorem}

\begin{proof}
By \rulename{I0} and \rulename{IS},
\begin{equation}
 W(0,k)=0,\qquad W(\Suc t,k)=\beta(k,W(t,k)). \label{base:eq:W}
\end{equation}
By \cref{base:lem:cmp,base:lem:interp}(ii),
\begin{align}
 \beta(\Suc k,\Suc u)
 &=\sel\bigl(L(\Suc\Suc\Suc u,\Suc k),\Suc u,\Suc\Suc u\bigr)
  =\sel\bigl(L(\Suc\Suc u,k),\Suc u,\Suc\Suc u\bigr)
  =\Suc\beta(k,u), \label{base:eq:beta-conj}\\
 \beta(\Suc\Suc k,0)
 &=\sel\bigl(L(\Suc\Suc 0,\Suc\Suc k),0,1\bigr)
  =\sel\bigl(L(0,k),0,1\bigr)=\sel(1,0,1)=1. \label{base:eq:beta-base}
\end{align}
\emph{Conjugacy.}  $W(\Suc t,\Suc\Suc k)=\Suc W(t,\Suc k)$, by \rulename{EI}
on $t$ with $k$ fixed.  The base is
$W(\Suc 0,\Suc\Suc k)=\beta(\Suc\Suc k,W(0,\Suc\Suc k))
 =\beta(\Suc\Suc k,0)=1=\Suc W(0,\Suc k)$, by \eqref{base:eq:W} and
\eqref{base:eq:beta-base}.  The step is
\begin{align*}
 W(\Suc\Suc t,\Suc\Suc k)
 &=\beta\bigl(\Suc\Suc k,\ W(\Suc t,\Suc\Suc k)\bigr)
  =\beta\bigl(\Suc\Suc k,\ \Suc W(t,\Suc k)\bigr)\\
 &=\Suc\beta\bigl(\Suc k,\ W(t,\Suc k)\bigr)=\Suc W(\Suc t,\Suc k),
\end{align*}
where the second step uses the induction equation and the third is
\eqref{base:eq:beta-conj}.

\emph{Below its cap the counter does not saturate.}  $W(n,\Suc n)=n$, by
\rulename{EI} on $n$.  The base is $W(0,\Suc 0)=0$.  In the step,
$W(\Suc n,\Suc\Suc n)=\Suc W(n,\Suc n)=\Suc n$, by the conjugacy
instantiated at $t=k=n$ and the induction equation.  The instantiation is
made in an identity that is already derived with independent variables; no
parameter of an active induction hypothesis is changed.

\emph{The two equations.}  $P(0)=W(0,0)=0$.  Substituting $\Suc x$ for both
variables in \eqref{base:eq:W},
\begin{align*}
 P(\Suc x)=W(\Suc x,\Suc x)
 &=\beta\bigl(\Suc x,\ W(x,\Suc x)\bigr)=\beta(\Suc x,x)\\
 &=\sel\bigl(L(\Suc\Suc x,\Suc x),\,x,\,\Suc x\bigr)
  =\sel\bigl(L(\Suc x,x),\,x,\,\Suc x\bigr)=\sel(0,x,\Suc x)=x,
\end{align*}
by \cref{base:lem:cmp}.  Finally, from $\Suc x=\Suc y$ congruence gives
$P(\Suc x)=P(\Suc y)$, hence $x=y$.  The transfer to $\Nat$ and $\Teq$ is
\cref{base:lem:E}(iv).
\end{proof}

\begin{remark}[Sources and the use of coverage in the inherited proof]
\label{base:rem:predsources}
The construction (parity, three-cycle, ramp, factor map, comparator,
saturating counter) is the scalar predecessor of the reference statement
filed with \pkg{certified\_zero}.  That proof works in the theory with
numerical cases.  It uses them to show that $Z(n)$ and $N(n)$ take the
values $0$ and $1$ only, and then splits the value $N(n)$ under the
induction hypothesis in the ramp invariant \cref{base:lem:ramp}(i), again in
\cref{base:lem:ramp}(iv), and in $H(3u)=3H(u)$; it also proves
$\sel(c,\Suc a,\Suc b)=\Suc\sel(c,a,b)$ by cases, which is an equational
case.  The splits are replaced here by the equation $Z(Z(N(n)))=N(n)$ and
the two-point rule, applied to identities that are derived in advance with
independent variables, which is the contribution of \pkg{uniform\_raw},
\pkg{newest\_parameter}, \pkg{parameter\_witness} and \pkg{no\_witness}.
The four notes differ in presentation.  \pkg{uniform\_raw} states the
interpolation as $F(Z(c))=\sel(c,F(1),F(0))$, as above;
\pkg{no\_witness} states it as $F(Z(Z(y)))=\sel(y,F(0),F(1))$;
\pkg{parameter\_witness} uses a finite-value transfer lemma for $Z(t)$ and
$N(t)$ proved by equational cases; \pkg{newest\_parameter} instantiates two
identities in $Z(v)$ at $v=Z(N(n))$ and replaces the split in $H(3u)$ by
$\sel(c,a+b,a'+b')=\sel(c,a,a')+\sel(c,b,b')$ and $\sel(c,c,0)=0$.  The
reference defines the counter with the body $\sel(L(\Suc u,m),u,\Suc u)$ and
then identifies the body of $P(\Suc m)$ with it by a separate induction; the
form with $L(\Suc\Suc u,k)$ and the conjugacy between the caps $\Suc\Suc k$
and $\Suc k$, taken from \pkg{no\_witness}, avoids that induction.  The
stronger statement matters where a calculus is closed under substitution
only on its equational part, or restricts the witnesses of \rulename{DI}:
the fragment $\mathcal E$ has no rule for $D$, so the predecessor is
available before any naming of terms.

Both forms are formalized in Lean.  The fragment $\mathcal E$ is the
calculus \Lean{EqPrf} of the modules \Lean{EqKernel}, \Lean{EqArith},
\Lean{EqParity}, \Lean{EqCompare} and \Lean{EqPred}; it has exactly the five
rules of \cref{base:def:E}: hypothesis, the two computations, transport in
one term template, and induction for a literal equation with parameters.
\Cref{base:thm:pred} is \Lean{EqPrf.pred\_zero}, \Lean{EqPrf.pred\_succ}
(in the form \Lean{EqPrf $\Gamma$ (Tm.pred (succ t)) t}, for every term and
every context) and \Lean{EqPrf.succ\_inj}.  The calculus embeds in the
kernel under any restriction of the witness of \rulename{DI}
(\Lean{EqPrf.toDW}), because its derivations contain no introduction of $D$.
The older modules \Lean{TeqParity}, \Lean{TeqCompare} and \Lean{TeqPred}
prove the two equations in $\Teq$ and split the Boolean value by coverage.
\end{remark}

\subsection{A smaller predecessor}
\label{base:sec:pred57}

This subsection draws on \pkg{479\_predecessor57}, which shortens the
three-cycle and the ramp of \cref{base:thm:pred} and reverses the
comparison in the saturating counter, on \pkg{absorbing\_orbit}, which
replaces the factor map by one that sends the odd numbers to $0$ and tests
the orbit itself, on \pkg{absorbing\_predecessor}, which has the same
absorbing factor map and replaces its parity control by a shorter one, on
\pkg{pred50}, which tests the orbit with a different orientation, and on
\pkg{fused\_iteration}, which shortens the three-cycle and fuses the parity
control with the selection of the factor map, and on \pkg{orbit\_equality},
whose diagonal counter lemma shortens the derivation of the $45$-node
predecessor.
The size of the resulting term is measured as follows.

\begin{definition}[Size of a term]
\label{base:def:termsize}
$|x|\defeq|0|\defeq 1$, $|\Suc t|\defeq 1+|t|$ and
$|\It{n}{a}{u}{h}|\defeq 1+|n|+|a|+|h|$.  Thus one node is counted for
every occurrence of a variable, free or bound, of $0$, of $\Suc$ and of the
iterator; the name of the bound variable is not counted.  The size of an
abbreviation is the size of its full expansion, with every repeated
subterm counted at each occurrence.
\end{definition}

Hence $|\sel(c,a,b)|=1+|c|+|a|+|b|$, $|Z(c)|=|c|+4$, $|N(c)|=|c|+7$,
$|a+b|=|a|+|b|+3$ and $|2^{x}|=|x|+8$.  Two three-cycles and ramps
occur.  The \emph{ceiling} pair, of \pkg{479\_predecessor57}, is
\[
 J(u)\defeq\It{u}{2}{w}{Z(w)},\qquad
 M_-(n)\defeq\It{n}{0}{u}{J(u)},\qquad
 G_-(n)\defeq\It{n}{0}{u}{\Suc\sel\bigl(M_-(u),\,\Suc u,\,u\bigr)},
\]
and the \emph{floor} pair, of \pkg{fused\_iteration}, is
\[
 J'(v)\defeq\It{\Suc v}{v}{w}{\sel(w,\Suc v,0)},\qquad
 M'(n)\defeq\It{n}{0}{v}{J'(v)},\qquad
 G'(n)\defeq\It{n}{0}{u}{\Suc\sel\bigl(M'(u),\,u,\,\Suc u\bigr)}.
\]
Put further
\[
 \tilde N(k,a)\defeq\It{k}{0}{v}{\sel(v,a,0)},\qquad
 \tilde N(n)\defeq\tilde N(n,n),\qquad
 H_-(n)\defeq\It{n}{n}{v}{\sel\bigl(v,G'(n),0\bigr)},
\]
\[
 H_{51}(u)\defeq\sel\bigl(N(u),G_-(u),0\bigr),\qquad
 H_{50}(u)\defeq\sel\bigl(\tilde N(u),G_-(u),0\bigr),
\]
and, for $H$ one of $H_-$, $H_{51}$, $H_{50}$,
\[
 O_H(y,a)\defeq\It{y}{a}{u}{H(u)},\qquad
 \Lambda_H(x,y)\defeq Z\bigl(O_H(\Suc y,2^{x})\bigr),
\]
\[
 \beta_H(k,u)\defeq\sel\bigl(O_H(\Suc\Suc u,2^{k}),\,u,\,\Suc u\bigr),\qquad
 W_H(t,k)\defeq\It{t}{0}{u}{\beta_H(k,u)},\qquad
 P_H(m)\defeq W_H(m,m).
\]
The main term is $P_-\defeq P_{H_-}$, and $P_{51}\defeq P_{H_{51}}$,
$P_{50}\defeq P_{H_{50}}$; we write $O$, $\Lambda$, $\beta_-$, $W_-$ for
$O_{H_-}$, $\Lambda_{H_-}$, $\beta_{H_-}$, $W_{H_-}$.  Written out,
\[
 P_-(m)=\It{m}{0}{u}{\sel\bigl(\It{\Suc\Suc u}{2^{m}}{v}{H_-(v)},\,u,\,\Suc u\bigr)},
\]
a term whose only free variable is $m$.

Against \cref{base:thm:pred}: the ceiling three-cycle $M_-$ iterates the zero
test from the seed $2$, and the floor three-cycle $M'$ iterates, from the
seed $v$ and $\Suc v$ times, the update $w\mapsto\sel(w,\Suc v,0)$, instead of
evaluating $Z(u)+2N(u)$; the ramps add $1$ or $2$ according to whether the
state of the three-cycle is zero, instead of adding one more than the state;
the factor maps select $0$ instead of their argument; $H_{50}$ uses the
\emph{weighted parity} $\tilde N$, which is not normalized to $0$ and $1$,
and $H_-$ \emph{fuses} the parity test with the selection: one iteration
seeded by its argument alternates between $0$ and $G'(n)$; and the counter
tests the orbit $O_H(\Suc\Suc u,2^k)$ itself, not its parity, with the two
outputs of the selector in the order of \cref{base:thm:pred}.  On the
natural numbers $J'$ and $M$ map $0\mapsto 1\mapsto 2\mapsto 0$, $M_-$ runs
through $0,2,1,0,\dots$, $G'(n)=G(n)=\lfloor 3n/2\rfloor$ and
$G_-(n)=\lceil 3n/2\rceil$; $\tilde N(n)$ is $0$ at even $n$ and $n$ at odd
$n$; each of $H_-$, $H_{51}$, $H_{50}$ sends $2u$ to $3u$ and every odd
number to $0$; so the orbit $O_H(y,2^x)$ is $3^y2^{x-y}$ for $y\le x$ and
$0$ for $y>x$, and $\Lambda_H(x,y)$ is the indicator of $x\le y$.  These
readings are not used.

\begin{lemma}[Ceiling three-cycle and ramp]
\label{base:lem:ramp57}
$\mathcal E$ derives
\begin{enumerate}[label=\textup{(\roman*)}]
\item $J(0)=2$, $J(1)=0$, $J(2)=1$, $M_-(0)=0$ and $M_-(\Suc n)=J(M_-(n))$;
\item $M_-(G_-(n))=N(n)$;
\item $G_-(0)=0$, $G_-(\Suc n)=\Suc\bigl(G_-(n)+Z(N(n))\bigr)$ and
 $G_-(\Suc\Suc n)=\Suc\Suc\Suc G_-(n)$;
\item $G_-(2n)=3n$;
\item $G_-(3n)+N(n)=3G_-(n)$.
\end{enumerate}
\end{lemma}

\begin{proof}
(i) $J(0)=2$ is \rulename{I0}.  By \rulename{IS} and
\cref{base:lem:interp}(i), $J(1)=Z(J(0))=Z(\Suc 1)=0$ and
$J(2)=Z(J(1))=Z(0)=1$.  The last two equations are \rulename{I0} and
\rulename{IS}.

(ii) As in the proof of \cref{base:lem:ramp}(i), two identities are derived
first, for variables $g$ and $b$, in the context that consists of the one
equation $Z(Z(b))=b$:
\begin{align}
 M_-\bigl(\Suc\sel(b,\Suc g,g)\bigr)
 &=\sel\bigl(b,\ J(J(M_-(g))),\ J(M_-(g))\bigr), \label{base:eq:ramp57-a}\\
 \sel\bigl(b,\ J(J(b)),\ J(b)\bigr)&=Z(b). \label{base:eq:ramp57-b}
\end{align}
Both are instances of the two-point rule.  For \eqref{base:eq:ramp57-a} the
templates are $F(v)\defeq M_-(\Suc\sel(v,\Suc g,g))$ and
$F'(v)\defeq\sel(v,J(J(M_-(g))),J(M_-(g)))$:
$F(0)=M_-(\Suc\Suc g)=J(M_-(\Suc g))=J(J(M_-(g)))=F'(0)$ and
$F(1)=M_-(\Suc g)=J(M_-(g))=F'(1)$, by $\sel(0,a,b)=a$, $\sel(1,a,b)=b$
and (i).  For \eqref{base:eq:ramp57-b} the templates are
$\sel(v,J(J(v)),J(v))$ and $Z(v)$: at $0$ both sides are $J(J(0))=J(2)=1=Z(0)$,
and at $1$ they are $J(1)=0=Z(1)$.  By \cref{base:lem:parity} the hypothesis
is derivable for $b=N(n)$, so both identities hold at $g=G_-(n)$ and
$b=N(n)$ in every context.  Now \rulename{EI} on $n$.  The base is
$M_-(G_-(0))=M_-(0)=0=N(0)$.  In the step the induction equation is
$M_-(G_-(n))=N(n)$, and
\begin{align*}
 M_-(G_-(\Suc n))
 &=M_-\bigl(\Suc\sel(M_-(G_-(n)),\Suc G_-(n),G_-(n))\bigr)
  =M_-\bigl(\Suc\sel(N(n),\Suc G_-(n),G_-(n))\bigr)\\
 &=\sel\bigl(N(n),\ J(J(M_-(G_-(n)))),\ J(M_-(G_-(n)))\bigr)\\
 &=\sel\bigl(N(n),\ J(J(N(n))),\ J(N(n))\bigr)=Z(N(n))=N(\Suc n).
\end{align*}
The second and fourth steps use the induction equation by congruence, the
third is \eqref{base:eq:ramp57-a} and the fifth is \eqref{base:eq:ramp57-b}.

(iii) The first equation is \rulename{I0}.  The two-point rule gives, for
variables $g$ and $b$ in the context $Z(Z(b))=b$,
$\sel(b,\Suc g,g)=g+Z(b)$: at $0$ both sides are $\Suc g=g+1$, and at $1$
both sides are $g=g+0$, by \cref{base:lem:add}.  Instantiating at
$g=G_-(n)$ and $b=N(n)$, \rulename{IS} and (ii) give
\[
 G_-(\Suc n)=\Suc\sel\bigl(M_-(G_-(n)),\Suc G_-(n),G_-(n)\bigr)
 =\Suc\sel\bigl(N(n),\Suc G_-(n),G_-(n)\bigr)=\Suc\bigl(G_-(n)+Z(N(n))\bigr).
\]
Applying it twice, with $N(\Suc n)=Z(N(n))$, $Z(Z(N(n)))=N(n)$ and
$N(n)+Z(N(n))=1$ (\cref{base:lem:parity}),
$G_-(\Suc\Suc n)=\Suc\bigl(\Suc(G_-(n)+Z(N(n)))+N(n)\bigr)
 =\Suc\Suc\bigl(G_-(n)+(N(n)+Z(N(n)))\bigr)=\Suc\Suc(G_-(n)+1)
 =\Suc\Suc\Suc G_-(n)$.

(iv) \rulename{EI}: $G_-(2\cdot 0)=G_-(0)=0=3\cdot 0$, and
$G_-(2(\Suc n))=G_-(\Suc\Suc(2n))=\Suc\Suc\Suc G_-(2n)=\Suc\Suc\Suc(3n)
 =3(\Suc n)$.

(v) \rulename{EI}.  The base is $G_-(3\cdot 0)+N(0)=G_-(0)+0=0=3G_-(0)$.  In
the step write $b\defeq N(n)$ and $z\defeq Z(N(n))$.  By (iii) at
$\Suc\Suc(3n)$ and at $3n$, and by $N(\Suc\Suc(3n))=N(3n)=N(n)$ and
$N(\Suc n)=z$ (\cref{base:lem:parity}),
\[
 G_-(3(\Suc n))+N(\Suc n)
 =G_-\bigl(\Suc\Suc\Suc(3n)\bigr)+z
 =\Suc\bigl(\Suc\Suc\Suc G_-(3n)+z\bigr)+z
 =\Suc\Suc\Suc\Suc\bigl(G_-(3n)+z+z\bigr).
\]
By (iii), \cref{base:lem:add}, the induction equation
$3G_-(n)=G_-(3n)+b$ and $b+z=1$,
\[
 3G_-(\Suc n)=3\,\Suc\bigl(G_-(n)+z\bigr)
 =\Suc\Suc\Suc\bigl(3G_-(n)+3z\bigr)
 =\Suc\Suc\Suc\bigl(G_-(3n)+(b+z)+z+z\bigr)
 =\Suc\Suc\Suc\Suc\bigl(G_-(3n)+z+z\bigr).
\]
No cancellation is used.
\end{proof}

\begin{lemma}[Floor three-cycle and ramp]
\label{base:lem:ramp46}
$\mathcal E$ derives
\begin{enumerate}[label=\textup{(\roman*)}]
\item $J'(0)=1$, $J'(1)=2$, $J'(2)=0$, $M'(0)=0$ and $M'(\Suc n)=J'(M'(n))$;
\item $M'(G'(n))=N(n)$;
\item $G'(0)=0$, $G'(\Suc n)=\Suc\bigl(G'(n)+N(n)\bigr)$ and
 $G'(\Suc\Suc n)=\Suc\Suc\Suc G'(n)$;
\item $G'(2n)=3n$;
\item $G'(3n)=3G'(n)+N(n)$.
\end{enumerate}
\end{lemma}

\begin{proof}
(i) Write $k_c(w)\defeq\sel(w,c,0)$.  By \rulename{IS}, \rulename{I0} and the
selector equations, $J'(0)=k_1(0)=1$;
$J'(1)=k_2(k_2(1))=k_2(0)=2$; and
$J'(2)=k_3(k_3(k_3(2)))=k_3(k_3(0))=k_3(3)=0$.  The last two equations are
\rulename{I0} and \rulename{IS}.

(ii) As in the proof of \cref{base:lem:ramp57}(ii), the two-point rule gives,
for variables $g$ and $b$ in the context $Z(Z(b))=b$,
\[
 M'\bigl(\Suc\sel(b,g,\Suc g)\bigr)=\sel\bigl(b,\ J'(M'(g)),\ J'(J'(M'(g)))\bigr),
 \qquad
 \sel\bigl(b,\ J'(b),\ J'(J'(b))\bigr)=Z(b):
\]
at $b=0$ the first is $M'(\Suc g)=J'(M'(g))$ and at $b=1$ it is
$M'(\Suc\Suc g)=J'(J'(M'(g)))$; the second is $J'(0)=1=Z(0)$ at $0$ and
$J'(J'(1))=J'(2)=0=Z(1)$ at $1$.  \rulename{EI} on $n$: the base is
$M'(G'(0))=M'(0)=0=N(0)$, and in the step, by the induction equation and
the two identities at $g=G'(n)$, $b=N(n)$,
\[
 M'(G'(\Suc n))=M'\bigl(\Suc\sel(N(n),G'(n),\Suc G'(n))\bigr)
 =\sel\bigl(N(n),J'(N(n)),J'(J'(N(n)))\bigr)=Z(N(n))=N(\Suc n).
\]

(iii) The first equation is \rulename{I0}.  The two-point rule gives
$\sel(b,g,\Suc g)=g+b$ in the context $Z(Z(b))=b$: at $0$ both sides are
$g$, and at $1$ both are $\Suc g=g+1$.  With (ii), $G'(\Suc n)=\Suc\sel(N(n),G'(n),\Suc G'(n))=\Suc(G'(n)+N(n))$.
These are the first two equations of \cref{base:lem:ramp}(ii) for $G'$ in
place of $G$.  The proofs of the third equation of \cref{base:lem:ramp}(ii)
and of \cref{base:lem:ramp}(iii) and (iv) use of $G$ only those two
equations, \cref{base:lem:parity} and \cref{base:lem:add}; with $G'$ in place
of $G$ they give the rest of (iii), and (iv) and (v).
\end{proof}

\begin{lemma}[Weighted parity]
\label{base:lem:wparity}
$\mathcal E$ derives, for all terms,
\begin{enumerate}[label=\textup{(\roman*)}]
\item $\sel(c,c,0)=0$ and $\sel(Z(c),a,b)=\sel(c,b,a)$;
\item $\tilde N(k,a)=\sel\bigl(N(k),0,a\bigr)$;
\item $\sel\bigl(\tilde N(n),a,b\bigr)=\sel\bigl(N(n),a,b\bigr)$.
\end{enumerate}
\end{lemma}

\begin{proof}
(i) Equational cases on $c$: $\sel(0,0,0)=0$ and $\sel(\Suc z,\Suc z,0)=0$;
$\sel(Z(0),a,b)=\sel(1,a,b)=b=\sel(0,b,a)$ and
$\sel(Z(\Suc z),a,b)=\sel(0,a,b)=a=\sel(\Suc z,b,a)$.

(ii) \rulename{EI} on $k$ with $a$ fixed.  The base is
$\tilde N(0,a)=0=\sel(0,0,a)=\sel(N(0),0,a)$.  The two-point rule gives, for
variables $a$ and $b$ in the context $Z(Z(b))=b$,
$\sel(\sel(b,0,a),a,0)=\sel(Z(b),0,a)$: at $b=0$ both sides are $a$, since
$\sel(0,a,0)=a=\sel(1,0,a)$, and at $b=1$ both sides are $0$, since
$\sel(\sel(1,0,a),a,0)=\sel(a,a,0)=0$ by (i) and $\sel(Z(1),0,a)=\sel(0,0,a)=0$.
Instantiated at $b=N(k)$, with \cref{base:lem:parity}, it gives the step:
$\tilde N(\Suc k,a)=\sel(\tilde N(k,a),a,0)=\sel(\sel(N(k),0,a),a,0)
 =\sel(Z(N(k)),0,a)=\sel(N(\Suc k),0,a)$, the second step by the induction
equation.

(iii) Equational cases on $n$.  At $0$ both controls are $0$:
$\tilde N(0)=\tilde N(0,0)=0=N(0)$.  At $\Suc z$, (ii) at $k=a=\Suc z$ gives
$\tilde N(\Suc z)=\sel(N(\Suc z),0,\Suc z)$, and the two-point rule gives,
for a variable $q$ in the context $Z(Z(q))=q$,
$\sel(\sel(q,0,\Suc z),a,b)=\sel(q,a,b)$: at $q=0$ both sides are
$\sel(0,a,b)=a$, and at $q=1$ the left side is $\sel(\Suc z,a,b)=b$ and the
right side is $\sel(1,a,b)=b$.  Instantiating at $q=N(\Suc z)$,
$\sel(\tilde N(\Suc z),a,b)=\sel(N(\Suc z),a,b)$.
\end{proof}

The equation $\tilde N(n)=N(n)$ is not claimed and is false on the natural
numbers ($\tilde N(3)=3$, $N(3)=1$); (iii) says that the two terms are
interchangeable as the control of a selector.  The size
$|\tilde N(t)|=2|t|+5$ is smaller than $|N(t)|=|t|+7$ only when $|t|=1$, as
in $H_{50}$.

\begin{lemma}[Absorbing factor maps and comparators]
\label{base:lem:factor57}
$\mathcal E$ derives
\begin{enumerate}[label=\textup{(\roman*)}]
\item $H_-(n)=\sel\bigl(N(n),G'(n),0\bigr)$ and $H_{50}(u)=H_{51}(u)$;
\end{enumerate}
and, for each $H$ among $H_-$, $H_{51}$, $H_{50}$,
\begin{enumerate}[label=\textup{(\roman*)},resume]
\item $H(0)=0$, $H(1)=0$, $H(2u)=3u$ and $H(3u)=3H(u)$;
\item $O_H(\Suc y,a)=O_H(y,H(a))$, $O_H(y,0)=0$, $O_H(y,3a)=3\,O_H(y,a)$ and
 $Z(2^{x})=0$;
\item $\Lambda_H(0,y)=1$, $\Lambda_H(\Suc x,0)=0$,
 $\Lambda_H(\Suc x,\Suc y)=\Lambda_H(x,y)$ and $\Lambda_H(n,n)=1$.
\end{enumerate}
\end{lemma}

\begin{proof}
(i) \emph{Fusion.}  Equational cases on $n$.  At $0$,
$H_-(0)=0=G'(0)=\sel(N(0),G'(0),0)$.  At $\Suc z$, write
$h(v)\defeq\sel(v,G'(\Suc z),0)$, so that $h(\Suc z)=0$.  By rotation
(\cref{base:cor:rot}(i)), \cref{base:lem:wparity}(ii) at $k=z$,
$a=G'(\Suc z)$, and \cref{base:lem:wparity}(i),
\begin{align*}
 H_-(\Suc z)&=\It{\Suc z}{\Suc z}{v}{h(v)}=\It{z}{h(\Suc z)}{v}{h(v)}
  =\It{z}{0}{v}{h(v)}=\tilde N\bigl(z,G'(\Suc z)\bigr)\\
 &=\sel\bigl(N(z),0,G'(\Suc z)\bigr)=\sel\bigl(Z(N(z)),G'(\Suc z),0\bigr)
  =\sel\bigl(N(\Suc z),G'(\Suc z),0\bigr).
\end{align*}
The second equation is \cref{base:lem:wparity}(iii) at $a=G_-(u)$, $b=0$.

(ii) By (i) it suffices to treat $\sel(N(u),G(u),0)$ for $G$ either $G'$ or
$G_-$.  Then $H(0)=\sel(0,G(0),0)=G(0)=0$,
$H(1)=\sel(1,G(1),0)=0$ and $H(2u)=\sel(N(2u),G(2u),0)=\sel(0,3u,0)=3u$, by
\cref{base:lem:parity} and \cref{base:lem:ramp46}(iv), respectively
\cref{base:lem:ramp57}(iv).  For $H_-$, the identity
$\sel(b,3a+b,3c)=3\,\sel(b,a,c)$ of the proof of \cref{base:lem:factor}, at
$b=N(u)$, $a=G'(u)$, $c=0$, with $3\cdot 0=0$, \cref{base:lem:ramp46}(v) and
$N(3u)=N(u)$, gives
\[
 H_-(3u)=\sel\bigl(N(3u),G'(3u),0\bigr)=\sel\bigl(N(u),3G'(u)+N(u),3\cdot 0\bigr)
 =3\,\sel\bigl(N(u),G'(u),0\bigr)=3H_-(u).
\]
For $H_{51}$ the two-point rule gives, for variables $a,c,d,e$ and $b$ in the
context $Z(Z(b))=b$,
\[
 3\,\sel(b,a,c)=\sel(b,3a,3c),\qquad \sel(b,d+b,e)=\sel(b,d,e):
\]
at $b=0$ the sides are $3a$, respectively $d+0=d$, and at $b=1$ they are
$3c$, respectively $e$.  Instantiating at $b=N(u)$, $a=G_-(u)$, $c=0$,
$d=G_-(3u)$ and $e=0$, and using $3\cdot 0=0$, \cref{base:lem:ramp57}(v)
and $N(3u)=N(u)$,
\begin{align*}
 3H_{51}(u)&=3\,\sel\bigl(N(u),G_-(u),0\bigr)=\sel\bigl(N(u),3G_-(u),0\bigr)
  =\sel\bigl(N(u),G_-(3u)+N(u),0\bigr)\\
 &=\sel\bigl(N(u),G_-(3u),0\bigr)=\sel\bigl(N(3u),G_-(3u),0\bigr)=H_{51}(3u).
\end{align*}
The equations for $H_{50}$ follow from those for $H_{51}$ by (i).

(iii) The first equation is rotation (\cref{base:cor:rot}(i)).  The second
and third are equation inductions on $y$ with $a$ fixed: the bases are
\rulename{I0}, and the steps are $O_H(\Suc y,0)=H(O_H(y,0))=H(0)=0$ and
$O_H(\Suc y,3a)=H(O_H(y,3a))=H(3\,O_H(y,a))=3H(O_H(y,a))=3\,O_H(\Suc y,a)$.
Equational cases on $a$ give $Z(2a)=Z(a)$: $Z(2\cdot 0)=Z(0)$, and
$Z(2(\Suc z))=Z(\Suc\Suc(2z))=0=Z(\Suc z)$; likewise $Z(3a)=Z(a)$, with
$3(\Suc z)=\Suc\Suc\Suc(3z)$ (\cref{base:lem:add}).  Then $Z(2^{x})=0$ is an
equation induction: $Z(2^{0})=Z(1)=0$, and $Z(2^{\Suc x})=Z(2\cdot 2^{x})
=Z(2^{x})=0$, since $2^{\Suc x}=2\cdot 2^{x}$ is \rulename{IS}.

(iv) Note $O_H(\Suc 0,a)=H(a)$ by \rulename{IS} and \rulename{I0}.  Now
\begin{align*}
 \Lambda_H(0,y)&=Z\bigl(O_H(\Suc y,1)\bigr)=Z\bigl(O_H(y,H(1))\bigr)=Z\bigl(O_H(y,0)\bigr)=Z(0)=1,\\
 \Lambda_H(\Suc x,0)&=Z\bigl(H(2\cdot 2^{x})\bigr)=Z(3\cdot 2^{x})=Z(2^{x})=0,\\
 \Lambda_H(\Suc x,\Suc y)&=Z\bigl(O_H(\Suc\Suc y,2\cdot 2^{x})\bigr)
  =Z\bigl(O_H(\Suc y,H(2\cdot 2^{x}))\bigr)=Z\bigl(O_H(\Suc y,3\cdot 2^{x})\bigr)\\
 &=Z\bigl(3\,O_H(\Suc y,2^{x})\bigr)=Z\bigl(O_H(\Suc y,2^{x})\bigr)=\Lambda_H(x,y).
\end{align*}
The last equation is an equation induction: $\Lambda_H(0,0)=1$ by the first,
and $\Lambda_H(\Suc n,\Suc n)=\Lambda_H(n,n)$ by the third.
\end{proof}

\begin{proposition}[A predecessor of $46$ nodes]
\label{base:prop:pred57}
\begin{enumerate}[label=\textup{(\roman*)}]
\item $\proves_{\mathcal E}P_-(0)=0$ and $\proves_{\mathcal E}P_-(\Suc x)=x$;
 hence the same sequents are derivable in $\Nat$ and in $\Teq$.  The same
 holds for $P_{51}$ and $P_{50}$.
\item For a variable $m$, $|P_-(m)|=46$, $|P_{50}(m)|=50$ and
 $|P_{51}(m)|=51$, while the predecessor $P(m)$ of \cref{base:thm:pred} has
 $|P(m)|=75$.
\end{enumerate}
\end{proposition}

\begin{proof}
(i) Let $H$ be one of $H_-$, $H_{51}$, $H_{50}$, and write $O$, $\Lambda$,
$\beta_-$, $W_-$, $P_-$ for $O_H$, $\Lambda_H$, $\beta_H$, $W_H$, $P_H$; only
\cref{base:lem:factor57}(iv) is used.  By \rulename{I0} and \rulename{IS},
$W_-(0,k)=0$ and $W_-(\Suc t,k)=\beta_-(k,W_-(t,k))$.  Since
$\Lambda(k,\Suc u)=Z(O(\Suc\Suc u,2^{k}))$,
\cref{base:lem:wparity}(i) gives
\[
 \beta_-(k,u)=\sel\bigl(Z(O(\Suc\Suc u,2^{k})),\Suc u,u\bigr)
 =\sel\bigl(\Lambda(k,\Suc u),\Suc u,u\bigr).
\]
By \cref{base:lem:factor57} and \cref{base:lem:interp}(ii),
\begin{align*}
 \beta_-(\Suc k,\Suc u)
 &=\sel\bigl(\Lambda(\Suc k,\Suc\Suc u),\Suc\Suc u,\Suc u\bigr)
  =\sel\bigl(\Lambda(k,\Suc u),\Suc\Suc u,\Suc u\bigr)=\Suc\beta_-(k,u),\\
 \beta_-(\Suc\Suc k,0)
 &=\sel\bigl(\Lambda(\Suc\Suc k,\Suc 0),1,0\bigr)
  =\sel\bigl(\Lambda(\Suc k,0),1,0\bigr)=\sel(0,1,0)=1.
\end{align*}
\emph{Conjugacy.}  $W_-(\Suc t,\Suc\Suc k)=\Suc W_-(t,\Suc k)$, by
\rulename{EI} on $t$ with $k$ fixed.  The base is
$W_-(\Suc 0,\Suc\Suc k)=\beta_-(\Suc\Suc k,0)=1=\Suc W_-(0,\Suc k)$.  The
step is
\[
 W_-(\Suc\Suc t,\Suc\Suc k)
 =\beta_-\bigl(\Suc\Suc k,\ \Suc W_-(t,\Suc k)\bigr)
 =\Suc\beta_-\bigl(\Suc k,\ W_-(t,\Suc k)\bigr)=\Suc W_-(\Suc t,\Suc k),
\]
by the induction equation and the first displayed identity.

\emph{Below its cap the counter does not saturate.}  $W_-(n,\Suc n)=n$, by
\rulename{EI} on $n$: $W_-(0,\Suc 0)=0$, and
$W_-(\Suc n,\Suc\Suc n)=\Suc W_-(n,\Suc n)=\Suc n$ by the conjugacy,
derived with independent variables and instantiated at $t=k=n$, and the
induction equation.

\emph{The two equations.}  $P_-(0)=W_-(0,0)=0$, and
\begin{align*}
 P_-(\Suc x)=W_-(\Suc x,\Suc x)
 &=\beta_-\bigl(\Suc x,W_-(x,\Suc x)\bigr)=\beta_-(\Suc x,x)
  =\sel\bigl(\Lambda(\Suc x,\Suc x),\Suc x,x\bigr)\\
 &=\sel\bigl(\Lambda(x,x),\Suc x,x\bigr)=\sel(1,\Suc x,x)=x,
\end{align*}
by \cref{base:lem:factor57}.  The transfer is \cref{base:lem:E}(iv).

(ii) By \cref{base:def:termsize}, with variables as arguments,
$|J'(v)|=1+2+1+(1+1+2+1)=9$, $|M'(n)|=12$, $|G'(n)|=20$ (its body has $17$
nodes), $|H_-(n)|=1+1+1+(1+1+20+1)=26$, $|O(y,a)|=|y|+|a|+27$,
$|\Lambda(x,y)|=4+|O(\Suc y,2^{x})|=|x|+|y|+40=42$,
$|\beta_-(k,u)|=1+|O(\Suc\Suc u,2^{k})|+1+2=1+39+3=43$ and
$|P_-(m)|=1+1+1+43=46$.  For the ceiling pair, $|J(u)|=10$, $|M_-(n)|=13$
and $|G_-(n)|=21$; with $|\tilde N(u)|=7$, $|H_{50}(u)|=1+7+21+1=30$ and
$|H_{51}(u)|=1+8+21+1=31$, so that $|O_H(y,a)|=|y|+|a|+|H(u)|+1$ and
$|P_H(m)|=3+1+(3+9+|H(u)|+1)+3=|H(u)|+20$, which is $50$ and $51$.  For the
terms of \cref{base:thm:pred}, $|j(u)|=5+19+3=27$, $|M(n)|=30$, $|G(n)|=38$,
$|H(u)|=48$, $|L(x,y)|=|x|+|y|+64=66$,
$|\beta(k,u)|=1+|L(\Suc\Suc u,k)|+1+2=72$ and $|P(m)|=75$.
\end{proof}

The other terms of the sources are
\begin{align*}
 H_{57}(u)&\defeq\sel\bigl(N(u),G_-(u),u\bigr),\qquad
 L_{57}(x,y)\defeq N\bigl(\It{y}{2^{x}}{u}{H_{57}(u)}\bigr),\\
 P_{57}(m)&\defeq\It{m}{0}{u}{\sel\bigl(L_{57}(m,\Suc u),\,\Suc u,\,u\bigr)}.
\end{align*}
The orientation of \pkg{pred50} runs the orbit the other way.  For $H$ one
of $H_-$, $H_{51}$ put
\[
 K_H(x,y)\defeq Z\bigl(O_H(y,2^{x})\bigr),\qquad
 \beta'_H(k,u)\defeq\sel\bigl(O_H(k,2^{\Suc u}),\,\Suc u,\,u\bigr),
\]
\[
 W'_H(t,k)\defeq\It{t}{0}{u}{\beta'_H(k,u)},\qquad P'_H(m)\defeq W'_H(m,m),
\]
and $P_{45}\defeq P'_{H_-}$, $P'_{50}\defeq P'_{H_{51}}$.  Written out,
\[
 P_{45}(m)=\It{m}{0}{u}{\sel\bigl(\It{m}{2^{\Suc u}}{v}{H_-(v)},\,\Suc u,\,u\bigr)}.
\]
On the natural numbers $K_H(x,y)$ is the indicator of $x<y$.  The term
$P'_{50}$ is the one of \pkg{pred50}; the term $P_{45}$ combines its
orientation with the fused factor map of \pkg{fused\_iteration}.  It was
first assembled in this article; \pkg{abstraction\_and\_refinements} and
\pkg{computation\_leaves} later deliver the same term independently, with the
same reduction to the factor interface $F(0)=F(1)=0$, $F(2a)=3a$,
$F(3a)=3F(a)$ and the reversed comparator.  The derivation below reuses the
arguments of \pkg{pred50} and \pkg{fused\_iteration} for the factor map and
the comparator, and concludes with the following lemma of
\pkg{orbit\_equality} on counters of the form $\It{t}{0}{u}{f(k,u)}$.

\begin{lemma}[Diagonal counter]
\label{base:lem:diag}
Let $f(k,u)$ be a term template in the variables $k$ and $u$, whose other
free variables are distinct from $t$ and $x$, and put
\[
 W_f(t,k)\defeq\It{t}{0}{u}{f(k,u)},\qquad P_f(m)\defeq W_f(m,m).
\]
If $\mathcal E$ derives
\[
 f(1,0)=0,\qquad f(\Suc\Suc k,0)=1,\qquad f(\Suc k,\Suc u)=\Suc f(k,u),
\]
then $\proves_{\mathcal E}P_f(0)=0$ and $\proves_{\mathcal E}P_f(\Suc x)=x$.
\end{lemma}

\begin{proof}
By \rulename{I0} and \rulename{IS}, $W_f(0,k)=0$ and
$W_f(\Suc t,k)=f(k,W_f(t,k))$.  By \cref{base:lem:E}(i) the three
hypotheses may be instantiated at arbitrary terms.

\emph{Conjugacy.}  $W_f(\Suc t,\Suc\Suc k)=\Suc W_f(t,\Suc k)$, by
\rulename{EI} on $t$ with $k$ fixed.  The base is
$W_f(\Suc 0,\Suc\Suc k)=f(\Suc\Suc k,W_f(0,\Suc\Suc k))=f(\Suc\Suc k,0)=1
 =\Suc W_f(0,\Suc k)$.  In the step the induction equation is
$W_f(\Suc t,\Suc\Suc k)=\Suc W_f(t,\Suc k)$, and
\[
 W_f(\Suc\Suc t,\Suc\Suc k)=f\bigl(\Suc\Suc k,W_f(\Suc t,\Suc\Suc k)\bigr)
 =f\bigl(\Suc\Suc k,\Suc W_f(t,\Suc k)\bigr)
 =\Suc f\bigl(\Suc k,W_f(t,\Suc k)\bigr)=\Suc W_f(\Suc t,\Suc k),
\]
by \rulename{IS}, the induction equation by congruence, the third hypothesis
at $k\defeq\Suc k$ and $u\defeq W_f(t,\Suc k)$, and \rulename{IS}.

\emph{The diagonal.}  The conjugacy is derived in the empty context with the
independent variables $t$ and $k$; its instance at $t\defeq\Suc x$ and
$k\defeq x$ is $W_f(\Suc\Suc x,\Suc\Suc x)=\Suc W_f(\Suc x,\Suc x)$, that is
\[
 P_f(\Suc\Suc x)=\Suc P_f(\Suc x).
\]
Also $P_f(1)=W_f(\Suc 0,\Suc 0)=f(1,W_f(0,1))=f(1,0)=0$.  Now
$P_f(\Suc x)=x$ by \rulename{EI} on $x$: the base is $P_f(1)=0$, and in the
step $P_f(\Suc\Suc x)=\Suc P_f(\Suc x)=\Suc x$, by the displayed instance
and the induction equation by congruence.  Finally $P_f(0)=W_f(0,0)=0$ is
\rulename{I0}.  The first induction is complete before it is instantiated,
so no parameter of an active induction hypothesis is changed.
\end{proof}

\begin{proposition}[A predecessor of $45$ nodes]
\label{base:prop:pred45}
\begin{enumerate}[label=\textup{(\roman*)}]
\item For $H$ one of $H_-$, $H_{51}$, $\proves_{\mathcal E}P'_H(0)=0$ and
 $\proves_{\mathcal E}P'_H(\Suc x)=x$; hence the same sequents are derivable
 in $\Nat$ and in $\Teq$.
\item For a variable $m$, $|P_{45}(m)|=45$ and $|P'_{50}(m)|=50$.
\end{enumerate}
\end{proposition}

\begin{proof}
(i) Of the factor map $H$ the argument uses exactly the following
equations, all proved for $H_-$ and for $H_{51}$ in \cref{base:lem:factor57}:
$H(1)=0$ and $H(2a)=3a$ (part (ii)); rotation $O_H(\Suc y,a)=O_H(y,H(a))$,
$O_H(y,0)=0$ and $O_H(y,3a)=3\,O_H(y,a)$ (part (iii), the last two proved
there from $H(0)=0$ and $H(3a)=3H(a)$ of part (ii)).  It uses, besides,
$Z(2^{x})=0$ and $Z(3a)=Z(a)$ (proof of part (iii)), which do not involve
$H$, the selector laws \cref{base:lem:interp}(ii) and
\cref{base:lem:wparity}(i), and the computation equations.  It uses no
property of the parity $N$ and no other property of $H$; in particular the
form $\sel(N(u),\cdot,0)$ of $H_{51}$, and the fusion of
\cref{base:lem:factor57}(i), are not used.  Write $K$, $\beta'$, $W'$, $P'$
for $K_H$, $\beta'_H$, $W'_H$, $P'_H$, and $O$ for $O_H$.

\emph{Comparator.}  By \rulename{I0}, $K(x,0)=Z(2^{x})=0$.  By $2^{0}=1$
(\rulename{I0}) and rotation, $K(0,\Suc y)=Z(O(\Suc y,1))=Z(O(y,H(1)))=Z(O(y,0))=Z(0)=1$.  By
$2^{\Suc x}=2\cdot 2^{x}$ (\rulename{IS}), rotation, $H(2a)=3a$,
$O(y,3a)=3\,O(y,a)$ and $Z(3a)=Z(a)$,
\[
 K(\Suc x,\Suc y)=Z\bigl(O(\Suc y,2\cdot 2^{x})\bigr)
 =Z\bigl(O(y,H(2\cdot 2^{x}))\bigr)=Z\bigl(O(y,3\cdot 2^{x})\bigr)
 =Z\bigl(3\,O(y,2^{x})\bigr)=Z\bigl(O(y,2^{x})\bigr)=K(x,y).
\]
In particular $K(1,1)=K(0,0)=0$, by the third and the first equation.

\emph{Counter.}  Since $K(\Suc u,k)=Z(O(k,2^{\Suc u}))$,
\cref{base:lem:wparity}(i) gives
$\beta'(k,u)=\sel\bigl(Z(O(k,2^{\Suc u})),u,\Suc u\bigr)=\sel\bigl(K(\Suc u,k),u,\Suc u\bigr)$.
Hence, by the comparator equations and \cref{base:lem:interp}(ii),
\begin{align*}
 \beta'(1,0)&=\sel\bigl(K(1,1),0,1\bigr)=\sel(0,0,1)=0,\\
 \beta'(\Suc\Suc k,0)&=\sel\bigl(K(\Suc 0,\Suc\Suc k),0,1\bigr)
  =\sel\bigl(K(0,\Suc k),0,1\bigr)=\sel(1,0,1)=1,\\
 \beta'(\Suc k,\Suc u)&=\sel\bigl(K(\Suc\Suc u,\Suc k),\Suc u,\Suc\Suc u\bigr)
  =\sel\bigl(K(\Suc u,k),\Suc u,\Suc\Suc u\bigr)=\Suc\beta'(k,u).
\end{align*}

\emph{The two equations.}  These are the three hypotheses of
\cref{base:lem:diag} for $f(k,u)\defeq\beta'(k,u)$, for which
$W_f=W'$ and $P_f=P'$.  The lemma gives $\proves_{\mathcal E}P'(0)=0$ and
$\proves_{\mathcal E}P'(\Suc x)=x$.  The transfer is \cref{base:lem:E}(iv).

(ii) By \cref{base:def:termsize}, $|2^{\Suc u}|=1+2+2+5=10$ and
$|O_H(m,2^{\Suc u})|=1+1+10+|H(v)|$.  For $H_-$, with $|H_-(v)|=26$
(\cref{base:prop:pred57}(ii)), this is $38$; the body
$\sel(O_{H_-}(m,2^{\Suc u}),\Suc u,u)$ has $1+38+2+1=42$ nodes; and
$|P_{45}(m)|=1+1+1+42=45$.  For $H_{51}$, with $|H_{51}(v)|=31$, the same
count gives $43$, $47$ and $|P'_{50}(m)|=50$.  The comparators have
$|K_{H_-}(x,y)|=4+1+1+9+26=41$ and $|K_{H_{51}}(x,y)|=46$ nodes.
\end{proof}

\begin{remark}[The diagonal counter; sources]
\label{base:rem:diag}
The route through \cref{base:lem:diag} is that of \pkg{orbit\_equality}.
The derivation that it replaces, which follows \pkg{pred50} and
\pkg{fused\_iteration}, proves the comparator diagonal $K(n,n)=0$ by an
equation induction, the conjugacy as in \cref{base:lem:diag}, the invariant
$W'(n,\Suc n)=n$ below the cap by a second induction, and then
$P'(\Suc x)=\beta'(\Suc x,W'(x,\Suc x))=\beta'(\Suc x,x)
 =\sel(K(x,x),x,\Suc x)=x$.  The lemma needs only the instance
$K(1,1)=0$, which is a computation, and replaces the invariant and the final
calculation by one induction on $P'(\Suc x)=x$: two equation inductions
instead of three.  The term $P_{45}$ and its size are unchanged.  The lemma
applies in the same way to the counters of \cref{base:thm:pred} and
\cref{base:prop:pred57}: by \cref{base:lem:cmp},
$\beta(1,0)=\sel(L(2,1),0,1)=\sel(L(1,0),0,1)=\sel(0,0,1)=0$, and by
\cref{base:lem:factor57}(iv),
$\beta_-(1,0)=\sel(\Lambda(1,1),1,0)=\sel(\Lambda(0,0),1,0)=\sel(1,1,0)=0$;
the other two hypotheses are \eqref{base:eq:beta-conj},
\eqref{base:eq:beta-base} and their counterparts in the proof of
\cref{base:prop:pred57}.  The proofs given there are kept.  The Lean
modules \Lean{TeqPred} and \Lean{EqPred} prove the predecessor equation of
\cref{base:thm:pred} by this route: \Lean{Q.pred\_succ} is one equation
induction on $P(\Suc x)=x$ from the conjugacy \Lean{Q.bump\_sat\_succ} and
$\beta(1,0)=0$, so neither the comparator diagonal nor the invariant below the
cap is formalized there.  \pkg{orbit\_equality} checks
\cref{base:lem:diag} by a symbolic certificate of $37$ equality-proof nodes
with exactly two equation inductions, in the schematic equational theory
whose only axioms are $W(0,n)=0$, $W(\Suc t,n)=F(n,W(t,n))$ and the three
hypotheses of the lemma for $F$ in place of $f$.  The lemma itself, for an
arbitrary template $f$, is not formalized in Lean or Rocq.
\end{remark}

The sizes of the pieces, with variables as arguments, are
\begin{center}
\begin{tabular}{lrrrrrrr}
\hline
term & $P$ & $P_{57}$ & $P_{51}$ & $P_{50}$ & $P'_{50}$ & $P_-$ & $P_{45}$\\
\hline
three-cycle & $30$ & $13$ & $13$ & $13$ & $13$ & $12$ & $12$\\
ramp & $38$ & $21$ & $21$ & $21$ & $21$ & $20$ & $20$\\
parity control of the factor map & $8$ & $8$ & $8$ & $7$ & $8$ & fused & fused\\
factor map & $48$ & $31$ & $31$ & $30$ & $31$ & $26$ & $26$\\
comparator & $66$ & $49$ & $47$ & $46$ & $46$ & $42$ & $41$\\
predecessor & $75$ & $57$ & $51$ & $50$ & $50$ & $46$ & $45$\\
\hline
\end{tabular}
\end{center}
The comparators are $L$, $L_{57}$, $\Lambda_{H_{51}}$, $\Lambda_{H_{50}}$,
$K_{H_{51}}$, $\Lambda$ and $K_{H_-}$.  The orientation of \pkg{pred50} saves
one node against the orientation of $\Lambda$, with either factor map: the
orbit counter $m$ and the seed $2^{\Suc u}$ have $1+10$ nodes against $3+9$
for $\Suc\Suc u$ and $2^{m}$.  In $P_{51}$, $P_{50}$, $P'_{50}$, $P_-$ and
$P_{45}$ the comparator is
not a subterm: the counter tests the orbit without the outer zero test, by
\cref{base:lem:wparity}(i).  Against $P_{57}$ this removes the parity wrapper
$N(\cdot)$ of $L_{57}$ ($7$ nodes) and costs one successor in the orbit;
the weighted parity of $H_{50}$ saves one more node; and against $P_{50}$,
the floor three-cycle saves one node in the clock and so in the ramp, and the
fused factor map $H_-$ saves four nodes against $H_{50}$ ($26$ against
$1+7+20+1=29$ with the floor ramp).  The ceiling three-cycle and ramp alone,
with the counter of \cref{base:thm:pred} ($\sel(L_{57}(\Suc\Suc u,m),u,\Suc u)$
as body), give $58$ nodes; the counter of $P_{57}$ alone, with $L$, gives
$74$.  The saving of $30$ nodes, from $P$ to $P_{45}$, concerns one expanded
term, not the size of a checker or of a derivation.

\begin{proposition}[A lower bound for the parity control]
\label{base:prop:control-bound}
For a term $t$ whose only free variable is $n$ write $t(k)$ for its value in
the natural numbers at $n=k$.
\begin{enumerate}[label=\textup{(\roman*)}]
\item If $t$ contains at most one iterator, then $t(k)$ is nondecreasing on
 the positive numbers $k$.
\item If $c(1)>0$ and $c(2)=0$, then $|c|\ge 7$; the term $\tilde N(n)$
 attains the bound.
\item If $c$ is such a term and $\sel(c,G_-(n),0)(1)=0$ and
 $\sel(c,G_-(n),0)(2)=3$, then $|\sel(c,G_-(n),0)|\ge 30$; the term
 $H_{50}(n)$ attains the bound.
\end{enumerate}
\end{proposition}

\begin{proof}
(i) Above the only iterator of $t$ there are only successors, so $t$ is
$\Suc^{\ell}n$, $\Suc^{\ell}0$ or $\Suc^{\ell}\It{k}{a}{u}{h}$ with $k$, $a$
and $h$ free of iterators.  Then $k$ and $a$ are $\Suc^{i}n$ or $\Suc^{i}0$,
and $h$ is $\Suc^{r}u$, $\Suc^{r}n$ or $\Suc^{r}0$.  If $k$ is $0$, the
value is that of $a$.  Otherwise $k(j)\ge 1$ for every positive $j$, and the
three bodies give the values $a(j)+r\,k(j)$, $j+r$ and $r$.  Each is
nondecreasing, and so is its $\ell$-fold successor.

(ii) A term with $p$ iterators and $s$ successors has $1+2p$ leaves, since
each iterator has three children and each successor one; so
$|t|=1+3p+s$.  A term with $|t|\le 6$ has $p\le 1$, and by (i)
$t(1)\le t(2)$.  Hence $c(1)>0=c(2)$ forces $|c|\ge 7$.  For $\tilde N$:
$|\tilde N(n)|=7$, $\tilde N(1)=\sel(0,1,0)=1$ and
$\tilde N(2)=\sel(\sel(0,2,0),2,0)=\sel(2,2,0)=0$.

(iii) $G_-(1)=2$ and $G_-(2)=3$.  The value at $1$ is $0$ only if $c(1)>0$,
and the value at $2$ is $3$ only if $c(2)=0$.  By (ii) $|c|\ge 7$, so the
size is $1+|c|+21+1\ge 30$.  $H_{50}(n)$ has $30$ nodes, and
$H_{50}(1)=0$, $H_{50}(2)=3$ by \cref{base:lem:factor57} and $2=2\cdot 1$.
\end{proof}

The bound concerns the control alone and the factor map of the form
$\sel(c,G_-(n),0)$ with the ramp $G_-$ fixed; it is not a lower bound for a
factor map of another form or with another ramp, or for a predecessor.  The
fused factor map $H_-$, a single iteration seeded by its argument, is of
another form, and its $26$ nodes do not contradict (iii).

\begin{proposition}[A lower bound for the three-cycle]
\label{base:prop:clock-bound}
With the notation of \cref{base:prop:control-bound}, let $f$ be a term whose
only free variable is $n$, and suppose that its values have a three-cycle
through $0$: $f(0)=p$, $f(p)=q$ and $f(q)=0$ for distinct positive numbers
$p$ and $q$.  Then $|f|\ge 9$.  The term $J'(n)$ attains the bound, and so a
three-cycle of the form $\It{n}{0}{v}{f(v)}$ has at least $12$ nodes,
attained by $M'$.
\end{proposition}

\begin{proof}
By the count $|t|=1+3j+s$, for $j$ iterators and $s$ successors, of the
proof of \cref{base:prop:control-bound}(ii), a term with at most $8$ nodes
has $j\le 2$, and $s\le 1$ if $j=2$.

\emph{At most one iterator.}  Then $f$ is $\Suc^\ell n$, $\Suc^\ell 0$ or
$\Suc^\ell\It{k}{a}{u}{h}$ with $k,a\in\{\Suc^i n,\Suc^i 0\}$ and
$h\in\{\Suc^r u,\Suc^r n,\Suc^r 0\}$.  A constant has no such cycle, since
$f(0)=f(p)$ would give $p=q$.  If $\ell>0$ or $f=n$, then $f(q)>0$.  So
$f=\It{k}{a}{u}{h}$.  If $k$ is the leaf $0$, then $f=a$, treated already.
Otherwise $k(i)\ge 1$ for every positive $i$, and at positive $i$ the body
$\Suc^r u$ gives $a(i)+r\,k(i)$, which vanishes at $q$ only if $a$ is the leaf
$0$ and $r=0$, and then $f$ is identically $0$; the body $\Suc^r n$ gives
$i+r>0$; and the body $\Suc^r 0$ gives the constant $r$, so that
$q=f(p)=r=f(q)=0$.  None has the cycle.

\emph{Two iterators and no successor.}  By induction on terms, a term
without $\Suc$ has the value $0$ when all its free variables have the value
$0$; so $f(0)=0\neq p$.

\emph{Two iterators and one successor.}  Suppose that the successor is in the
body of an iterator $X$, and let $Y$ be the other iterator.  If $X$ lies in
the body of $Y$, then $Y$ is the whole term, its counter has neither a
successor nor a bound variable, so it has the value $0$ at $n=0$, and $Y$
returns its seed, which has no successor: $f(0)=0$.  Otherwise the counter of
$X$ has neither a successor nor a bound variable, so it has the value $0$ at
$n=0$, $X$ returns its seed, which has no successor, and every subterm has
the value $0$ at $n=0$: again $f(0)=0$.  Hence the successor is in no body.
If $f=\Suc t'$, then $f(q)>0$.  So $f=\It{k}{a}{u}{h}$ with the successor in
$k$ or in $a$.  If it is in $k$, then $a$ and $h$ have no successor and
$f(0)=0$ as before.  So the successor is in $a$, and $k$ and $h$ have none.
The second iterator is in $k$, in $a$ or in $h$.
(1) In $k$: then $a\in\{1,\Suc n\}$ and $h\in\{0,n,u\}$, and $k$ is an
iterator over the leaves $0$, $n$ and its own state, whose value at a
positive $i$ is uniformly $0$ or $i$; so at positive $i$ the value of $f$ is
uniformly one of $0$, $i$, $1$, $i+1$, and $f(q)=0$ forces $f(i)=0$ for all
positive $i$, whence $q=f(p)=0$.
(2) In $a$: then $k\in\{0,n\}$ and $h\in\{0,n,u\}$.  If $k=0$ or $h=u$, then
$f=a$, a term with one iterator, treated above.  Otherwise $k=n$,
$h\in\{0,n\}$, and $f(i)$ is uniformly $0$ or $i$ at positive $i$, so that
$q=f(p)$ is $0$ or $p$.
(3) In $h$: then $a\in\{1,\Suc n\}$ and $k\in\{0,n\}$.  If $k=0$, then
$f=a>0$.  If $k=n$, then $f(0)=a(0)=1$, so $p=1$; and $h$ is an iterator
over the leaves $0$, $n$, $u$ and its own state, without successor, whose
value at a positive $n$ and a positive $u$ is uniformly $0$, $n$ or $u$.
Then $f(1)=h(1,a(1))\in\{0,1,a(1)\}$, and $q=f(1)\notin\{0,1\}$ forces $a=\Suc n$
and a body that returns $u$; but then at every positive $i$ the state stays
$i+1$, and $f(q)=q+1\neq 0$.

So no term with at most $8$ nodes has the cycle.  The term
$J'(n)=\It{\Suc n}{n}{w}{\sel(w,\Suc n,0)}$ has $9$ nodes and the cycle
$0\mapsto 1\mapsto 2\mapsto 0$ by \cref{base:lem:ramp46}(i), which holds in
the natural numbers.  The last claim follows from
$|\It{n}{0}{v}{f(v)}|=3+|f|$.
\end{proof}

The bound concerns the update of a three-cycle of this form; it is not a
lower bound for a clock of another form or for a predecessor.

\begin{remark}[Use of the smaller predecessor; sources]
\label{base:rem:pred57}
The term $P_-$ has the two properties of $P$ that the later sections use:
its only free variable is its argument, and it satisfies the two equations
of \cref{base:thm:pred} in $\mathcal E$, hence in every context and at every
term.  The constructions of the rest of this section, and the single-binder
universal of \cref{base:sec:single-binder}, use $P$ only through these
properties, and so do $P_-$ and $P_{45}$, which may replace $P$ in each of
them; with $P_{45}$ every occurrence of a quantified variable in
\cref{base:def:single-binder} costs $|P_{45}(z)|-1=44$ added nodes instead of
$74$ ($45$ with $P_-$).  Moreover, if two terms $p$ and
$p'$ both satisfy the two equations in $\mathcal E$, then equational cases
give $p(x)=p'(x)$, so a replacement inside a term yields a provably equal
term (inside an iterator body by \cref{base:cor:rot}(ii)).  The article keeps
$P$ as the reference term because its derivation is formalized in Lean.

The terms of the table come from five packages: $P_{57}$ from
\pkg{479\_predecessor57}, $P_{51}$ from \pkg{absorbing\_orbit}, $P_{50}$ from
\pkg{absorbing\_predecessor}, $P'_{50}$ from \pkg{pred50}, and $P_-$ from
\pkg{fused\_iteration}; $P_{45}$ is the combination of
\cref{base:prop:pred45}.  $P_{50}$ and $P'_{50}$ are different terms of the
same size: $P'_{50}$ keeps the parity $N$ in its factor map and saves its
node in the comparator, whose orbit runs $m$ steps from the seed
$2^{\Suc u}$ instead of $\Suc\Suc u$ steps from $2^m$.  The derivations of
$P_{51}$ and $P_{50}$ are part of \cref{base:prop:pred57}, and that of
$P'_{50}$ is part of \cref{base:prop:pred45}.  For $P_{57}$, the
two identities in the proof of \cref{base:lem:factor57}(ii), instantiated at
$c=u$ and $e=3u$, give $3H_{57}(u)=H_{57}(3u)$; with
$H_{57}(2u)=\sel(0,3u,2u)=3u$ and $H_{57}(1)=\sel(1,G_-(1),1)=1$ these are the
three equations of \cref{base:lem:factor}, so the proofs of
\cref{base:lem:orbit,base:lem:cmp} give the four comparator equations of
\cref{base:lem:factor57}(iv) for $L_{57}$ in place of $\Lambda_H$, the fourth
by the same equation induction; the proof of \cref{base:prop:pred57}(i) then
applies with $L_{57}(k,\Suc u)$ in place of $\Lambda(k,\Suc u)$, without the
use of \cref{base:lem:wparity}(i).

The packages prove the ramp laws by the same inductions; the main report of
\pkg{479\_predecessor57} splits the Boolean value $N(n)$ in the induction
step, and the detailed notes of all five replace those splits by the
equational two-point argument, as in the proofs above.
\pkg{absorbing\_predecessor} and \pkg{fused\_iteration} prove $H(3u)=3H(u)$
by distributing the selector over a sum and $\sel(c,c,0)=0$ instead of the
identities used here.  \pkg{absorbing\_predecessor} obtains $P_{50}$ from
$P_{51}$ by two applications of the congruence of iterator bodies, using
\cref{base:lem:wparity}(iii).  \pkg{fused\_iteration} states the fusion of
\cref{base:lem:factor57}(i) for every template $A(n)$ with
$\proves_{\mathcal E}A(0)=0$ in place of $G'(n)$, with the same proof.  All
five prove the predecessor equations through a counter of the form
$W(t,m)=\It{t}{0}{u}{\sel(L'(m,u),\Suc u,u)}$ (for \pkg{pred50},
$\sel(K(u,m),u,\Suc u)$), with $L'$ the respective comparator, the shift
$W(\Suc t,\Suc m)=\Suc W(t,m)$, $L'(n,n)=1$ and $W(n,n)=n$, followed by a
further equation induction that identifies the body of the predecessor at
$\Suc m$ with that of $W(\cdot,m)$ (\cref{base:cor:rot}(ii)), and by a
further body congruence for the step to the final orbit test; the counters
$W_H(t,k)$ and the conjugacy between the caps $\Suc\Suc k$ and $\Suc k$ used
above avoid those inductions, as in \cref{base:thm:pred}.
\pkg{absorbing\_predecessor} also proves \cref{base:prop:control-bound}, and
\pkg{fused\_iteration} proves \cref{base:prop:clock-bound}.  None of the new
terms is formalized in Lean or Rocq.
\end{remark}

\subsection{Coverage from a zero-test split}
\label{base:sec:split}

This subsection draws on \pkg{zero\_test\_split}.  With the equational
predecessor available, the coverage axiom (T2) of \cref{base:def:teq} can
be replaced by a disjunction of two equations without an existential
quantifier.  Let $\Teq^{-}$ be the theory with the axioms (T1) and (T4)
of \cref{base:def:teq}, and consider the following schemas, for every
term $n$:
\[
 \mathsf{Cov}:\ n=0\vee\exists m\,(n=\Suc m),\qquad
 \mathsf{Split}:\ n=0\vee Z(n)=0,\qquad
 \mathsf{BR}:\ Z(n)=0\vee Z(n)=1 .
\]
$\mathsf{Cov}$ is (T2), and $\mathsf{Split}$ written out is
$n=0\vee\It{n}{1}{u}{0}=0$.  The proof of \cref{base:lem:E}(iv) derives
\rulename{EI} from (T4) by \cref{base:lem:teqcases}(ii), which does not use
(T2) or (T3), and \rulename{Tr} from the Leibniz rule.  Hence every
sequent derivable in $\mathcal E$ is derivable in $\Teq^{-}$, in particular
the equations of \cref{base:thm:pred}.

\begin{lemma}[Decomposition]
\label{base:lem:dec}
$\mathcal E$ derives $n=\sel\bigl(Z(n),\Suc P(n),0\bigr)$.
\end{lemma}

\begin{proof}
Equational cases (\cref{base:lem:E}(iii)) on $n$.  At $0$ the right side is
$\sel(Z(0),\Suc P(0),0)=\sel(1,\Suc P(0),0)=0$.  At $\Suc x$ it is
$\sel(Z(\Suc x),\Suc P(\Suc x),0)=\sel(0,\Suc P(\Suc x),0)=\Suc P(\Suc x)
=\Suc x$, by \cref{base:thm:pred}.
\end{proof}

\begin{theorem}[Three equivalent case axioms]
\label{base:thm:split}
Over $\Teq^{-}$ the schemas $\mathsf{Cov}$, $\mathsf{Split}$ and
$\mathsf{BR}$ are pairwise interderivable: each instance of one is
derivable from the instances of another at the same term.  Separation is
not used.
\end{theorem}

\begin{proof}
$\mathsf{Cov}\Rightarrow\mathsf{Split}$: eliminate the disjunction.  The
branch $n=0$ is the first disjunct.  In the other branch eliminate the
existential quantifier at a fresh $m$; from $n=\Suc m$ the Leibniz rule
gives $Z(n)=Z(\Suc m)=0$.

$\mathsf{Split}\Rightarrow\mathsf{Cov}$: the branch $n=0$ is the first
disjunct.  In the branch $Z(n)=0$, \cref{base:lem:dec} and the Leibniz rule
give $n=\sel(0,\Suc P(n),0)=\Suc P(n)$, and $\exists$I at the witness $P(n)$
gives the second disjunct.

$\mathsf{Split}\Rightarrow\mathsf{BR}$: in the branch $n=0$, $Z(n)=Z(0)=1$;
the branch $Z(n)=0$ is the first disjunct.

$\mathsf{BR}\Rightarrow\mathsf{Split}$: the branch $Z(n)=0$ is the second
disjunct.  In the branch $Z(n)=1$, \cref{base:lem:dec} gives
$n=\sel(1,\Suc P(n),0)=0$.

Each argument is one elimination of a disjunction, with at most one
elimination of an existential quantifier.
\end{proof}

\begin{corollary}[Coverage replaced by the split]
\label{base:cor:split}
Let $\Teq^{\mathrm{Z}}$ be $\Teq$ with (T2) replaced by $\mathsf{Split}$.
Then $\Teq^{\mathrm{Z}}$ and $\Teq$ have the same theorems, and
\cref{base:thm:bridge} holds with $\Teq^{\mathrm{Z}}$ in place of $\Teq$.
\end{corollary}

\begin{proof}
The two theories share the logic, the language and the axioms (T1), (T3)
and (T4), and so both contain $\Teq^{-}$; they differ only in (T2) and
$\mathsf{Split}$.  By \cref{base:thm:split} each of these is a theorem of
the other theory, so a derivation in one becomes a derivation in the other
by replacing each instance of the differing axiom with its derivation.  The identity translation is therefore
a bi-interpretation of the two theories, and \cref{base:thm:bridge} transfers
verbatim.
\end{proof}

\begin{proposition}[The split is not derivable without a case axiom]
\label{base:prop:split-indep}
The axioms (T1), (T3), (T4) do not derive $\mathsf{Split}$, $\mathsf{Cov}$ or
$\mathsf{BR}$ at a variable, even together with every instance of excluded
middle.  In particular $n=0\vee(n=0\to\bot)$ cannot replace $\mathsf{Split}$.
\end{proposition}

\begin{proof}
Interpret the language in the classical structure with domain
$\mathbb N\times\mathbb N$: $0$ is $(0,0)$, the successor acts on both
coordinates, and the value of a term under an assignment $\rho$ is the pair
of its values in $\mathbb N$ under the two coordinate assignments $\rho_0$
and $\rho_1$; this is well defined and commutes with substitution, because
the values in $\mathbb N$ do.  The computation equations (T1) hold in each
coordinate, and (T3) holds because $(0,0)\neq(1,1)$.  For (T4) fix the
parameters of an equation $E(z)$ and suppose that $E(0)$ holds and that
$E(z)$ implies $E(\Suc z)$ for every $z$.  An induction in $\mathbb N$ gives
$E((k,k))$ for every $k$.  An equation holds at a pair if and only if it
holds in each coordinate, and the coordinate $i$ of $E((a_0,a_1))$ depends
only on $a_i$ and the coordinate $i$ of the parameters.  Hence the
coordinate $0$ of $E((a_0,a_1))$ is that of $E((a_0,a_0))$ and the
coordinate $1$ is that of $E((a_1,a_1))$, and $E((a_0,a_1))$ holds.  The
structure is classical, so it satisfies every instance of excluded middle,
and intuitionistic derivations are sound in it.  At $n=(0,1)$, however,
$n\neq(0,0)$ and $Z(n)=(Z(0),Z(1))=(1,0)$ differs from $(0,0)$ and from
$(1,1)$, so $\mathsf{Split}$ and $\mathsf{BR}$ fail, and $\mathsf{Cov}$ fails
by \cref{base:thm:split}.
\end{proof}

\begin{remark}[Scope of the split; sources]
\label{base:rem:split}
$\mathsf{Split}$ is quantifier free and has no witness; coverage becomes a
theorem with the fixed witness $P(n)$.  Counting one node for each term node
of \cref{base:def:termsize}, each equation, connective and quantifier, the
universal closure of $\mathsf{Cov}$ has $10$ nodes and that of
$\mathsf{Split}$, expanded, has $12$; the existential-free alternative
$\forall n\,(n=0\vee n=\Suc P_{45}(n))$ has $8+45=53$ ($65$ with the term
$P_{57}$ of \cref{base:sec:pred57}, as \pkg{zero\_test\_split} counts it).  The change removes an
existential quantifier, not nodes.  Separation stays an axiom of
$\Teq^{\mathrm{Z}}$ for the reason of \cref{base:rem:sep}: the one-element
structure also satisfies $\mathsf{Split}$.  \pkg{zero\_test\_split} adds that
computation and the Leibniz rule derive $0=1\proves s=t$, by transport into
$\sel(\cdot,s,t)$, and hence, by induction on formulas with the witness $0$
at an existential quantifier, $0=1\proves\varphi$ for every formula $\varphi$
built from equations with $\wedge$, $\vee$, $\to$, $\forall$ and $\exists$;
this does not make (T3) derivable, since $\bot$ is a logical constant.  The
same product structure shows, in the chapter on the kernel rules, that
coverage does not follow from the other scalar principles.  The
equivalence is not formalized in Lean.
\end{remark}

\subsection{Successor laws, subtraction and the reflection of equality}

From here on the derivations are in $\Teq$ and use coverage and separation.
``By cases on $r$'' means the elimination of the disjunction (T2) at $r$
followed by the Leibniz rule; unlike equational cases it may be applied
under hypotheses that mention $r$.

\begin{lemma}[Successor laws]
\label{base:lem:succ}
$\Teq$ proves $\Suc x=\Suc y\to x=y$, $\Suc x=0\to\bot$ and
$x=0\vee x=\Suc P(x)$.
\end{lemma}

\begin{proof}
The first is \cref{base:thm:pred}.  From $\Suc x=0$ congruence gives
$Z(\Suc x)=Z(0)$, that is $0=\Suc 0$, and (T3) gives $\bot$.  For the third,
eliminate (T2): if $x=0$ the first disjunct holds, and if $x=\Suc y$ then
$P(x)=y$ and $x=\Suc P(x)$.
\end{proof}

Put $x\mathbin{\dot-}y\defeq\It{y}{x}{u}{P(u)}$ and
$K(u)\defeq\sel(u,1,P(u))$, and define the \emph{equality test}
\[
 \varepsilon(a,b)\defeq K\bigl(\Suc a\mathbin{\dot-}b\bigr).
\]

\begin{lemma}[Truncated subtraction]
\label{base:lem:monus}
$\Teq$ proves
\begin{enumerate}[label=\textup{(\roman*)}]
\item $x\mathbin{\dot-}0=x$, $x\mathbin{\dot-}\Suc y=P(x\mathbin{\dot-}y)$
 and $\Suc x\mathbin{\dot-}\Suc y=x\mathbin{\dot-}y$;
\item $\Suc a\mathbin{\dot-}a=1$;
\item $\sel\bigl(x\mathbin{\dot-}y,\ x,\ y+(x\mathbin{\dot-}y)\bigr)=x$;
\item $\Suc a\mathbin{\dot-}b=1\leftrightarrow a=b$.
\end{enumerate}
\end{lemma}

\begin{proof}
(i) The first two are computations.  The equation
$P(\Suc x\mathbin{\dot-}y)=x\mathbin{\dot-}y$ holds by induction on $y$:
the base is $P(\Suc x)=x$, and
$P(\Suc x\mathbin{\dot-}\Suc y)=P(P(\Suc x\mathbin{\dot-}y))
 =P(x\mathbin{\dot-}y)=x\mathbin{\dot-}\Suc y$.  Its left side is
$\Suc x\mathbin{\dot-}\Suc y$.

(ii) Induction on $a$: $\Suc 0\mathbin{\dot-}0=1$, and
$\Suc\Suc a\mathbin{\dot-}\Suc a=\Suc a\mathbin{\dot-}a$ by (i).

(iii) Induction on $y$ with $x$ fixed.  The base is
$\sel(x,x,0+x)=\sel(x,x,x)=x$, by \cref{base:lem:interp}(ii).  For the step
write $r\defeq x\mathbin{\dot-}y$; the induction equation is
$\sel(r,x,y+r)=x$ and the goal is $\sel(P(r),x,\Suc y+P(r))=x$.  Argue by
cases on $r$.  If $r=0$, then $P(r)=0$ and the goal is
$\sel(0,x,\dots)=x$.  If $r=\Suc z$, the induction equation gives
$y+\Suc z=x$, and $P(r)=z$, so
$\Suc y+P(r)=\Suc y+z=y+\Suc z=x$ and the goal is $\sel(z,x,x)=x$.

(iv) If $\Suc a\mathbin{\dot-}b=1$, then (iii) at $x=\Suc a$, $y=b$ gives
$\sel(1,\Suc a,b+1)=\Suc a$, that is $\Suc b=\Suc a$, and
\cref{base:lem:succ} gives $a=b$.  The converse is (ii).
\end{proof}

The equation (iii) says, in one literal equation, that a nonzero difference
is an actual difference.

\begin{theorem}[Reflection of equality]
\label{base:thm:reflect}
$\Teq$ proves, for all terms $a$ and $b$,
\[
 \varepsilon(a,b)=0\ \leftrightarrow\ a=b .
\]
Consequently $\Teq\proves a=b\vee(a=b\to\bot)$.
\end{theorem}

\begin{proof}
First, $K(u)=0\leftrightarrow u=1$.  If $u=\Suc 0$ then
$K(u)=P(\Suc 0)=0$.  Conversely assume $K(u)=0$ and argue by cases on $u$.
If $u=0$ then $K(u)=1$, so $\Suc 0=0$, which is contradictory by
\cref{base:lem:succ}.  If $u=\Suc z$ then $K(u)=P(\Suc z)=z$, so $z=0$ and
$u=\Suc 0$.  Hence $\varepsilon(a,b)=0$ is equivalent to
$\Suc a\mathbin{\dot-}b=1$, and to $a=b$ by \cref{base:lem:monus}(iv).  For
the consequence argue by cases on $\varepsilon(a,b)$: if it is $0$ then
$a=b$; if it is $\Suc z$, then $a=b$ would give $\Suc z=0$.
\end{proof}

The theorem concerns the equality of the logic on the whole domain.  It is
not restricted to numerals, and no faithful comparison symbol is assumed.

\begin{remark}[Variants]
\label{base:rem:reflectvariants}
The reference statement uses the sign $\operatorname{sg}(u)\defeq\sel(u,0,1)$,
the invariant
$\sel(\operatorname{sg}(x\mathbin{\dot-}y),x,y+(x\mathbin{\dot-}y))=x$ and
the test $\operatorname{sg}(K(\Suc a\mathbin{\dot-}b))$, which takes the
values $0$ and $1$ only, and it carries the disjunction
``$b_\varphi=0\vee b_\varphi=1$'' through the construction below.  The form
above omits the sign: the selector tests its first argument for zero, so a
two-valued characteristic is not needed.  The Lean module \Lean{TeqReflect}
uses the symmetric test
$\varepsilon'(a,b)\defeq\sel(a\mathbin{\dot-}b,\ b\mathbin{\dot-}a,\ 1)$
instead, whose reflection needs neither $K$ nor \cref{base:lem:monus}(iv).  If
$\varepsilon'(a,b)=0$, then $a\mathbin{\dot-}b=0$, since a successor would
give $1=0$; hence $\varepsilon'(a,b)=b\mathbin{\dot-}a=0$, and
\cref{base:lem:order}(iii) gives $a=a+(b\mathbin{\dot-}a)=b$.  Conversely
$\varepsilon'(a,a)=\sel(0,0,1)=0$ by \cref{base:lem:order}(i).  The
test $\varepsilon$ contains each argument three times and $\varepsilon'$
twice; \cref{base:sec:orbiteq} gives a test that contains each once.
\end{remark}

\subsection{Equality from the absorbing orbit}
\label{base:sec:orbiteq}

This subsection draws on \pkg{orbit\_equality}.  The absorbing factor map
$H_-$ of \cref{base:sec:pred57} gives a second equality test.  With
$O=O_{H_-}$, put
\[
 \delta(a,b)\defeq N\bigl(O(b,2^{a})\bigr),\qquad
 \eta(a,b)\defeq\It{O(b,2^{a})}{1}{v}{Z(v)} .
\]
Each of $a$ and $b$ occurs once in either term.  On the natural numbers the
orbit $O(b,2^{a})$ is $3^{b}2^{a-b}$ for $b\le a$ and $0$ for $b>a$
(\cref{base:sec:pred57}), whose parity is $1$ exactly when $b=a$; so
$\delta$ is the indicator of $a=b$ and $\eta=1-\delta$.  This reading is not
used.

\begin{lemma}[Recurrences of the orbit test]
\label{base:lem:orbiteq}
$\mathcal E$ derives, for all terms,
\begin{enumerate}[label=\textup{(\roman*)}]
\item $\delta(0,0)=1$, $\delta(\Suc a,0)=0$, $\delta(0,\Suc b)=0$ and
 $\delta(\Suc a,\Suc b)=\delta(a,b)$;
\item $\delta(a,a)=1$, $\delta(0,b)=Z(b)$ and $\delta(a,0)=Z(a)$;
\item $\eta(a,b)=Z\bigl(\delta(a,b)\bigr)$.
\end{enumerate}
\end{lemma}

\begin{proof}
(i) The equations $2^{0}=1$, $2^{\Suc a}=2\cdot 2^{a}$ and $O(0,c)=c$ are
\rulename{I0} and \rulename{IS}.  With \cref{base:lem:factor57}(ii) and (iii)
for $H_-$ and \cref{base:lem:parity},
\begin{align*}
 \delta(0,0)&=N\bigl(O(0,1)\bigr)=N(1)=1,\qquad
 \delta(\Suc a,0)=N\bigl(2\cdot 2^{a}\bigr)=0,\\
 \delta(0,\Suc b)&=N\bigl(O(\Suc b,1)\bigr)=N\bigl(O(b,H_-(1))\bigr)
  =N\bigl(O(b,0)\bigr)=N(0)=0,\\
 \delta(\Suc a,\Suc b)&=N\bigl(O(\Suc b,2\cdot 2^{a})\bigr)
  =N\bigl(O(b,H_-(2\cdot 2^{a}))\bigr)=N\bigl(O(b,3\cdot 2^{a})\bigr)\\
 &=N\bigl(3\,O(b,2^{a})\bigr)=N\bigl(O(b,2^{a})\bigr)=\delta(a,b),
\end{align*}
by rotation, $H_-(1)=0$, $O(y,0)=0$, $H_-(2c)=3c$, $O(y,3c)=3\,O(y,c)$,
$N(2n)=0$ and $N(3n)=N(n)$.

(ii) $\delta(a,a)=1$ is an equation induction on $a$: the base is
$\delta(0,0)=1$, and $\delta(\Suc a,\Suc a)=\delta(a,a)=1$ by (i) and the
induction equation.  $\delta(0,b)=Z(b)$ is equational cases on $b$:
$\delta(0,0)=1=Z(0)$ and $\delta(0,\Suc z)=0=Z(\Suc z)$, by (i) and
\cref{base:lem:interp}(i).  Likewise $\delta(a,0)=Z(a)$ by cases on $a$.

(iii) Put $B(q)\defeq\It{q}{1}{v}{Z(v)}$ for a variable $q$.  Then
$B(q)=Z(N(q))$ by \rulename{EI} on $q$: the base is $B(0)=1=Z(0)=Z(N(0))$,
and $B(\Suc q)=Z(B(q))=Z(Z(N(q)))=Z(N(\Suc q))$, by \rulename{IS}, the
induction equation by congruence, and $N(\Suc q)=Z(N(q))$
(\cref{base:lem:parity}).  The instance at $q\defeq O(b,2^{a})$ is
$\eta(a,b)=Z(\delta(a,b))$.
\end{proof}

The derivation uses no predecessor, no order and no equality test; it uses
of $H_-$ only the equations of \cref{base:lem:factor57}(ii) and (iii), and so
holds for $H_{51}$ in place of $H_-$ as well.

\begin{theorem}[Reflection of equality by the orbit test]
\label{base:thm:orbitreflect}
$\Teq$ proves, for all terms $a$ and $b$,
\begin{enumerate}[label=\textup{(\roman*)}]
\item $\delta(a,b)=1\leftrightarrow a=b$ and $\eta(a,b)=0\leftrightarrow a=b$;
\item $\delta(a,b)=0\vee\delta(a,b)=1$ and $\eta(a,b)=0\vee\eta(a,b)=1$.
\end{enumerate}
\end{theorem}

\begin{proof}
The proof uses the predecessor $P$, truncated subtraction and
\cref{base:lem:monus}(i)--(iii); it does not use \cref{base:thm:reflect},
\cref{base:lem:monus}(iv) or induction for compound formulas.  Every
induction is \cref{base:lem:teqcases}(ii) for a literal equation.

\emph{Zero fibre.}  $Z(r)=1\to r=0$.  By cases on $r$: if $r=0$ there is
nothing to prove; if $r=\Suc y$, then $1=Z(r)=Z(\Suc y)=0$, so $0=\Suc 0$
and (T3) gives $\bot$, from which $r=0$ follows.  By
\cref{base:lem:orbiteq}(ii), $\delta(0,r)=1\to r=0$.

\emph{One fibre.}  $\delta(1,r)=1\to r=1$.  By cases on $r$: if $r=0$, then
$\delta(1,0)=0$ by \cref{base:lem:orbiteq}(i), so $1=0$, a contradiction by
(T3); if $r=\Suc s$, then
$1=\delta(\Suc 0,\Suc s)=\delta(0,s)=Z(s)$ by
\cref{base:lem:orbiteq}(i),(ii), so $s=0$ by the zero fibre, and $r=1$.

\emph{Residual invariant.}  For variables $a$, $b$ and $t$,
\[
 \delta(a,b)=1\ \proves_{\Teq}\ \delta\bigl(a\mathbin{\dot-}t,\ b\mathbin{\dot-}t\bigr)=1,
\]
by \cref{base:lem:teqcases}(ii) on $t$ in the context
$\Gamma\defeq\{\delta(a,b)=1\}$, in which $t$ is not free.  The base is
the hypothesis, since $a\mathbin{\dot-}0=a$ and $b\mathbin{\dot-}0=b$.  In
the step write $A\defeq a\mathbin{\dot-}t$ and $B\defeq b\mathbin{\dot-}t$;
the induction equation is $\delta(A,B)=1$, and by
\cref{base:lem:monus}(i) the goal is $\delta(P(A),P(B))=1$.  Argue by cases
on $A$ and on $B$.  If $A=B=0$, then $P(A)=P(B)=0$ and the goal is
$\delta(0,0)=1$.  If $A=0$ and $B=\Suc d$, the induction equation becomes
$\delta(0,\Suc d)=1$, while $\delta(0,\Suc d)=0$, so $0=\Suc 0$, and (T3)
gives $\bot$ and hence the goal; the case $A=\Suc c$, $B=0$ is the same
with $\delta(\Suc c,0)=0$.  If $A=\Suc c$ and $B=\Suc d$, then
$\delta(P(A),P(B))=\delta(c,d)=\delta(\Suc c,\Suc d)=\delta(A,B)=1$, by
\cref{base:thm:pred} and \cref{base:lem:orbiteq}(i).  The induction formula
is the one equation displayed; $a$, $b$ and the hypothesis stay fixed.

(i) Assume $\delta(a,b)=1$ and argue by cases on $a$.  If $a=0$, the
zero fibre gives $b=0=a$.  If $a=\Suc c$, the residual invariant at
$t\defeq c$ gives $\delta(\Suc c\mathbin{\dot-}c,\ b\mathbin{\dot-}c)=1$,
that is $\delta(1,b\mathbin{\dot-}c)=1$ by \cref{base:lem:monus}(ii), and
the one fibre gives $b\mathbin{\dot-}c=1$.  \Cref{base:lem:monus}(iii) at
$x\defeq b$, $y\defeq c$ then gives
\[
 b=\sel\bigl(b\mathbin{\dot-}c,\ b,\ c+(b\mathbin{\dot-}c)\bigr)
  =\sel(1,b,c+1)=c+1=\Suc c=a .
\]
Conversely, $a=b$ gives $\delta(a,b)=\delta(a,a)=1$ by the Leibniz rule and
\cref{base:lem:orbiteq}(ii).

(ii) Put $q\defeq O(b,2^{a})$, so that $\delta(a,b)=N(q)$, and recall
$N(q)+Z(N(q))=1$ (\cref{base:lem:parity}).  By cases on $N(q)$: if
$N(q)=0$, the first disjunct holds; if $N(q)=\Suc y$, then
$Z(N(q))=0$, and $N(q)=N(q)+0=1$.  By \cref{base:lem:orbiteq}(iii),
$\eta(a,b)=Z(\delta(a,b))$ is $Z(0)=1$ in the first case and $Z(1)=0$ in
the second.

The second equivalence of (i) follows: if $a=b$, then $\delta(a,b)=1$ and
$\eta(a,b)=Z(1)=0$; if $\eta(a,b)=0$, then by (ii) either $\delta(a,b)=0$,
whence $\eta(a,b)=Z(0)=1$ and $1=0$, a contradiction by (T3), or
$\delta(a,b)=1$, whence $a=b$.
\end{proof}

Like \cref{base:thm:reflect}, the theorem concerns the equality of the logic
on the whole domain, and it gives $\Teq\proves a=b\vee(a=b\to\bot)$ by
(ii).  It is not a predecessor-free reconstruction: the residual invariant
and the last step use $P$ and \cref{base:lem:monus}.

\begin{corollary}[A predecessor from the orbit test]
\label{base:cor:pdelta}
Put $P_\delta(m)\defeq\It{m}{0}{u}{\sel\bigl(\delta(\Suc u,m),\Suc u,u\bigr)}$.
Then $\proves_{\mathcal E}P_\delta(0)=0$ and
$\proves_{\mathcal E}P_\delta(\Suc x)=x$, and $|P_\delta(m)|=52$.
\end{corollary}

\begin{proof}
Take $f(k,u)\defeq\sel(\delta(\Suc u,k),\Suc u,u)$ in
\cref{base:lem:diag}, so that $P_f=P_\delta$.  By
\cref{base:lem:orbiteq}(i) and \cref{base:lem:interp}(ii),
\begin{align*}
 f(1,0)&=\sel\bigl(\delta(1,1),1,0\bigr)=\sel\bigl(\delta(0,0),1,0\bigr)=\sel(1,1,0)=0,\\
 f(\Suc\Suc k,0)&=\sel\bigl(\delta(1,\Suc\Suc k),1,0\bigr)
  =\sel\bigl(\delta(0,\Suc k),1,0\bigr)=\sel(0,1,0)=1,\\
 f(\Suc k,\Suc u)&=\sel\bigl(\delta(\Suc\Suc u,\Suc k),\Suc\Suc u,\Suc u\bigr)
  =\sel\bigl(\delta(\Suc u,k),\Suc\Suc u,\Suc u\bigr)=\Suc f(k,u).
\end{align*}
The lemma gives the two equations.  By \cref{base:prop:orbitsize}(i),
$|\delta(\Suc u,m)|=42+2+1=45$; the body has $1+45+2+1=49$ nodes and
$|P_\delta(m)|=1+1+1+49=52$.
\end{proof}

The derivation uses only \cref{base:lem:orbiteq}(i), not the reflection of
equality.  The body of $P_\delta$ is that of $P_{45}$ with the parity term
$N$ around the orbit $O(m,2^{\Suc u})$, which accounts for the $7$ further
nodes; $P_\delta$ is larger than $P_{45}$ and is not used below.

\begin{proposition}[Sizes of the equality tests]
\label{base:prop:orbitsize}
In the measure of \cref{base:def:termsize}, for all terms $s$ and $t$:
\begin{enumerate}[label=\textup{(\roman*)}]
\item $|\delta(s,t)|=42+|s|+|t|$ and $|\eta(s,t)|=43+|s|+|t|$; at variables
 $44$ and $45$;
\item with $P_{45}$ in place of $P$, $|\varepsilon(s,t)|=187+3|s|+3|t|$ and
 $|\varepsilon'(s,t)|=95+2|s|+2|t|$; at variables $193$ and $99$;
\item with $P$ itself, $|\varepsilon(s,t)|=307+3|s|+3|t|$ and
 $|\varepsilon'(s,t)|=155+2|s|+2|t|$; at variables $313$ and $159$.
\end{enumerate}
\end{proposition}

\begin{proof}
All bound variables of the abbreviations are fresh, so a size is additive in
the sizes of the arguments.  (i) $|2^{s}|=|s|+8$ and $|H_-(v)|=26$
(\cref{base:prop:pred57}(ii)), so
$|O(t,2^{s})|=1+|t|+(|s|+8)+26=35+|s|+|t|$.  Then
$|\delta(s,t)|=7+|O(t,2^{s})|$, since $|N(c)|=|c|+7$, and
$|\eta(s,t)|=1+|O(t,2^{s})|+|1|+|Z(v)|=1+|O(t,2^{s})|+2+5$.
(ii) The term $P_{45}(m)$ has $45$ nodes and two occurrences of $m$, so
$|P_{45}(c)|=43+2|c|$ and $|P_{45}(u)|=45$.  Hence
$|r|=1+|t|+(1+|s|)+45=47+|s|+|t|$ for $r\defeq\Suc s\mathbin{\dot-}t$, and
$|\varepsilon(s,t)|=|\sel(r,1,P_{45}(r))|=1+|r|+2+43+2|r|=46+3|r|$.
Further $|s\mathbin{\dot-}t|=1+|t|+|s|+45$, and
$|\varepsilon'(s,t)|=1+|s\mathbin{\dot-}t|+|t\mathbin{\dot-}s|+2$.
(iii) The same count with $|P(m)|=75$ (\cref{base:prop:pred57}(ii)), whose
argument also occurs twice, so that $|P(c)|=73+2|c|$.
\end{proof}

\begin{remark}[Use of the orbit test; sources]
\label{base:rem:orbiteq}
By \cref{base:rem:pred57} the predecessor $P_{45}$ may replace $P$ in the
definition of truncated subtraction, and so in $\varepsilon$.
\pkg{orbit\_equality} compares $\eta$ with $\varepsilon$ expanded with
$P_{45}$, the $193$ nodes of \cref{base:prop:orbitsize}(ii): $\eta$ saves
$148$ nodes, and contains each argument once instead of three times.  The
figure $193$ is not stated elsewhere in the article; it is the expansion of
the test $\varepsilon$ defined before \cref{base:lem:monus}.  Against the
symmetric test $\varepsilon'$ of \cref{base:rem:reflectvariants}, also
expanded with $P_{45}$, the saving is $54$ nodes.
% ed.: the sizes of varepsilon' and the sizes with P are computed here; the package compares eta only with varepsilon expanded with P_45.
Neither $45$ nor any other of these sizes is claimed to be minimal, and the
counts concern one expanded term, not a checker or a derivation.

In \cref{base:def:chi} the atomic clause may be
$\chi_{s=t}\defeq\eta(s,t)$ instead of $\varepsilon(s,t)$; the proof of
\cref{base:prop:chi} is unchanged, its atomic case being
\cref{base:thm:orbitreflect}(i), and so is that of \cref{base:thm:qfi}.  No
axiom, rule, domain or interpretation changes, so \cref{base:thm:bridge}
holds verbatim with either atomic clause.  The order of the construction is
kept: the recurrences of \cref{base:lem:orbiteq} are available after
\cref{base:lem:factor57}, and the reflection after
\cref{base:lem:monus}(iii).  The checkers delivered with
\pkg{orbit\_equality} implement the inherited guarded calculus and do not use
$\eta$; they are recorded with the other checkers in the chapter on
certificates.  None of the terms of this subsection is formalized in Lean or
Rocq.
\end{remark}

\subsection{Induction for quantifier-free formulas}

\begin{definition}[Characteristic term]
\label{base:def:chi}
For a quantifier-free ordinary formula $\varphi$ the term $\chi_\varphi$ is
\begin{gather*}
 \chi_{s=t}\defeq\varepsilon(s,t),\qquad \chi_\bot\defeq 1,\qquad
 \chi_{\varphi\wedge\psi}\defeq\sel(\chi_\varphi,\chi_\psi,1),\\
 \chi_{\varphi\vee\psi}\defeq\sel(\chi_\varphi,0,\chi_\psi),\qquad
 \chi_{\varphi\to\psi}\defeq\sel(\chi_\varphi,\chi_\psi,0).
\end{gather*}
\end{definition}

The definition is a recursion on finite syntax, and
$\chi_{\varphi[\theta]}=\chi_\varphi[\theta]$.

\begin{proposition}[Reflection of quantifier-free formulas]
\label{base:prop:chi}
For every quantifier-free $\varphi$,
$\Teq\proves\chi_\varphi=0\leftrightarrow\varphi$.
\end{proposition}

\begin{proof}
By induction on $\varphi$.  An equation is \cref{base:thm:reflect}.  For
$\bot$: $\Suc 0=0$ implies $\bot$ by \cref{base:lem:succ}, and $\bot$
implies everything.

Conjunction.  Assume $\sel(\chi_\varphi,\chi_\psi,1)=0$ and argue by cases on
$\chi_\varphi$.  If $\chi_\varphi=0$ the hypothesis is $\chi_\psi=0$, and
the induction hypotheses give $\varphi$ and $\psi$.  If $\chi_\varphi$ is a
successor the hypothesis is $1=0$, a contradiction.  Conversely, $\varphi$
and $\psi$ give $\chi_\varphi=0$ and $\chi_\psi=0$, hence
$\sel(0,\chi_\psi,1)=\chi_\psi=0$.

Disjunction.  Assume $\sel(\chi_\varphi,0,\chi_\psi)=0$ and argue by cases
on $\chi_\varphi$.  If it is $0$, then $\varphi$.  If it is a successor, the
hypothesis is $\chi_\psi=0$, hence $\psi$.  Conversely, $\varphi$ gives
$\sel(0,0,\chi_\psi)=0$, and $\psi$ gives $\sel(\chi_\varphi,0,0)=0$ by
\cref{base:lem:interp}(ii).

Implication.  Assume $\sel(\chi_\varphi,\chi_\psi,0)=0$ and $\varphi$.  Then
$\chi_\varphi=0$, the hypothesis is $\chi_\psi=0$, and $\psi$ follows.
Conversely assume $\varphi\to\psi$ and argue by cases on $\chi_\varphi$.  If
it is $0$, then $\varphi$, hence $\psi$, hence $\chi_\psi=0$ and
$\sel(0,\chi_\psi,0)=0$.  If it is a successor, the selector is $0$.
\end{proof}

No excluded middle for arbitrary formulas is used: the case distinctions are
on terms.

\begin{theorem}[Quantifier-free induction]
\label{base:thm:qfi}
For every quantifier-free $\varphi(z)$,
\[
 \Teq\proves\varphi(0)\to\bigl(\forall z\,(\varphi(z)\to\varphi(\Suc z))
 \to\forall z\,\varphi(z)\bigr).
\]
Equivalently, if $\Gamma\proves_{\Teq}\varphi(0)$ and
$\Gamma,\varphi(z)\proves_{\Teq}\varphi(\Suc z)$ with $z$ not free in
$\Gamma$, then $\Gamma\proves_{\Teq}\varphi(t)$ for every $t$.  The context
$\Gamma$ may contain quantified formulas.
\end{theorem}

\begin{proof}
Assume the two antecedents.  By \cref{base:prop:chi} the first gives
$\chi_\varphi(0)=0$.  Under the further assumption $\chi_\varphi(z)=0$ the
proposition gives $\varphi(z)$, the second antecedent gives
$\varphi(\Suc z)$, and the proposition gives $\chi_\varphi(\Suc z)=0$.
\Cref{base:lem:teqcases}(ii) for the equation $\chi_\varphi(z)=0$ gives
$\chi_\varphi(z)=0$ for all $z$, and the proposition gives $\varphi(z)$.
\end{proof}

In Lean this is \Lean{Teq.qfInd}.

\subsection{Order, powers of two and pairing}

Write $x\le y$ for the equation $x\mathbin{\dot-}y=0$.

\begin{lemma}[Order]
\label{base:lem:order}
$\Teq$ proves
\begin{enumerate}[label=\textup{(\roman*)}]
\item $x\mathbin{\dot-}x=0$, $(a+b)\mathbin{\dot-}a=b$,
 $a\mathbin{\dot-}(a+b)=0$ and $\Suc(x+d)\mathbin{\dot-}x=\Suc d$;
\item $\sel\bigl(x\mathbin{\dot-}y,\ x+(y\mathbin{\dot-}x),\ y\bigr)=y$;
\item $x\le y\to x+(y\mathbin{\dot-}x)=y$;
\item $x\le y\to x\le y+z$;\quad $\Suc x\le y\to x\le y$;\quad
 $x\le y\wedge y\le z\to x\le z$.
\end{enumerate}
\end{lemma}

\begin{proof}
(i) By induction: $0\mathbin{\dot-}0=0$ and
$\Suc x\mathbin{\dot-}\Suc x=x\mathbin{\dot-}x$;
$(0+b)\mathbin{\dot-}0=b$ and
$(\Suc a+b)\mathbin{\dot-}\Suc a=\Suc(a+b)\mathbin{\dot-}\Suc a
 =(a+b)\mathbin{\dot-}a$, by induction on $a$;
$a\mathbin{\dot-}(a+0)=a\mathbin{\dot-}a=0$ and
$a\mathbin{\dot-}(a+\Suc b)=P(a\mathbin{\dot-}(a+b))=P(0)=0$, by induction
on $b$.  The fourth equation is the second for $x+\Suc d$.

(ii) Induction on $y$ with $x$ fixed.  The base
$\sel(x,x+(0\mathbin{\dot-}x),0)=0$ holds by cases on $x$: at $0$ the left
side is $0+(0\mathbin{\dot-}0)=0$, and at a successor it is $0$.  For the
step write $r\defeq x\mathbin{\dot-}y$; the induction equation is
$\sel(r,x+(y\mathbin{\dot-}x),y)=y$ and the goal is
\[
 \sel\bigl(P(r),\ x+(\Suc y\mathbin{\dot-}x),\ \Suc y\bigr)=\Suc y .
\]
Argue by cases on $r$.  If $r=0$, the induction equation gives
$x+(y\mathbin{\dot-}x)=y$.  By (i),
$\Suc y\mathbin{\dot-}x=\Suc(x+(y\mathbin{\dot-}x))\mathbin{\dot-}x
 =\Suc(y\mathbin{\dot-}x)$, so
$x+(\Suc y\mathbin{\dot-}x)=\Suc(x+(y\mathbin{\dot-}x))=\Suc y$, which is
the goal since $P(r)=0$.  If $r=\Suc z$, then $P(r)=z$ and
\cref{base:lem:monus}(iii) gives $y+\Suc z=x$.  Argue by cases on $z$.  If
$z=0$, then $x=\Suc y$ and the goal is
$\Suc y+(\Suc y\mathbin{\dot-}\Suc y)=\Suc y$, which holds by (i).  If $z$
is a successor, the selector returns $\Suc y$.

(iii) is (ii) with $x\mathbin{\dot-}y=0$.

(iv) By (iii), $y+z=x+((y\mathbin{\dot-}x)+z)$, and (i) applies.  If
$\Suc x\le y$, then $y=\Suc x+d=x+\Suc d$ with
$d\defeq y\mathbin{\dot-}\Suc x$, and (i) applies.  If $x\le y$ and
$y\le z$, then $z=y+(z\mathbin{\dot-}y)$ and the first part applies.
\end{proof}

\begin{remark}[The route of the Lean module]
\label{base:rem:orderlean}
The module \Lean{TeqOrder} proves (iii) through the maximum identity
$x+(y\mathbin{\dot-}x)=y+(x\mathbin{\dot-}y)$ instead of (ii): equation
induction on $y$, with base $x+0=0+x$ by $0\mathbin{\dot-}x=0$, and a step
by cases on $x\mathbin{\dot-}y$ that uses only (i) and
\cref{base:lem:monus}(iii); then $x\le y$ gives $x+(y\mathbin{\dot-}x)=y$
at once.  Of the transitivity laws in (iv) it needs only
$x+y\le k\to x\le k$, and it does not use the monotonicity of $2^{x}$ in
\cref{base:lem:pow}: the logarithm counter of \cref{base:thm:pair} is
proved from the invariant ``$2^{i}\le k$ implies $\mathrm{lgc}(i,k)=i$'',
whose step needs only $2^{i}\le 2^{i}+2^{i}\le k$.
\end{remark}

\begin{lemma}[Powers of two]
\label{base:lem:pow}
$\Teq$ proves $2^{0}=1$, $2^{\Suc x}=2^{x}+2^{x}$, $2^{x}\le 2^{x+d}$,
$x\le y\to 2^{x}\le 2^{y}$ and $\Suc x\le 2^{x}$.
\end{lemma}

\begin{proof}
The first two are computations.  $2^{x}\le 2^{x+d}$ by induction on $d$:
the base is \cref{base:lem:order}(i), and
$2^{x+\Suc d}=2^{x+d}+2^{x+d}$, so the step is \cref{base:lem:order}(iv).
Monotonicity follows by $y=x+(y\mathbin{\dot-}x)$.  $\Suc x\le 2^{x}$ by
induction on $x$: the base is $1\mathbin{\dot-}1=0$.  In the step let
$p\defeq 2^{x}$ and $d\defeq p\mathbin{\dot-}\Suc x$, so that
$p=\Suc x+d$ by the induction equation and \cref{base:lem:order}(iii).
Then $p+p=\Suc\Suc x+(x+d+d)$ by \cref{base:lem:add}, and
$\Suc\Suc x\le p+p$ by \cref{base:lem:order}(i).
\end{proof}

Put
\begin{gather*}
 \ell(k,v)\defeq\sel\bigl(2^{\Suc v}\mathbin{\dot-}k,\ \Suc v,\ v\bigr),
 \qquad
 \Lambda(t,k)\defeq\It{t}{0}{v}{\ell(k,v)},\qquad \lg(w)\defeq\Lambda(w,w),\\
 \langle a,b\rangle\defeq 2^{a+b}+a,\qquad
 \operatorname{fst}(w)\defeq w\mathbin{\dot-}2^{\lg(w)},\qquad
 \operatorname{snd}(w)\defeq\lg(w)\mathbin{\dot-}\operatorname{fst}(w).
\end{gather*}
The step $\ell(k,v)$ increments $v$ while $2^{v+1}\le k$, so $\Lambda(t,k)$
is a saturating counter for the binary logarithm of $k$.

\begin{theorem}[Pairing]
\label{base:thm:pair}
$\Teq$ proves $\lg\langle a,b\rangle=a+b$,
$\operatorname{fst}\langle a,b\rangle=a$ and
$\operatorname{snd}\langle a,b\rangle=b$.
\end{theorem}

\begin{proof}
Fix $a$, $b$ and put $c\defeq a+b$ and $k\defeq 2^{c}+a$.

\emph{Below the exponent the counter follows its index:}
$\sel\bigl(i\mathbin{\dot-}c,\ \Lambda(i,k),\ i\bigr)=i$, by induction on
$i$.  The base is $\sel(0\mathbin{\dot-}c,0,0)=0$.  In the step argue by
cases on $\Suc i\mathbin{\dot-}c$.  If it is a successor, the goal
$\sel(\Suc i\mathbin{\dot-}c,\Lambda(\Suc i,k),\Suc i)=\Suc i$ is a
computation.  If it is $0$, then $\Suc i\le c$, hence $i\le c$ by
\cref{base:lem:order}(iv), and the induction equation gives
$\Lambda(i,k)=i$.  Moreover $2^{\Suc i}\le 2^{c}\le k$ by
\cref{base:lem:pow,base:lem:order}, so
$\Lambda(\Suc i,k)=\ell(k,i)=\sel(0,\Suc i,i)=\Suc i$, which is the goal.
At $i=c$ this gives
\begin{equation}
 \Lambda(c,k)=c. \label{base:eq:lg-low}
\end{equation}

\emph{At the exponent the counter stops:} $\ell(k,c)=c$.  By
\cref{base:lem:pow}, $2^{c}=\Suc c+d$ with
$d\defeq 2^{c}\mathbin{\dot-}\Suc c$, hence
$2^{\Suc c}=2^{c}+2^{c}=(2^{c}+a)+\Suc(b+d)$ and
$2^{\Suc c}\mathbin{\dot-}k=\Suc(b+d)$ by \cref{base:lem:order}(i).  The
selector in $\ell(k,c)$ therefore returns $c$.

\emph{Beyond the exponent:} $\Lambda(c+j,k)=c$, by induction on $j$.  The
base is \eqref{base:eq:lg-low}, and
$\Lambda(\Suc(c+j),k)=\ell(k,\Lambda(c+j,k))=\ell(k,c)=c$.

\emph{Conclusion.}  $\Suc c\le 2^{c}$ gives $c\le 2^{c}\le k$, so
$k=c+(k\mathbin{\dot-}c)$ and $\lg(k)=\Lambda(k,k)=c$.  Then
$\operatorname{fst}(k)=(2^{c}+a)\mathbin{\dot-}2^{c}=a$ and
$\operatorname{snd}(k)=(a+b)\mathbin{\dot-}a=b$ by
\cref{base:lem:order}(i).
\end{proof}

Surjectivity of the pairing is not needed and not claimed.  The numerical
domain is never replaced by the range of the coding; a pair only stores the
state of one recursion.

\begin{remark}[The pairing of the reference statement]
\label{base:rem:sqrtpair}
The reference statement filed with \pkg{certified\_zero} uses the pairing
$(a+b)^{2}+a$, with multiplication $a\cdot b\defeq\It{b}{0}{u}{u+a}$ and a
bounded search for the integer square root,
$q(n)\defeq\It{n}{0}{v}{\sel((\Suc v)^{2}\mathbin{\dot-}n,\Suc v,v)}$, whose
invariants are proved by quantifier-free induction; it appeals to ``the
usual linear-order and semiring identities'' without deriving them.  The
pairing above is the one of the Lean modules \Lean{TeqOrder} and
\Lean{TeqPair}.  It needs no multiplication, and all its inductions are on
literal equations, so \cref{base:thm:qfi} is not used for it.  Either
pairing serves the next theorem.
\end{remark}

\subsection{Primitive recursion}

\begin{definition}[The derived recursor]
\label{base:def:rec}
Let $a$ be a term and $g(y,v)$ a term with two distinguished variables.
With a fresh variable $w$ put
\[
 s_g(w)\defeq\bigl\langle\Suc\operatorname{fst}(w),\
   g(\operatorname{fst}(w),\operatorname{snd}(w))\bigr\rangle,\qquad
 R(c,a;\,y\,v.\,g)\defeq
   \operatorname{snd}\bigl(\It{c}{\langle 0,a\rangle}{w}{s_g(w)}\bigr).
\]
The recursor binds $y$ and $v$ in $g$.
\end{definition}

The iterator carries the pair of the counter and the value.  For a function
symbol $f$ of $\iPRA$ of arity $k$ the term $\ulcorner f\urcorner$, with free
variables among the fixed variables $x_1,\dots,x_k$, is defined by recursion
on $f$; $\ulcorner f\urcorner(t_1,\dots,t_k)$ abbreviates
$\ulcorner f\urcorner\subst{\vec t}{\vec x}$.
\begin{gather*}
 \ulcorner\mathrm{zero}\urcorner\defeq 0,\qquad
 \ulcorner\mathrm{succ}\urcorner\defeq\Suc x_1,\qquad
 \ulcorner\pi^k_i\urcorner\defeq x_i,\qquad
 \ulcorner\mathrm{comp}(f;\vec g)\urcorner\defeq
   \ulcorner f\urcorner\bigl(\ulcorner g_1\urcorner,\dots,
   \ulcorner g_m\urcorner\bigr),\\
 \ulcorner\mathrm{prec}(f,g)\urcorner\defeq
   R\bigl(x_1,\ \ulcorner f\urcorner(x_2,\dots,x_{k+1});\
     y\,v.\,\ulcorner g\urcorner(y,v,x_2,\dots,x_{k+1})\bigr).
\end{gather*}
The translation $\sigma$ of the terms of $\iPRA$ into the terms of $\Teq$ is
\[
 \sigma(x)\defeq x,\qquad
 \sigma\bigl(f(t_1,\dots,t_k)\bigr)\defeq
   \ulcorner f\urcorner\bigl(\sigma(t_1),\dots,\sigma(t_k)\bigr).
\]
It commutes with substitution literally,
$\sigma(t[\theta])=\sigma(t)[\sigma\circ\theta]$, and
$\sigma(0)=0$, $\sigma(\Suc t)=\Suc\sigma(t)$.

\begin{theorem}[Primitive recursion]
\label{base:thm:primrec}
\begin{enumerate}[label=\textup{(\roman*)}]
\item $\Teq$ proves $R(0,a;\,y\,v.\,g)=a$ and
 $R(\Suc c,a;\,y\,v.\,g)=g\bigl(c,\ R(c,a;\,y\,v.\,g)\bigr)$.
\item For every defining equation $s=t$ of a function symbol of $\iPRA$
 (\cref{base:def:ipra}), $\Teq\proves\sigma(s)=\sigma(t)$.
\end{enumerate}
\end{theorem}

\begin{proof}
(i) Write $A(c)\defeq\It{c}{\langle 0,a\rangle}{w}{s_g(w)}$.  By induction
on $c$, $\operatorname{fst}(A(c))=c$: the base is
$\operatorname{fst}\langle 0,a\rangle=0$, and
$\operatorname{fst}(A(\Suc c))=\operatorname{fst}(s_g(A(c)))
 =\Suc\operatorname{fst}(A(c))=\Suc c$, by \cref{base:thm:pair}.  Then
$R(0,a;\,y\,v.\,g)=\operatorname{snd}\langle 0,a\rangle=a$, and
\[
 R(\Suc c,a;\,y\,v.\,g)=\operatorname{snd}\bigl(s_g(A(c))\bigr)
 =g\bigl(\operatorname{fst}(A(c)),\operatorname{snd}(A(c))\bigr)
 =g\bigl(c,\ R(c,a;\,y\,v.\,g)\bigr).
\]
(ii) For a projection, $\sigma(\pi^k_i(\vec x))$ is the term $x_i$.  For a
composition, $\sigma(\mathrm{comp}(f;\vec g)(\vec x))$ and
$\sigma(f(g_1(\vec x),\dots,g_m(\vec x)))$ are the same term.  For a
recursion the two equations are (i) for the seed
$\ulcorner f\urcorner(\vec x)$ and the step
$\ulcorner g\urcorner(y,v,\vec x)$.
\end{proof}

The construction is an external recursion on each finite description.  It
is not an internal quantification over functions, and no universal
evaluator is introduced.  In Lean the recursor is the content of the module
\Lean{TeqRec}.

% ---------------------------------------------------------------------
\section{The bridge}
\label{base:sec:bridge}

This section draws on the reference statements filed with
\pkg{certified\_zero} (the inherited basis, the inherited ordinary bridge and
the section on primitive recursion) and follows the Lean modules
\Lean{NativeTranslate}, \Lean{NativeComposite}, \Lean{NativeRoundTrip},
\Lean{ScalarToPRA}, \Lean{PRAToScalar}, \Lean{RoundTripScalar} and
\Lean{RoundTripPRA} in the organization of the proofs.  The bridge has two
spans: between $\Nat$ and $\Teq$ the terms are the same and the formulas are
translated; between $\Teq$ and $\iPRA$ the logic is the same and the terms
are translated.

\subsection{The native core and \texorpdfstring{$\Teq$}{Teq}}

\begin{definition}[Native reading of an ordinary formula]
\label{base:def:star}
For an ordinary formula $\varphi$ over the terms of \cref{base:def:terms},
the native formula $\varphi^{*}$ replaces every connective by its native
abbreviation (\cref{base:def:macros}):
\begin{gather*}
 (s=t)^{*}\defeq(s=t),\qquad \bot^{*}\defeq\bot_*,\qquad
 (\varphi\wedge\psi)^{*}\defeq\varphi^{*}\wedge_*\psi^{*},\qquad
 (\varphi\vee\psi)^{*}\defeq\varphi^{*}\vee_*\psi^{*},\\
 (\varphi\to\psi)^{*}\defeq\varphi^{*}\Rightarrow_*\psi^{*},\qquad
 (\forall x\,\varphi)^{*}\defeq\forall_* x\,\varphi^{*},\qquad
 (\exists x\,\varphi)^{*}\defeq\exists_* x\,\varphi^{*}.
\end{gather*}
\end{definition}

By \eqref{base:eq:sel-subst}, $(\varphi[\theta])^{*}=\varphi^{*}[\theta]$.

\begin{proposition}[Soundness of the ordinary reading]
\label{base:prop:sound}
If $\Gamma\proves_{\Nat}A$, then $\Gamma^\circ\proves_{\Teq}A^\circ$.
\end{proposition}

\begin{proof}
By induction on the derivation, using \eqref{base:eq:reading-subst} for the
instances.  \rulename{Hyp} is hypothesis; \rulename{I0} and \rulename{IS}
are (T1).  \rulename{Tr}: from $s=t$ and $h(s)=r$ the equality rules give
$r=h(t)$.  \rulename{QE}: the premises read $B^\circ(0)\to\forall y\,B^\circ$
and $B^\circ(0)$, which give $B^\circ(t)$.  \rulename{DE}: the first premise
reads $\exists x\,(B^\circ(x,0)\to\forall y\,B^\circ(x,y))$ and the
hypothesis $\Q{y}B(b,y)$ of the second reads
$B^\circ(b,0)\to\forall y\,B^\circ(b,y)$, so the rule is existential
elimination.  \rulename{DI}: the premise reads
$\Gamma^\circ,B^\circ(t,0)\proves B^\circ(t,\Suc z)$.  Over
$\Gamma^\circ,B^\circ(t,0)$, \cref{base:lem:teqcases}(i) for the formula
$B^\circ(t,y)$ gives $\forall y\,B^\circ(t,y)$; discharge $B^\circ(t,0)$ and
introduce the existential quantifier at $t$.  \rulename{QI}: the premise
reads $\Gamma^\circ,E(z)\proves E(\Suc z)$; over $\Gamma^\circ,E(0)$,
\cref{base:lem:teqcases}(ii) gives $E(y)$ for a fresh $y$; generalize and
discharge $E(0)$.
\end{proof}

Coverage is used for \rulename{DI} only and equation induction for
\rulename{QI} only.

\begin{proposition}[Completeness of the native reading]
\label{base:prop:complete}
If $\Gamma\proves_{\Teq}\varphi$, then $\Gamma^{*}\proves_{\Nat}\varphi^{*}$.
\end{proposition}

\begin{proof}
By induction on the derivation.  Every logical rule of
\cref{base:def:logic} for a connective translates into the rule of the same
name in \cref{base:thm:logic}; the translation commutes with the
substitutions in the quantifier rules.  Reflexivity is
\cref{base:lem:eq}(i), and the Leibniz rule is \cref{base:thm:leibniz}.  It
remains to derive the native reading of each axiom in the empty context; by
\cref{base:lem:struct} all substitution instances then follow in every
context.

(T1) The readings are the conclusions of \rulename{I0} and \rulename{IS}.

(T2) Apply \cref{base:cor:cases} to the formula
$C(y)\defeq(y=0)\vee_*\exists_* v\,(y=\Suc v)$.  $C(0)$ follows from $0=0$
by $\vee_*$I.  $C(\Suc z)$ follows from $\Suc z=\Suc z$ by $\exists_*$I with
the witness $z$ and $\vee_*$I.  Hence $\proves C(x)$.

(T3) The reading is $(0=\Suc 0)\Rightarrow_*(0=\Suc 0)$, which is
$\Rightarrow_*$I applied to \rulename{Hyp}.

(T4) Let $U\defeq\forall_* z\,(E(z)\Rightarrow_* E(\Suc z))$.  By
$\Rightarrow_*$I twice it suffices to derive
$E(0),U\proves E(0)\wedge_*\Q{z}E(z)$, the conclusion being
$\forall_* z\,E(z)$ by definition.  The first component is a hypothesis.
For the second apply \rulename{QI}: over $E(0),U,E(z)$ with $z$ fresh,
$\forall_*$E gives $E(z)\Rightarrow_* E(\Suc z)$ and $\Rightarrow_*$E gives
$E(\Suc z)$.
\end{proof}

\begin{proposition}[The two formula round trips]
\label{base:prop:roundtrip}
\begin{enumerate}[label=\textup{(\roman*)}]
\item For every native formula $A$: $(A^\circ)^{*}\dashv\vdash A$ in $\Nat$.
\item For every ordinary formula $\varphi$:
 $\Teq\proves\varphi\leftrightarrow(\varphi^{*})^\circ$.
\end{enumerate}
\end{proposition}

\begin{proof}
(i) By induction on $A$.  An equation is its own round trip.

Let $A$ be $\Q{y}B$ and put $C\defeq(B^\circ)^{*}$, so that
$C\dashv\vdash B$ by the induction hypothesis and
$(A^\circ)^{*}=C(0)\Rightarrow_*\forall_* y\,C$.  From this formula derive
$\Q{y}B$ by \cref{base:lem:qI}: under the assumption $B(0)$ one has $C(0)$,
hence $\forall_* y\,C$, hence $C(\Suc z)$ by $\forall_*$E, hence
$B(\Suc z)$.  Conversely, from $\Q{y}B$ and the assumption $C(0)$ one has
$B(0)$, hence $B(y)$ for a fresh $y$ by \rulename{QE}, hence $C(y)$;
$\forall_*$I and $\Rightarrow_*$I give $(A^\circ)^{*}$.

Let $A$ be $\D{x}{y}B$.  Then $(A^\circ)^{*}=\exists_* x\,Y$, where
$Y\defeq((\Q{y}B)^\circ)^{*}$ and $Y\dashv\vdash\Q{y}B$ by the previous
case, with $x$ as a parameter.  From $\exists_* x\,Y$: $\exists_*$E at a
fresh $b$ gives $Y(b)$, hence $\Q{y}B(b,y)$, and \cref{base:lem:dI} gives
$\D{x}{y}B$.  From $\D{x}{y}B$: \rulename{DE} at a fresh $b$ gives
$\Q{y}B(b,y)$, hence $Y(b)$, and $\exists_*$I gives $\exists_* x\,Y$.

(ii) A derivation in $\Nat$ can be used inside $\Teq$: if
$\Delta\proves_{\Nat}A$ and $\Gamma\proves_{\Teq}B^\circ$ for every
$B\in\Delta$, then $\Gamma\proves_{\Teq}A^\circ$, by
\cref{base:prop:sound} and cut.  This is called \emph{transfer} below.  The
following equivalences are provable in $\Teq$ for all native formulas $X$,
$Y$.
\begin{enumerate}[label=\textup{(\alph*)}]
\item $(\bot_*)^\circ\leftrightarrow\bot$: by (T3) and $\bot$E.
\item $(X\wedge_* Y)^\circ\leftrightarrow X^\circ\wedge Y^\circ$: by transfer
 of $\wedge_*$E and of $\wedge_*$I.
\item $(\exists_* x\,X)^\circ\leftrightarrow\exists x\,X^\circ$: the left
 side is $\exists x\,(\Q{y}\operatorname{gate}(y,X))^\circ$, and
 $(\Q{y}\operatorname{gate}(y,X))^\circ\leftrightarrow X^\circ$ by transfer
 of \cref{base:lem:gate}(iii).
\item $(X\Rightarrow_* Y)^\circ\leftrightarrow(X^\circ\to Y^\circ)$: from
 left to right by transfer of $\Rightarrow_*$E.  From right to left, the
 goal $(X\Rightarrow_* Y)^\circ$ is
 $\sel(0,X,Y)^\circ\to\forall y\,\sel(y,X,Y)^\circ$.  Assume
 $\sel(0,X,Y)^\circ$.  Transfer of \cref{base:lem:gate}(iv) gives
 $X^\circ$, hence $Y^\circ$.  Transfer of $Y\proves X\Rightarrow_* Y$ gives
 $(X\Rightarrow_* Y)^\circ$, and its antecedent is the assumption.
\item $(X\vee_* Y)^\circ\leftrightarrow X^\circ\vee Y^\circ$: by (c) the left
 side is equivalent to $\exists x\,\sel(x,X,Y)^\circ$.  It suffices to
 derive $\forall x\,(\sel(x,X,Y)^\circ\to X^\circ\vee Y^\circ)$, and by
 \cref{base:lem:teqcases}(i) its two instances at $0$ and at $\Suc z$
 suffice.  They are transfers of $\sel(0,X,Y)\proves X$ and
 $\sel(\Suc z,X,Y)\proves Y$, which hold by \cref{base:lem:gate}(iv), (v),
 followed by $\vee$I.  No coverage of the witness $x$ and no transport along
 $x=0$ or $x=\Suc z$ is needed; this is the route of the Lean module
 \Lean{NativeRoundTrip}.  The converse is transfer of $\vee_*$I.
\item $(\forall_* x\,X)^\circ\leftrightarrow\forall x\,X^\circ$: transfer of
 $\wedge_*$E gives $X^\circ(0)$ and
 $(\Q{x}X)^\circ=X^\circ(0)\to\forall x\,X^\circ$ from the left side, hence
 the right side.  Conversely $\forall x\,X^\circ$ gives $X^\circ(0)$ and
 $(\Q{x}X)^\circ$, and transfer of $\wedge_*$I gives the left side.
\end{enumerate}
The claim follows by induction on $\varphi$, because provable equivalence is
a congruence for the connectives and quantifiers of the logic.
\end{proof}

\begin{corollary}[$\Nat$ and $\Teq$ prove the same]
\label{base:cor:nteq}
$\Gamma\proves_{\Nat}A$ if and only if $\Gamma^\circ\proves_{\Teq}A^\circ$,
and $\Gamma\proves_{\Teq}\varphi$ if and only if
$\Gamma^{*}\proves_{\Nat}\varphi^{*}$.
\end{corollary}

\begin{proof}
If $\Gamma^\circ\proves_{\Teq}A^\circ$, then
$(\Gamma^\circ)^{*}\proves_{\Nat}(A^\circ)^{*}$ by
\cref{base:prop:complete}, and \cref{base:prop:roundtrip}(i) with cut
gives $\Gamma\proves_{\Nat}A$.  If $\Gamma^{*}\proves_{\Nat}\varphi^{*}$,
then $(\Gamma^{*})^\circ\proves_{\Teq}(\varphi^{*})^\circ$ by
\cref{base:prop:sound}, and \cref{base:prop:roundtrip}(ii) gives
$\Gamma\proves_{\Teq}\varphi$.  The other two implications are the two
propositions.
\end{proof}

In Lean these results are \Lean{NPrf.sound}, \Lean{NPrf.ofPrf},
\Lean{NEquiv.ofFm\_toFm}, \Lean{Teq.iff\_toFm\_ofFm}, \Lean{NPrf.iff\_teq}
and \Lean{Teq.iff\_native}.

\subsection{Term interpretations}

\begin{definition}[Term interpretation]
\label{base:def:interp}
Let $T$ and $T'$ be theories in intuitionistic logic with equality over term
languages $\mathcal L$ and $\mathcal L'$.  A \emph{term interpretation} of
$T$ in $T'$ is a map $F$ from the terms of $\mathcal L$ to the terms of
$\mathcal L'$ such that
\begin{enumerate}[label=\textup{(\roman*)}]
\item $F(x)=x$ for every variable $x$, and $\FV(F(t))\subseteq\FV(t)$;
\item $T'\proves F(t[\theta])=F(t)[F\circ\theta]$ for every term $t$ and
 substitution $\theta$;
\item $T'\proves\varphi^{F}$ for every axiom $\varphi$ of $T$,
\end{enumerate}
where $\varphi^{F}$ is $\varphi$ with each equation $s=t$ replaced by
$F(s)=F(t)$.
\end{definition}

A term interpretation does not relativize quantifiers and reads equality as
equality: its domain is the whole domain.

\begin{lemma}[Term interpretations preserve derivability]
\label{base:lem:interp-sound}
Let $F$ be a term interpretation of $T$ in $T'$.
\begin{enumerate}[label=\textup{(\roman*)}]
\item $T'\proves(\varphi[\theta])^{F}\leftrightarrow\varphi^{F}[F\circ\theta]$.
\item If $\Gamma\proves_{T}\varphi$, then $\Gamma^{F}\proves_{T'}\varphi^{F}$.
\end{enumerate}
\end{lemma}

\begin{proof}
(i) By induction on $\varphi$.  For an equation the two sides are
$F(s[\theta])=F(t[\theta])$ and $F(s)[F\circ\theta]=F(t)[F\circ\theta]$,
which are equivalent by condition (ii) of the definition.  The connectives
and quantifiers are handled by the congruence of provable equivalence, the
bound variables being fixed by $\theta$ and by $F$.
(ii) By induction on the derivation.  The propositional rules, $\forall$I
and $\exists$E translate literally.  For $\forall$E the premise gives
$\forall x\,\varphi^{F}$, hence $\varphi^{F}\subst{F(t)}{x}$, which is
equivalent to $(\varphi\subst{t}{x})^{F}$ by (i); $\exists$I is similar.
Reflexivity is reflexivity.  For the Leibniz rule the premises give
$F(s_i)=F(t_i)$ and, by (i), $\varphi^{F}\subst{F(\vec s)}{\vec x}$; the
Leibniz rule of $T'$ and (i) give $(\varphi\subst{\vec t}{\vec x})^{F}$.  An
instance $\varphi[\theta]$ of an axiom translates to a formula equivalent to
$\varphi^{F}[F\circ\theta]$, an instance of a theorem of $T'$.
\end{proof}

\begin{lemma}[Round trips give faithfulness]
\label{base:lem:faithful}
Let $F$ be a term interpretation of $T$ in $T'$ and $G$ one of $T'$ in $T$,
with $T\proves G(F(t))=t$ and $T'\proves F(G(t'))=t'$ for all terms.  Then
\begin{enumerate}[label=\textup{(\roman*)}]
\item $T\proves\varphi\leftrightarrow(\varphi^{F})^{G}$ and
 $T'\proves\psi\leftrightarrow(\psi^{G})^{F}$;
\item $\Gamma\proves_{T}\varphi$ if and only if
 $\Gamma^{F}\proves_{T'}\varphi^{F}$, and $\Gamma\proves_{T'}\psi$ if and
 only if $\Gamma^{G}\proves_{T}\psi^{G}$.
\end{enumerate}
\end{lemma}

\begin{proof}
(i) By induction on the formula; for an equation $s=t$ the round trip is
$G(F(s))=G(F(t))$, equivalent to $s=t$ by the hypothesis.  (ii) One
direction is \cref{base:lem:interp-sound}.  For the other, if
$\Gamma^{F}\proves_{T'}\varphi^{F}$ then
$(\Gamma^{F})^{G}\proves_{T}(\varphi^{F})^{G}$, and (i) gives
$\Gamma\proves_{T}\varphi$.
\end{proof}

\subsection{From \texorpdfstring{$\Teq$}{Teq} to
 \texorpdfstring{$\iPRA$}{iPRA}}

A term of $\iPRA$ has no binder.  An iterator term is therefore compiled
into a function symbol of its free variables, and the iterator becomes a
primitive recursion whose step ignores the counter.  For a symbol $A$ of
arity $k$ and a symbol $H$ of arity $k+1$ put
\[
 \iota(A,H)\defeq\mathrm{prec}\bigl(A,\ \mathrm{comp}(H;\pi^{k+2}_2,
   \pi^{k+2}_3,\dots,\pi^{k+2}_{k+2})\bigr),
\]
a symbol of arity $k+1$.  By the axioms for recursion, composition and
projection,
\begin{equation}
 \iPRA\proves\iota(A,H)(0,\vec x)=A(\vec x),\qquad
 \iPRA\proves\iota(A,H)(\Suc y,\vec x)=H\bigl(\iota(A,H)(y,\vec x),\vec x\bigr).
 \label{base:eq:iota}
\end{equation}

\begin{definition}[The translation $\tau$]
\label{base:def:tau}
Fix an enumeration of the variables.  For a list
$\vec x=x_1,\dots,x_k$ of distinct variables and a term $t$ with
$\FV(t)\subseteq\vec x$, the symbol $[t]_{\vec x}$ of arity $k$ is
\begin{gather*}
 {}[x_i]_{\vec x}\defeq\pi^k_i,\qquad
 {}[0]_{\vec x}\defeq\mathrm{comp}(\mathrm{zero};\,),\qquad
 [\Suc t]_{\vec x}\defeq\mathrm{comp}(\mathrm{succ};[t]_{\vec x}),\\
 {}[\It{c}{a}{u}{h}]_{\vec x}\defeq
  \mathrm{comp}\bigl(\iota([a]_{\vec x},[h]_{u\vec x});\
   [c]_{\vec x},\pi^k_1,\dots,\pi^k_k\bigr)\qquad(u\notin\vec x).
\end{gather*}
The translation of terms is
\[
 \tau(x)\defeq x,\qquad\tau(0)\defeq 0,\qquad
 \tau(\Suc t)\defeq\Suc\tau(t),\qquad
 \tau(\It{c}{a}{u}{h})\defeq
  \iota\bigl([a]_{\vec y},[h]_{u\vec y}\bigr)\bigl(\tau(c),\vec y\bigr),
\]
where $\vec y$ lists the variables of $\FV(a)\cup(\FV(h)\setminus\{u\})$ in
the order of the enumeration.
\end{definition}

The translation does not commute with substitution literally: a substitution
into the body of an iterator changes the symbol.  It commutes provably.

\begin{lemma}[Extensionality of the simulated iterator]
\label{base:lem:iota-ext}
Let $\vec z$ be a list of distinct variables, $\vec s$ a list of terms, and
$u$ a variable that is not in $\vec z$ and not free in $\vec s$.  If
$\iPRA\proves A'(\vec z)=A(\vec s)$ and
$\iPRA\proves H'(u,\vec z)=H(u,\vec s)$, then
$\iPRA\proves\iota(A',H')(t,\vec z)=\iota(A,H)(t,\vec s)$ for every term
$t$.
\end{lemma}

\begin{proof}
By quantifier-free induction on a fresh variable $n$ in the equation
$\iota(A',H')(n,\vec z)=\iota(A,H)(n,\vec s)$.  The base is
$A'(\vec z)=A(\vec s)$ by \eqref{base:eq:iota}.  In the step,
\begin{align*}
 \iota(A',H')(\Suc n,\vec z)
 &=H'\bigl(\iota(A',H')(n,\vec z),\vec z\bigr)
  =H\bigl(\iota(A',H')(n,\vec z),\vec s\bigr)\\
 &=H\bigl(\iota(A,H)(n,\vec s),\vec s\bigr)=\iota(A,H)(\Suc n,\vec s),
\end{align*}
by \eqref{base:eq:iota}, the second hypothesis instantiated at
$u=\iota(A',H')(n,\vec z)$, and the induction equation.  Instantiate $n$ at
$t$.
\end{proof}

\begin{lemma}[Evaluation and substitution]
\label{base:lem:tau}
\begin{enumerate}[label=\textup{(\roman*)}]
\item If $\FV(t)\subseteq\vec x$, then
 $\iPRA\proves[t]_{\vec x}(\vec r)=\tau(t)\subst{\vec r}{\vec x}$ for all
 terms $\vec r$.  In particular $\iPRA\proves[t]_{\vec x}(\vec x)=\tau(t)$.
\item $\iPRA\proves\tau(t[\theta])=\tau(t)[\tau\circ\theta]$ for every
 substitution $\theta$.
\end{enumerate}
\end{lemma}

\begin{proof}
(i) By induction on $t$, for all lists $\vec x$.  For a variable, $0$ and a
successor the claim follows from the axioms for projection and composition
and the induction hypothesis.  Let $t$ be $\It{c}{a}{u}{h}$, and let
$\vec y$ be the list of \cref{base:def:tau}, a sublist of $\vec x$.  By the
induction hypothesis for $a$ and $h$ with variables as arguments,
$[a]_{\vec x}(\vec x)=\tau(a)=[a]_{\vec y}(\vec y)$ and
$[h]_{u\vec x}(u,\vec x)=\tau(h)=[h]_{u\vec y}(u,\vec y)$.  By
\cref{base:lem:iota-ext},
\[
 \iota\bigl([a]_{\vec x},[h]_{u\vec x}\bigr)(n,\vec x)
 =\iota\bigl([a]_{\vec y},[h]_{u\vec y}\bigr)(n,\vec y)
\]
for a fresh variable $n$.  By the axioms for composition and projection and
the induction hypothesis for $c$,
$[t]_{\vec x}(\vec r)
 =\iota([a]_{\vec x},[h]_{u\vec x})(\tau(c)\subst{\vec r}{\vec x},\vec r)$.
Substituting $\vec r$ for $\vec x$ and $\tau(c)\subst{\vec r}{\vec x}$ for
$n$ in the displayed equation identifies this term with
$\tau(t)\subst{\vec r}{\vec x}$.

(ii) By induction on $t$.  The cases other than the iterator are immediate.
Let $t$ be $\It{c}{a}{u}{h}$ with $u$ fixed by $\theta$ and not free in its
values, so that $t[\theta]=\It{c[\theta]}{a[\theta]}{u}{h[\theta]}$.  Let
$\vec y$ be the list for $a$, $h$ and $\vec z$ the list for $a[\theta]$,
$h[\theta]$.  By (i) and the induction hypothesis,
\begin{align*}
 {}[a[\theta]]_{\vec z}(\vec z)&=\tau(a[\theta])=\tau(a)[\tau\circ\theta]
   =[a]_{\vec y}\bigl(\tau(\theta(\vec y))\bigr),\\
 {}[h[\theta]]_{u\vec z}(u,\vec z)&=\tau(h[\theta])=\tau(h)[\tau\circ\theta]
   =[h]_{u\vec y}\bigl(u,\tau(\theta(\vec y))\bigr).
\end{align*}
\Cref{base:lem:iota-ext} with $\vec s\defeq\tau(\theta(\vec y))$ and the
induction hypothesis for $c$ give
\begin{align*}
 \tau(t[\theta])
 &=\iota\bigl([a[\theta]]_{\vec z},[h[\theta]]_{u\vec z}\bigr)
   \bigl(\tau(c[\theta]),\vec z\bigr)\\
 &=\iota\bigl([a]_{\vec y},[h]_{u\vec y}\bigr)
   \bigl(\tau(c)[\tau\circ\theta],\tau(\theta(\vec y))\bigr)
  =\tau(t)[\tau\circ\theta].
\end{align*}
\end{proof}

\begin{proposition}[$\Teq$ in $\iPRA$]
\label{base:prop:tau}
$\tau$ is a term interpretation of $\Teq$ in $\iPRA$.
\end{proposition}

\begin{proof}
Conditions (i) and (ii) of \cref{base:def:interp} are the definition and
\cref{base:lem:tau}(ii).  The axioms:

(T1) With the notation of \cref{base:def:tau},
$\tau(\It{0}{a}{u}{h})=\iota(\dots)(0,\vec y)=[a]_{\vec y}(\vec y)=\tau(a)$
by \eqref{base:eq:iota} and \cref{base:lem:tau}(i).  Writing
$q\defeq\tau(\It{n}{a}{u}{h})$,
\[
 \tau(\It{\Suc n}{a}{u}{h})=\iota(\dots)(\Suc\tau(n),\vec y)
 =[h]_{u\vec y}(q,\vec y)=\tau(h)\subst{q}{u}
 =\tau\bigl(h\subst{\It{n}{a}{u}{h}}{u}\bigr),
\]
by \eqref{base:eq:iota} and both parts of \cref{base:lem:tau}.

(T2) Let $\mathrm{pd}\defeq\mathrm{prec}(\mathrm{zero},\pi^2_1)$, so that
$\mathrm{pd}(0)=0$ and $\mathrm{pd}(\Suc y)=y$.  Quantifier-free induction
on $x$ proves $x=0\vee x=\Suc\mathrm{pd}(x)$: the base is the left disjunct,
and the step is the right disjunct, since
$\Suc x=\Suc\mathrm{pd}(\Suc x)$.  Existential introduction gives the image
of (T2).

(T3) is the axiom $\Suc x=0\to\bot$ at $x=0$, with symmetry.

(T4) For an equation $E(z)$, written $s=t$, the image of the axiom is
\[
 \tau(s\subst{0}{z})=\tau(t\subst{0}{z})\to\Bigl(\forall z\,\bigl(\tau(s)=\tau(t)\to
 \tau(s\subst{\Suc z}{z})=\tau(t\subst{\Suc z}{z})\bigr)\to
 \forall z\,\tau(s)=\tau(t)\Bigr),
\]
which by \cref{base:lem:tau}(ii) is equivalent to the induction axiom of
$\iPRA$ for the equation $\tau(s)=\tau(t)$.
\end{proof}

\subsection{From \texorpdfstring{$\iPRA$}{iPRA} to
 \texorpdfstring{$\Teq$}{Teq}}

\begin{proposition}[$\iPRA$ in $\Teq$]
\label{base:prop:sigma}
The translation $\sigma$ of \cref{base:sec:arith} is a term interpretation
of $\iPRA$ in $\Teq$.
\end{proposition}

\begin{proof}
$\sigma$ fixes variables and commutes with substitution literally.  The
defining equations hold by \cref{base:thm:primrec}(ii).  The images of
$\Suc x=0\to\bot$ and $\Suc x=\Suc y\to x=y$ are those formulas themselves,
which hold by \cref{base:lem:succ}.  For a quantifier-free $\varphi$ the
formula $\varphi^{\sigma}$ is quantifier free and
$(\varphi\subst{t}{z})^{\sigma}=\varphi^{\sigma}\subst{\sigma(t)}{z}$, so the
image of the induction axiom for $\varphi$ is the instance of
\cref{base:thm:qfi} for $\varphi^{\sigma}$.
\end{proof}

\subsection{The two round trips on terms}

\begin{lemma}[A recursion that ignores its counter is the iterator]
\label{base:lem:prc-iter}
If $y$ is not free in $h(v)$, then
$\Teq\proves R(c,a;\,y\,v.\,h(v))=\It{c}{a}{u}{h(u)}$.
\end{lemma}

\begin{proof}
\Cref{base:lem:uniq} for $f(z)\defeq R(z,a;\,y\,v.\,h(v))$: by
\cref{base:thm:primrec}(i), $f(0)=a$ and $f(\Suc z)=h(f(z))$.
\end{proof}

\begin{proposition}[Round trips on terms]
\label{base:prop:termrt}
\begin{enumerate}[label=\textup{(\roman*)}]
\item $\Teq\proves\sigma(\tau(t))=t$ for every term $t$ of $\Teq$.
\item $\iPRA\proves\tau(\sigma(t))=t$ for every term $t$ of $\iPRA$.
\end{enumerate}
\end{proposition}

\begin{proof}
(i) By the definition of $\ulcorner\cdot\urcorner$ the term
$\ulcorner\mathrm{comp}(H;\pi^{k+2}_2,\dots,\pi^{k+2}_{k+2})\urcorner(y,v,\vec x)$
is $\ulcorner H\urcorner(v,\vec x)$, so by \cref{base:lem:prc-iter}
\begin{equation}
 \Teq\proves\ulcorner\iota(A,H)\urcorner(c,\vec x)
 =\It{c}{\ulcorner A\urcorner(\vec x)}{u}{\ulcorner H\urcorner(u,\vec x)}.
 \label{base:eq:code-iota}
\end{equation}
First, $\Teq\proves\ulcorner[t]_{\vec x}\urcorner(\vec x)=t$, by induction
on $t$.  For a variable, $0$ and a successor both sides are the same term
up to the induction hypothesis.  For an iterator,
\begin{align*}
 \ulcorner[\It{c}{a}{u}{h}]_{\vec x}\urcorner(\vec x)
 &=\ulcorner\iota([a]_{\vec x},[h]_{u\vec x})\urcorner
   \bigl(\ulcorner[c]_{\vec x}\urcorner(\vec x),\vec x\bigr)\\
 &=\It{\ulcorner[c]_{\vec x}\urcorner(\vec x)}
     {\ulcorner[a]_{\vec x}\urcorner(\vec x)}{u}
     {\ulcorner[h]_{u\vec x}\urcorner(u,\vec x)},
\end{align*}
by \eqref{base:eq:code-iota}, and the last term equals $\It{c}{a}{u}{h}$ by
the induction hypothesis, congruence in the counter and the seed, and
\cref{base:cor:rot}(ii) for the body.  Now the claim follows by induction on
$t$; for an iterator, with the notation of \cref{base:def:tau},
\begin{align*}
 \sigma(\tau(\It{c}{a}{u}{h}))
 &=\ulcorner\iota([a]_{\vec y},[h]_{u\vec y})\urcorner
   \bigl(\sigma(\tau(c)),\vec y\bigr)\\
 &=\It{\sigma(\tau(c))}{\ulcorner[a]_{\vec y}\urcorner(\vec y)}{u}
     {\ulcorner[h]_{u\vec y}\urcorner(u,\vec y)}
  =\It{c}{a}{u}{h}.
\end{align*}

(ii) First, $\iPRA\proves\tau(\ulcorner f\urcorner)=f(x_1,\dots,x_k)$ for
every symbol $f$ of arity $k$, by induction on $f$.  For $\mathrm{zero}$ and
$\mathrm{succ}$ the two sides are the same term, and for a projection the
claim is its axiom.  For a composition, by \cref{base:lem:tau}(ii) and the
induction hypotheses,
\[
 \tau\bigl(\ulcorner f\urcorner(\ulcorner g_1\urcorner,\dots,
   \ulcorner g_m\urcorner)\bigr)
 =\tau(\ulcorner f\urcorner)\subst{\tau(\ulcorner\vec g\urcorner)}{\vec x}
 =f\bigl(g_1(\vec x),\dots,g_m(\vec x)\bigr)
 =\mathrm{comp}(f;\vec g)(\vec x).
\]
For a recursion $p\defeq\mathrm{prec}(f,g)$ write
$q\defeq\tau(\ulcorner p\urcorner)$ and $\vec x'\defeq x_2,\dots,x_{k+1}$.
By \cref{base:thm:primrec}(i), $\Teq$ proves
$\ulcorner p\urcorner\subst{0}{x_1}=\ulcorner f\urcorner(\vec x')$ and
$\ulcorner p\urcorner\subst{\Suc x_1}{x_1}
 =\ulcorner g\urcorner(x_1,\ulcorner p\urcorner,\vec x')$.  By
\cref{base:prop:tau}, \cref{base:lem:interp-sound} and
\cref{base:lem:tau}(ii), $\iPRA$ proves the images of these equations, and
with the induction hypotheses for $f$ and $g$ they read
\[
 q\subst{0}{x_1}=f(\vec x'),\qquad
 q\subst{\Suc x_1}{x_1}=g(x_1,q,\vec x').
\]
Quantifier-free induction on $x_1$ in the equation
$q=p(x_1,\vec x')$ now succeeds: the base is
$q\subst{0}{x_1}=f(\vec x')=p(0,\vec x')$, and under the induction equation
$q\subst{\Suc x_1}{x_1}=g(x_1,q,\vec x')=g(x_1,p(x_1,\vec x'),\vec x')
 =p(\Suc x_1,\vec x')$.

The claim follows by induction on $t$:
$\tau(\sigma(f(\vec t)))=\tau(\ulcorner f\urcorner)
 \subst{\tau(\sigma(\vec t))}{\vec x}=f(\tau(\sigma(\vec t)))=f(\vec t)$,
by \cref{base:lem:tau}(ii), the first step and the induction hypothesis.
\end{proof}

\begin{remark}[The route of the Lean modules]
\label{base:rem:tau-lean}
The modules \Lean{ScalarToPRA} and \Lean{RoundTripScalar} compile terms of
$\iPRA$ rather than iterator terms.  For a term $p$ of $\iPRA$ the symbol
$\langle p\rangle_{\vec x}$ is $\pi^k_i$ at the variable $x_i$ and
$\mathrm{comp}(f;\langle p_1\rangle_{\vec x},\dots,\langle p_m\rangle_{\vec x})$
at $f(p_1,\dots,p_m)$, and the iterator case of $\tau$ reads
$\iota\bigl(\langle\tau(a)\rangle_{\vec y},\langle\tau(h)\rangle_{u\vec y}\bigr)
(\tau(c),\vec y)$.  Then $\langle p\rangle_{\vec x}(\vec s)=p\subst{\vec s}{\vec x}$
is a fact about $\iPRA$ alone, proved by an induction with two cases, so
\cref{base:lem:tau}(i) needs no induction on iterator terms; the hypotheses of
\cref{base:lem:iota-ext} are equations between terms; and
$\sigma(\langle p\rangle_{\vec x})$ coincides syntactically with $\sigma(p)$, so
the first step of the proof of \cref{base:prop:termrt}(i) disappears and
$\sigma(\tau(t))=t$ is one induction on $t$.  The recursion equations of the
code of $\mathrm{prec}(f,g)$ are proved once, from those of the recursor of
\cref{base:thm:primrec}, and the translations of the two defining axioms of
$\mathrm{prec}(f,g)$ are derived from them.
\end{remark}

\subsection{The bi-interpretation theorem}

\begin{theorem}[Bi-interpretation]
\label{base:thm:bridge}
The systems $\Nat$, $\Teq$ and $\iPRA$ are pairwise bi-interpretable.
Precisely:
\begin{enumerate}[label=\textup{(\roman*)}]
\item \emph{$\Nat$ and $\Teq$.}  $\Gamma\proves_{\Nat}A$ if and only if
 $\Gamma^\circ\proves_{\Teq}A^\circ$; $\Gamma\proves_{\Teq}\varphi$ if and
 only if $\Gamma^{*}\proves_{\Nat}\varphi^{*}$;
 $(A^\circ)^{*}\dashv\vdash A$ in $\Nat$; and
 $\Teq\proves\varphi\leftrightarrow(\varphi^{*})^\circ$.
\item \emph{$\Teq$ and $\iPRA$.}  $\tau$ and $\sigma$ are term
 interpretations in opposite directions;
 $\Teq\proves\sigma(\tau(t))=t$ and $\iPRA\proves\tau(\sigma(t))=t$;
 $\Gamma\proves_{\Teq}\varphi$ if and only if
 $\Gamma^{\tau}\proves_{\iPRA}\varphi^{\tau}$;
 $\Gamma\proves_{\iPRA}\psi$ if and only if
 $\Gamma^{\sigma}\proves_{\Teq}\psi^{\sigma}$;
 $\Teq\proves\varphi\leftrightarrow(\varphi^{\tau})^{\sigma}$; and
 $\iPRA\proves\psi\leftrightarrow(\psi^{\sigma})^{\tau}$.
\item \emph{$\Nat$ and $\iPRA$.}  With $A^{\flat}\defeq(A^\circ)^{\tau}$
 and $\psi^{\sharp}\defeq(\psi^{\sigma})^{*}$:
 $\Gamma\proves_{\Nat}A$ if and only if
 $\Gamma^{\flat}\proves_{\iPRA}A^{\flat}$;
 $\Gamma\proves_{\iPRA}\psi$ if and only if
 $\Gamma^{\sharp}\proves_{\Nat}\psi^{\sharp}$;
 $(A^{\flat})^{\sharp}\dashv\vdash A$ in $\Nat$; and
 $\iPRA\proves\psi\leftrightarrow(\psi^{\sharp})^{\flat}$.
\item \emph{Shape of the interpretations.}  Every translation fixes the
 variables, leaves the range of every quantifier and binder the whole
 numerical domain, and reads equality as equality.  On terms the composites
 are provably the identity, so the definable map between the domain and its
 double interpretation is the identity, with graph $x=y$; it is total,
 bijective and commutes with every term former and function symbol.
\item \emph{Equiconsistency.}  $\proves_{\Nat}0=\Suc 0$ if and only if
 $\proves_{\Teq}\bot$ if and only if $\proves_{\iPRA}\bot$.
\end{enumerate}
\end{theorem}

\begin{proof}
(i) is \cref{base:cor:nteq,base:prop:roundtrip}.  (ii) is
\cref{base:prop:tau,base:prop:sigma,base:prop:termrt} and
\cref{base:lem:faithful}.  (iii) The two equivalences are the compositions
of those in (i) and (ii).  For the first round trip,
$(A^{\flat})^{\sharp}=(((A^\circ)^{\tau})^{\sigma})^{*}$.  By (ii),
$\Teq$ proves $A^\circ\leftrightarrow((A^\circ)^{\tau})^{\sigma}$, so by
\cref{base:prop:complete} the native readings of the two sides are
interderivable in $\Nat$, and $(A^\circ)^{*}\dashv\vdash A$ by (i).  For the
second, $(\psi^{\sharp})^{\flat}=(((\psi^{\sigma})^{*})^\circ)^{\tau}$.  By
(i), $\Teq$ proves
$\psi^{\sigma}\leftrightarrow((\psi^{\sigma})^{*})^\circ$; applying $\tau$,
$\iPRA$ proves
$(\psi^{\sigma})^{\tau}\leftrightarrow(\psi^{\sharp})^{\flat}$, and
$\iPRA\proves\psi\leftrightarrow(\psi^{\sigma})^{\tau}$ by (ii).
(iv) The first sentence holds by the definitions of $A^\circ$,
$\varphi^{*}$, $\tau$ and $\sigma$; the second is
\cref{base:prop:termrt}, where for the span between $\Nat$ and $\Teq$ the
terms are not translated at all.
(v) If $\proves_{\Nat}0=\Suc 0$ then $\proves_{\Teq}0=\Suc 0$ by
\cref{base:prop:sound}, and (T3) gives $\bot$.  If $\proves_{\Teq}\bot$
then $\proves_{\Nat}\bot^{*}$, which is $0=\Suc 0$.  Since
$\bot^{\tau}=\bot=\bot^{\sigma}$, (ii) gives the second equivalence.
\end{proof}

The translations act on finite derivations effectively.  No claim is made
that any of the systems proves its own consistency.  The interpretations use
no restriction to standard numerals, no retract or quotient of the domain,
and no assumed symbol for equality: the reflection of equality is
\cref{base:thm:reflect}, a theorem.

In Lean, part (ii) is \Lean{biInterp} with \Lean{teq\_iff\_pra},
\Lean{pra\_iff\_teq}, \Lean{teq\_roundTrip}, \Lean{pra\_roundTrip} and
\Lean{teq\_consistent\_iff}; part (iii) is \Lean{native\_iff\_pra},
\Lean{pra\_iff\_native}, \Lean{native\_roundTrip} and
\Lean{pra\_native\_roundTrip}; and part (v) is
\Lean{native\_inconsistent\_iff}.  The final theorems depend on no axiom
beyond \Lean{propext}, \Lean{Classical.choice} and \Lean{Quot.sound}, and
part (ii) on \Lean{propext} and \Lean{Quot.sound} alone.

\begin{remark}[Consistency]
\label{base:rem:consistent}
Read over the natural numbers, with the iterator as iteration, every rule
of $\Nat$ preserves truth, and $0=\Suc 0$ is false.  Hence $\Nat$ does not
derive $0=\Suc 0$, and by \cref{base:thm:bridge}(v) neither $\Teq$ nor
$\iPRA$ derives $\bot$.  This argument is carried out in the metatheory and
is formalized in Lean as \Lean{NPrf.sound\_std}, \Lean{native\_consistent},
\Lean{teq\_consistent} and \Lean{pra\_consistent}.
\end{remark}

\begin{remark}[How the packages state the bridge]
\label{base:rem:bridgepackages}
Every package repeats this theorem for its own calculus, by composing an
equivalence of its calculus with the inherited one and the inherited bridge.
The reference statement composes two steps: the native calculus against an
ordinary theory $T_I$, which is $\Teq$ with the iterators presented as staged
function symbols (\cref{base:rem:nolang}), and $T_I$ against $\iPRA$.  Its
proof of the second step identifies the double translation of an operation
with the operation by uniqueness of iteration and of primitive recursion,
as in \cref{base:prop:termrt}; the explicit treatment of substitution
(\cref{base:lem:iota-ext,base:lem:tau}) is spelled out in the Lean
formalization and is given here in full.
\end{remark}

The universal quantifier of the native reading may be replaced by the
single-binder universal of \cref{base:def:single-binder}, as
\pkg{single\_binder} proposes.  Let $\varphi^{\dagger}$ be defined as
$\varphi^{*}$ in \cref{base:def:star}, except that
$(\forall x\,\varphi)^{\dagger}\defeq\mathcal U_x\,\varphi^{\dagger}$.  By
the commutation of $\mathcal U$ with substitution,
$(\varphi[\theta])^{\dagger}=\varphi^{\dagger}[\theta]$.

\begin{corollary}[The bridge with the single-binder universal]
\label{base:cor:single-binder-bridge}
\begin{enumerate}[label=\textup{(\roman*)}]
\item $\varphi^{\dagger}\dashv\vdash\varphi^{*}$ in $\Nat$ for every
 ordinary formula $\varphi$.
\item $\Teq\proves(\mathcal U_xX)^\circ\leftrightarrow\forall x\,X^\circ$ for
 every native formula $X$.
\item \Cref{base:thm:bridge} holds with $\varphi^{\dagger}$ in place of
 $\varphi^{*}$ in part (i) and with $\psi^{\sharp}\defeq(\psi^{\sigma})^{\dagger}$
 in part (iii); parts (ii), (iv) and (v) are unchanged, except that in (iv)
 the translation $\dagger$ fixes the free variables and replaces a bound
 universal variable by the term $P(z)$ of its new bound variable.
\end{enumerate}
\end{corollary}

\begin{proof}
(i) By induction on $\varphi$.  Equations and $\bot$ have the same two
translations.  In the other cases the induction hypothesis gives
$\psi^{\dagger}\dashv\vdash\psi^{*}$ and $\chi^{\dagger}\dashv\vdash\chi^{*}$
for the immediate subformulas, and, by \cref{base:lem:struct}, for all their
substitution instances.  The rules cited are those of \cref{base:thm:logic}
and \cref{base:prop:single-binder}, and every step from one formula to an
interderivable one is a cut.  For $\psi\wedge\chi$: $\wedge_*$E gives
$\psi^{\dagger}$ and $\chi^{\dagger}$, hence $\psi^{*}$ and $\chi^{*}$, and
$\wedge_*$I gives $\psi^{*}\wedge_*\chi^{*}$.  For $\psi\vee\chi$: $\vee_*$E
with the branches $\psi^{\dagger}\proves\psi^{*}\proves\psi^{*}\vee_*\chi^{*}$
and $\chi^{\dagger}\proves\chi^{*}\proves\psi^{*}\vee_*\chi^{*}$, the last
steps by $\vee_*$I.  For $\psi\to\chi$: over the context
$\psi^{\dagger}\Rightarrow_*\chi^{\dagger},\ \psi^{*}$ one has
$\psi^{\dagger}$, hence $\chi^{\dagger}$ by $\Rightarrow_*$E, hence
$\chi^{*}$, and $\Rightarrow_*$I discharges $\psi^{*}$.  For
$\exists x\,\psi$: $\exists_*$E at a fresh $b$ with the premise
$\psi^{\dagger}(b)\proves\psi^{*}(b)\proves\exists_* x\,\psi^{*}$, the last
step by $\exists_*$I.  For $\forall x\,\psi$:
\cref{base:prop:single-binder}(ii) at the variable $x$ gives
$\mathcal U_x\psi^{\dagger}\proves\psi^{\dagger}(x)$, hence $\psi^{*}(x)$, and
$\forall_*$I gives $\forall_* x\,\psi^{*}$, because $x$ is not free in
$\mathcal U_x\psi^{\dagger}$; conversely $\forall_*$E at $x$, the induction
hypothesis and \cref{base:prop:single-binder}(i).  In each case the
converse is the same derivation with $\dagger$ and $*$ exchanged.

(ii) Transfer, in the sense of the proof of
\cref{base:prop:roundtrip}(ii), of the two sequents of
\cref{base:prop:single-binder}(iii) gives
$\Teq\proves(\mathcal U_xX)^\circ\leftrightarrow(\forall_* x\,X)^\circ$, and
\cref{base:prop:roundtrip}(ii)(f) gives
$(\forall_* x\,X)^\circ\leftrightarrow\forall x\,X^\circ$.

(iii) Part (i) of the theorem: $\Gamma\proves_{\Teq}\varphi$ if and only if
$\Gamma^{*}\proves_{\Nat}\varphi^{*}$, by \cref{base:cor:nteq}, if and only if
$\Gamma^{\dagger}\proves_{\Nat}\varphi^{\dagger}$, by (i) and
\cref{base:lem:cut} applied to each member of the context and to the
conclusion.  $(A^\circ)^{\dagger}\dashv\vdash(A^\circ)^{*}\dashv\vdash A$ by (i)
and \cref{base:prop:roundtrip}(i).  By \cref{base:prop:sound} applied to the
two sequents of (i), $\Teq\proves(\varphi^{\dagger})^\circ\leftrightarrow
(\varphi^{*})^\circ$, and \cref{base:prop:roundtrip}(ii) gives
$\Teq\proves\varphi\leftrightarrow(\varphi^{\dagger})^\circ$.  The proof of
part (iii) of the theorem uses only parts (i) and (ii) and therefore applies
verbatim.  Part (iv): the binders of $\varphi^{\dagger}$ range over the whole
domain under the ordinary reading, equality is read as equality, and the
term composites are those of \cref{base:prop:termrt}, which $\dagger$ does
not touch.  Part (v): $\bot^{\dagger}=\bot^{*}$ is $0=\Suc 0$.
\end{proof}

The translation $\dagger$ trades the exponential growth of the binder rank
and of the number of copies of a matrix under nested universals
(\cref{base:rem:single-binder-size}) for one occurrence of $P(z)$ per
occurrence of a quantified variable and the selectors of one gate per
universal.  The axioms and rules of $\Nat$, $\Teq$
and $\iPRA$ are unchanged.

% ---------------------------------------------------------------------
\section{Clause contexts}
\label{base:sec:clauses}

This section draws on the definition in the main report of
\pkg{raw\_clausal} and on the reference statement of the inherited basis
filed with \pkg{certified\_zero}.  Only the notion and its reading are fixed
here; the clause calculi are the subject of other chapters, beginning with
\cref{claus:sec:five}.

\begin{definition}[Clauses and clause contexts]
\label{base:def:clause}
A \emph{clause} is an expression
\[
 [A;\,y.\,B(y)],\qquad y\notin\FV(A),
\]
where $A$ and $B$ are native formulas.  The clause binds $y$ in its
consequent $B$ only.  It is a member of a context, not a formula: it cannot
occur inside a formula or as the conclusion of a sequent.  A \emph{clause
context} is a finite list of clause occurrences.  Its \emph{ordinary
reading} is
\[
 [A;\,y.\,B(y)]^\circ\defeq A^\circ\to\forall y\,B^\circ(y).
\]
An assumption of a formula $A$ is the \emph{constant clause}
\[
 [\,0=0;\,v.\,A\,],\qquad v\notin\FV(A),
\]
and $\Gamma,A$ abbreviates $\Gamma,[\,0=0;\,v.\,A\,]$.
\end{definition}

A clause is used by the inference that, for an occurrence
$[A;\,y.\,B(y)]$ in $\Gamma$, passes from $\Gamma\proves A$ to
$\Gamma\proves B(t)$ for a term $t$.  Three observations connect clauses to
the formulas of this chapter.
\begin{enumerate}[label=\textup{(\roman*)}]
\item The reading of the constant clause is
 $0=0\to\forall v\,A^\circ$, which is equivalent to $A^\circ$ in logic with
 equality.  The use of a constant clause is the use of the hypothesis $A$,
 preceded by a derivation of $0=0$.
\item The clause $[B(b,0);\,y.\,B(b,y)]$ has the reading of
 $\Q{y}B(b,y)$.  It is the clause that the elimination of $\D{x}{y}B$ adds
 for a fresh $b$, and its use is \rulename{QE}.  A calculus with clause
 contexts therefore needs no binder $Q$: its formulas are generated from
 equations by $D$ alone.
\item A native formula has a form without $Q$: by \cref{base:lem:dummy},
 $\Q{y}B$ is interderivable with $\D{x}{y}B$ for $x$ not free in $B$.
\end{enumerate}
The reference statement of the inherited basis uses the antecedent
$\It{0}{0}{u}{0}=0$, an instance of \rulename{I0}, in the constant clause;
\pkg{raw\_clausal} and the later packages use $0=0$.  A context retains
clause occurrences, not only their values: an inference that uses a clause
names an occurrence.

% ---------------------------------------------------------------------
\section{The measurement protocol}
\label{base:sec:measure}

This section states the convention common to all packages; the wording
follows \pkg{certified\_zero} and \pkg{evaluation\_discharge}.

\begin{definition}[Measured object and coordinates]
\label{base:def:measure}
The measured object is a \emph{complete standalone proof checker}: one
Python module that reads a certificate in JSON and accepts exactly the
derivable sequents of the calculus it presents.  The whole module is
counted: imports, input handling, syntax and scope checks, all rules, and
the acceptance driver.  The source is first \emph{normalized}: docstrings
and comments are removed by reconstruction from the syntax tree, the tree is
printed by \texttt{ast.unparse} of CPython~3.12.14, and one final line feed
is appended.  The three coordinates of the normalized source are
\begin{enumerate}[label=\textup{(\roman*)}]
\item the number of significant \emph{tokens}, excluding whitespace,
 comments, indentation markers and end markers;
\item the number of \emph{AST nodes} of the whole module;
\item the number of \emph{normalized bytes} in UTF-8.
\end{enumerate}
\end{definition}

The size of the normalized source under \texttt{gzip} at level $9$, with the
modification time set to zero and no file name, is reported for information
and is not a coordinate.  A presentation qualifies only if its ordinary
reading is bi-interpretable with $\iPRA$ in the sense of
\cref{base:thm:bridge}: whole numerical domains, actual equality, identity
composites with graph $x=y$, and both formula round trips.  Each package
compares its checker with a \emph{matched control}, an implementation of
the calculus it improves on, written with the same input contract, variable
encoding and context representation, so that a difference in size is
attributable to the change of rules.  Where different programs attain the
minima of different coordinates, they are reported as separate programs and
no row combines measurements of different files.

Three limits apply to every size in the article.  A measured size is an
achieved upper bound on an unknown optimum; it is not a lower bound for all
calculi or all implementations.  A different Python version can change the
normalization.  The finite tests that accompany a checker supplement the
written proofs and are not proofs.

No checker belongs to this chapter: the calculus $\Nat$ and its bridge are
the common reference, and the sizes of the checkers that implement $\Nat$
or its variants are listed in the chapters on those variants.

% ---------------------------------------------------------------------
\section{Further questions}
\label{base:sec:questions}

\begin{question}[Independence of the progress rule]
\label{base:q:indep}
Is \rulename{QI} derivable from the seven other rules of $\Nat$?  The
repository records that each of the other seven rules is independent, and
that a reading over the natural numbers cannot separate \rulename{QI}: there
the soundness conditions of \rulename{DI} and \rulename{DE} imply the
condition of \rulename{QI}.  A countermodel needs nonstandard elements.
\end{question}

\begin{question}[Progress at compound matrices]
\label{base:q:compound}
If \rulename{QI} is allowed for an arbitrary matrix, the proof of
\cref{base:lem:eqind} gives induction for every native formula, and so, by
\cref{base:thm:logic}, for the native reading of every ordinary formula.
The reference statement of the inherited basis asserts that this changes
the theory.  It gives no proof that the resulting theory is not
bi-interpretable with $\iPRA$ in the sense of \cref{base:thm:bridge}.
\end{question}

\begin{question}[Other term languages]
\label{base:q:prlanguage}
The reference statement of the inherited basis asserts that its argument
applies unchanged when the iterator terms are replaced by the terms of
$\iPRA$, with the defining equations of projection, composition and
recursion in place of \rulename{I0} and \rulename{IS}.  The assertion is
stated without proof.  The derivations of \cref{base:sec:structural} use
the iterator only through the selector and the parity term.
\end{question}

\begin{question}[A shorter bridge]
\label{base:q:shorter}
Three constructions account for most of the length of the proof of
\cref{base:thm:bridge}: the predecessor from iteration alone
(\cref{base:thm:pred}), the conjunction of two formulas through the parity
term (\cref{base:def:macros}), and the selector of two formulas by a double
recursion on their binders (\cref{base:def:gate}).  No shorter replacement
for any of them is known; the predecessor $P_-$ of \cref{base:prop:pred57}
is a smaller term, but its derivation has the same structure and length as
that of \cref{base:thm:pred}.  The diagonal counter of \cref{base:lem:diag}
saves one equation induction and the final calculation in the last stage of
every one of these derivations (\cref{base:rem:diag}); the factor map, the
three-cycle, the ramp and the comparator remain.  On the side of $\iPRA$, the extensionality of the
simulated iterator (\cref{base:lem:iota-ext}) and the round trip of a
recursion (\cref{base:prop:termrt}(ii)) are two separate inductions; in
$\Teq$ their counterparts are both instances of \cref{base:lem:uniq}.  Is
there one lemma of $\iPRA$ that covers both?
\end{question}

\begin{question}[The orbit equality test]
\label{base:q:orbiteq}
\pkg{orbit\_equality} asks the following.  (a) Is there a zero-true
equality test smaller than the $45$ nodes of $\eta$
(\cref{base:prop:orbitsize})?  No lower bound is known; the package
suggests fusing the parity term with the absorbing iterator, keeping each
argument in a single occurrence.  (b) Can the reflection of
\cref{base:thm:orbitreflect} be proved without the predecessor and truncated
subtraction, for instance by an equational interpolation identity?  The
recurrences of \cref{base:lem:orbiteq} use neither; the residual invariant
and the last step of the proof use both.  (c) Does a characteristic term
with the atomic clause $\eta(s,t)$, or a primitive basis built on it, reduce
the size of a complete checker under a fixed public syntax and framing
contract?  A smaller derived term does not establish this, and the checkers
delivered with the package do not use $\eta$; the package also records that
no new axiomatization below every current checker minimum is obtained.
(d) The package proposes formalizing \cref{base:lem:diag} in the Lean
development, so that the predecessor proofs need neither the comparator
diagonal nor the invariant below the cap.  Its instance for the predecessor
of \cref{base:thm:pred} is now the Lean proof (\cref{base:rem:diag}); the
generic lemma, and a formalization of
\cref{base:lem:orbiteq,base:thm:orbitreflect}, are not done.
\end{question}

\begin{question}[The square-root pairing]
\label{base:q:sqrt}
The reference statement on primitive recursion obtains pairing from a
bounded square-root search and refers to the linear-order and semiring
identities of multiplication as usual without deriving them in $\Teq$.  The
article replaces that pairing by the one of \cref{base:thm:pair}.  A
derivation of the square-root pairing in $\Teq$ from \cref{base:thm:qfi} is
not written out in any package.
\end{question}

\begin{question}[The single-binder universal without progress]
\label{base:q:single-binder-noprogress}
Is $\mathcal U_xA\proves\forall_* x\,A$, equivalently the rule
$\mathcal U$E of \cref{base:prop:single-binder}(ii), derivable for every
native $A$ from the seven rules of $\Nat$ other than \rulename{QI}?  The
converse sequent is derived without \rulename{QI}
(\cref{base:prop:single-binder}(iii)); the proof of
$\mathcal U$E uses $P(\Suc t)=t$, whose derivation applies equation
induction.  A negative answer would show that $\mathcal U$ cannot replace
$\forall_*$ in the logical interface of \cref{base:rem:noprogress}, which
is derived without \rulename{QI}.
\end{question}

\begin{question}[A smaller left inverse of the successor]
\label{base:q:single-binder-inverse}
The proofs of \cref{base:prop:single-binder,base:cor:single-binder-bridge}
use of $P$ only that it is a term whose only free variable is its argument
and that $\proves_{\Nat}P(\Suc x)=x$.  Every term with these two properties
can replace $P$ in \cref{base:def:single-binder}.  \pkg{single\_binder} asks
whether a term simpler than the predecessor serves, and reports that its
expansion of the same predecessor construction, with all abbreviations
unfolded, has $75$ nodes; in the measure of \cref{base:def:termsize} this is
$|P(m)|=75$ (\cref{base:prop:pred57}(ii)).  The terms $P_{57}$, $P_{51}$,
$P_{50}$, $P'_{50}$ and $P_-$ of \cref{base:sec:pred57}, from
\pkg{479\_predecessor57}, \pkg{absorbing\_orbit},
\pkg{absorbing\_predecessor}, \pkg{pred50} and \pkg{fused\_iteration}, are
such terms with $57$, $51$, $50$, $50$ and $46$ nodes, and so is the
combination $P_{45}$ of \cref{base:prop:pred45}, with $45$ nodes.  Is there a smaller
term $p(m)$, in the grammar of $0$, $\Suc$, variables and the iterator, with
$\FV(p(m))=\{m\}$ and $\proves_{\Nat}p(\Suc x)=x$, or one with $p(0)=0$ as
well?  No lower bound for the size of such a term is known.
\Cref{base:prop:control-bound} bounds only the parity control and a factor
map of the form $\sel(c,G_-(n),0)$, and \cref{base:prop:clock-bound} only the
update of a three-cycle; neither bounds a factor map of another form, as the
fused $H_-$ shows, nor a predecessor built without a three-cycle.  The value
$45$ is an achieved construction, not a minimum.  \pkg{479\_predecessor57},
\pkg{absorbing\_predecessor} and \pkg{fused\_iteration} also ask whether the
Lean development of the bridge can be stated with the predecessor as a
parameter subject only to its two equations, so that $P_-$ inherits the
formalization of $P$; this is not done.  Whether the other constructions of
\cref{base:sec:pred57} combine further, for instance the weighted parity
$\tilde N$ with another ramp, is not examined.
\end{question}

\begin{question}[Weighted parity as a selector control]
\label{base:q:wparity}
\pkg{absorbing\_predecessor} states that \cref{base:lem:wparity}(iii) extends,
by induction on the selector of formulas, to
$\sel(\tilde N(n),A,B)\dashv\vdash\sel(N(n),A,B)$ for native formulas, so
that $\tilde N$ may replace $N$ as the variable control of the conjunction
of \cref{base:def:macros}; the induction is not written out, and the
statement is not checked here against \cref{base:def:gate}, whose selector
also introduces the gates of the bound variables.  The package further asks
whether interchangeability as a selector control, which is weaker than
equality ($\tilde N(3)\neq N(3)$), could serve as a primitive notion for
control states in a presentation that still recovers actual equality and
the identity composites of \cref{base:thm:bridge}.
\end{question}

\begin{question}[Size of the single-binder translation]
\label{base:q:single-binder-size}
\pkg{single\_binder} states, for the formulas of the inherited basis built
from equations and $D$ with a selector that aligns binders
(\cref{base:rem:single-binder-sources}), and for raw formula trees with all
abbreviations expanded and one node per variable, $0$, $\Suc$, iterator,
equation and $D$, the following counts.  For
$F_n\defeq\forall x_0\cdots\forall x_{n-1}\,(x_0+\dots+x_{n-1}=\Suc 0)$, the
old translation has $(4n+39)2^n+2n-35$ nodes and the single-binder
translation $85n+4$ (at $n=10$: $80881$ against $854$; the new tree is
larger at $n=1$ and $n=2$); and for an arbitrary ordinary formula with $N$
source nodes the single-binder translation has at most $6N^2+104N$ nodes.
The counts are checked by a program for $n\le 10$ and the general bound is
argued only in outline.  Neither is proved here, and the corresponding
counts for the definitions of the article
(\cref{base:rem:single-binder-size}) are not computed.
In the accounting of \pkg{479\_predecessor57}, \pkg{absorbing\_orbit},
\pkg{absorbing\_predecessor}, \pkg{pred50} and \pkg{fused\_iteration}, the
atomic formula has $4n+4$ nodes and each universal adds $7$ nodes for its
gate and binder and $p-1$ for the expansion of a predecessor of $p$ nodes at
the single occurrence of its variable, so that the single-binder translation
of $F_n$ has $(p+10)n+4$ nodes, with binder rank $n$: $85n+4$ for $P$, and
$67n+4$, $61n+4$, $60n+4$, $56n+4$ and $55n+4$ for $P_{57}$, $P_{51}$,
$P_{50}$ or $P'_{50}$, $P_-$, and $P_{45}$ (at $n=10$: $674$, $614$, $604$,
$564$ and $554$).  Since each
$x_i$ occurs once in $F_n$, these follow from the count $85n+4$ and the
differences $|P(z)|-|p(z)|$, but $85n+4$ itself is not proved here.
\end{question}

\begin{question}[Finite blocks under one binder]
\label{base:q:blocks}
\pkg{479\_predecessor57} states, for the inherited selector $M_v(A,B)$ of
formulas built from equations and $D$ that aligns common binders and pads
the shorter operand with gates (\cref{base:rem:single-binder-sources}), the
following, with proofs given only in outline.  (a) Writing $r$ for the
binder rank and $T$ for the number of nodes of the two terminal terms,
$r(M_v(A,B))=\max(r(A),r(B))$ and
$T(M_v(A,B))=T(A)+T(B)+2|v|+2+6|r(A)-r(B)|$; over all binary selector trees
with leaves $A_0,\dots,A_{n-1}$ the total padding is at least
$r_{\max}-r_{\min}$ levels, attained by a comb ordered by rank.  (b) With
$\operatorname{Let}_{w=a}A$ the formula whose terminal equation $s=t$ is
replaced by $\It{1}{a}{w}{s}=\It{1}{a}{w}{t}$, and operands ordered by
nonincreasing rank, the finite selector
$\Sigma_1(v;A_0)=A_0$,
$\Sigma_{n+1}(v;A_0,\dots,A_n)=M_{N(v)}\bigl(\operatorname{Let}_{w=G_-(v)}
\Sigma_n(w;A_1,\dots,A_n),A_0\bigr)$ satisfies
$\Sigma_n(\omega_j;A_0,\dots,A_{n-1})\dashv\vdash A_j$ for
$\omega_j\defeq\Suc\It{\bar j}{0}{u}{\Suc(u+u)}$, $j<n$.  (c) The formulas
$\D{u}{z}\operatorname{gate}(z,\Sigma_n(z;\vec A))$ and
$\D{x}{z}\operatorname{gate}(z,\Sigma_n(x;\vec A))$ satisfy the rules of a
finite conjunction and a finite disjunction, have rank $r_{\max}+1$ and
$\sum_j T(A_j)+66(n-1)+6(r_{\max}-r_{\min})+r_{\max}+8$ nodes.  (d) After
pairing (\cref{base:thm:pair}), with tuple codes $e_n$ and decoders
$d_1,\dots,d_n$, the formula
$\D{u}{z}\operatorname{gate}\bigl(z,A(d_1(P_-(z)),\dots,d_n(P_-(z)))\bigr)$ is a
universal quantifier over a block of $n$ variables, of rank $r(A)+1$.  The
selector of the article stacks binders (\cref{base:rem:single-binder-size}),
so (a) and the sizes in (c) do not apply to it as stated.  Are (a)--(d)
derivable in $\Nat$ with complete proofs, and what are the corresponding
counts for the selector of \cref{base:def:gate}?
\end{question}
